% generated by GAPDoc2LaTeX from XML source (Frank Luebeck)
\documentclass[a4paper,11pt]{report}

\usepackage[top=37mm,bottom=37mm,left=27mm,right=27mm]{geometry}
\sloppy
\pagestyle{myheadings}
\usepackage{amssymb}
\usepackage[utf8]{inputenc}
\usepackage{makeidx}
\makeindex
\usepackage{color}
\definecolor{FireBrick}{rgb}{0.5812,0.0074,0.0083}
\definecolor{RoyalBlue}{rgb}{0.0236,0.0894,0.6179}
\definecolor{RoyalGreen}{rgb}{0.0236,0.6179,0.0894}
\definecolor{RoyalRed}{rgb}{0.6179,0.0236,0.0894}
\definecolor{LightBlue}{rgb}{0.8544,0.9511,1.0000}
\definecolor{Black}{rgb}{0.0,0.0,0.0}

\definecolor{linkColor}{rgb}{0.0,0.0,0.554}
\definecolor{citeColor}{rgb}{0.0,0.0,0.554}
\definecolor{fileColor}{rgb}{0.0,0.0,0.554}
\definecolor{urlColor}{rgb}{0.0,0.0,0.554}
\definecolor{promptColor}{rgb}{0.0,0.0,0.589}
\definecolor{brkpromptColor}{rgb}{0.589,0.0,0.0}
\definecolor{gapinputColor}{rgb}{0.589,0.0,0.0}
\definecolor{gapoutputColor}{rgb}{0.0,0.0,0.0}

%%  for a long time these were red and blue by default,
%%  now black, but keep variables to overwrite
\definecolor{FuncColor}{rgb}{0.0,0.0,0.0}
%% strange name because of pdflatex bug:
\definecolor{Chapter }{rgb}{0.0,0.0,0.0}
\definecolor{DarkOlive}{rgb}{0.1047,0.2412,0.0064}


\usepackage{fancyvrb}

\usepackage{mathptmx,helvet}
\usepackage[T1]{fontenc}
\usepackage{textcomp}


\usepackage[
            pdftex=true,
            bookmarks=true,        
            a4paper=true,
            pdftitle={Written with GAPDoc},
            pdfcreator={LaTeX with hyperref package / GAPDoc},
            colorlinks=true,
            backref=page,
            breaklinks=true,
            linkcolor=linkColor,
            citecolor=citeColor,
            filecolor=fileColor,
            urlcolor=urlColor,
            pdfpagemode={UseNone}, 
           ]{hyperref}

\newcommand{\maintitlesize}{\fontsize{50}{55}\selectfont}

% write page numbers to a .pnr log file for online help
\newwrite\pagenrlog
\immediate\openout\pagenrlog =\jobname.pnr
\immediate\write\pagenrlog{PAGENRS := [}
\newcommand{\logpage}[1]{\protect\write\pagenrlog{#1, \thepage,}}
%% were never documented, give conflicts with some additional packages

\newcommand{\GAP}{\textsf{GAP}}

%% nicer description environments, allows long labels
\usepackage{enumitem}
\setdescription{style=nextline}

%% depth of toc
\setcounter{tocdepth}{1}





%% command for ColorPrompt style examples
\newcommand{\gapprompt}[1]{\color{promptColor}{\bfseries #1}}
\newcommand{\gapbrkprompt}[1]{\color{brkpromptColor}{\bfseries #1}}
\newcommand{\gapinput}[1]{\color{gapinputColor}{#1}}


\begin{document}

\logpage{[ 0, 0, 0 ]}
\begin{titlepage}
\mbox{}\vfill

\begin{center}{\maintitlesize \textbf{ \textsf{ANUPQ} \mbox{}}}\\
\vfill

\hypersetup{pdftitle= \textsf{ANUPQ} }
\markright{\scriptsize \mbox{}\hfill  \textsf{ANUPQ}  \hfill\mbox{}}
{\Huge \textbf{ ANU p\texttt{\symbol{45}}Quotient \mbox{}}}\\
\vfill

{\Huge  3.3.3 \mbox{}}\\[1cm]
{ 25 November 2025 \mbox{}}\\[1cm]
\mbox{}\\[2cm]
{\Large \textbf{ Greg Gamble\\
  \mbox{}}}\\
{\Large \textbf{ Werner Nickel\\
 \mbox{}}}\\
{\Large \textbf{ Eamonn O'Brien\\
   \mbox{}}}\\
\hypersetup{pdfauthor= Greg Gamble\\
  ;  Werner Nickel\\
 ;  Eamonn O'Brien\\
   }
\mbox{}\\[2cm]
\begin{minipage}{12cm}\noindent
 \textsf{ANUPQ} is maintained by \href{mailto:mhorn@rptu.de} {Max Horn}. For support requests, please use \href{https://github.com/gap-packages/anupq/issues} {our issue tracker}. \end{minipage}

\end{center}\vfill

\mbox{}\\
{\mbox{}\\
\small \noindent \textbf{ Greg Gamble\\
  }  Email: \href{mailto://Greg.Gamble@uwa.edu.au} {\texttt{Greg.Gamble@uwa.edu.au}}}\\
{\mbox{}\\
\small \noindent \textbf{ Eamonn O'Brien\\
   }  Email: \href{mailto://obrien@math.auckland.ac.nz} {\texttt{obrien@math.auckland.ac.nz}}\\
  Homepage: \href{https://www.math.auckland.ac.nz/~obrien} {\texttt{https://www.math.auckland.ac.nz/\texttt{\symbol{126}}obrien}}}\\
\end{titlepage}

\newpage\setcounter{page}{2}
{\small 
\section*{Copyright}
\logpage{[ 0, 0, 1 ]}
 {\copyright} 2001\texttt{\symbol{45}}2016 by Greg Gamble

 {\copyright} 2001\texttt{\symbol{45}}2005 by Werner Nickel

 {\copyright} 1995\texttt{\symbol{45}}2001 by Eamon O'Brien

 The \textsf{GAP} package \textsf{ANUPQ} is licensed under the \href{https://opensource.org/licenses/artistic-license-2.0} {Artistic License 2.0}. \mbox{}}\\[1cm]
\newpage

\def\contentsname{Contents\logpage{[ 0, 0, 2 ]}}

\tableofcontents
\newpage

 
\chapter{\textcolor{Chapter }{Introduction}}\label{Introduction}
\logpage{[ 1, 0, 0 ]}
\hyperdef{L}{X7DFB63A97E67C0A1}{}
{
  
\section{\textcolor{Chapter }{Overview}}\logpage{[ 1, 1, 0 ]}
\hyperdef{L}{X8389AD927B74BA4A}{}
{
  \index{ANUPQ} The \textsf{GAP}{\nobreakspace}4 package \textsf{ANUPQ} provides an interface to the ANU \texttt{pq} C program written by Eamonn O'Brien, making the functionality of the C program
available to \textsf{GAP}. Henceforth, we shall refer to the \textsf{ANUPQ} package when referring to the \textsf{GAP} interface, and to the ANU \texttt{pq} program or just \texttt{pq} when referring to that C program. 

 The \texttt{pq} program consists of implementations of the following algorithms: 

 
\begin{enumerate}
\item  A \emph{$p$\texttt{\symbol{45}}quotient algorithm} to compute pc\texttt{\symbol{45}}presentations for $p$\texttt{\symbol{45}}factor groups of finitely presented groups. 

       
\item  A \emph{$p$\texttt{\symbol{45}}group generation algorithm} to generate pc presentations of groups of prime power order. 

    
\item  A \emph{standard presentation algorithm} used to compute a canonical pc\texttt{\symbol{45}}presentation of a $p$\texttt{\symbol{45}}group. 

  
\item  An algorithm which can be used to compute the \emph{automorphism group} of a $p$\texttt{\symbol{45}}group. 

 This part of the \texttt{pq} program is not accessible through the \textsf{ANUPQ} package. Instead, users are advised to consider the \textsf{GAP}{\nobreakspace}4 package \textsf{AutPGrp} by Bettina Eick and Eamonn O'Brien, which implements a better algorithm in \textsf{GAP} for the computation of automorphism groups of $p$\texttt{\symbol{45}}groups.  
\end{enumerate}
   The current version of the \textsf{ANUPQ} package requires \textsf{GAP}{\nobreakspace}4.5, and version 1.5 of the \textsf{AutPGrp} package. All code that made the package compatible with earlier versions of \textsf{GAP} has been removed. If you must use an older \textsf{GAP} version and cannot upgrade, then you may try using an older \textsf{ANUPQ} version. However, you should not use versions of the \textsf{ANUPQ} package older than 2.2, since they are known to have bugs. }

 
\section{\textcolor{Chapter }{How to read this manual}}\logpage{[ 1, 2, 0 ]}
\hyperdef{L}{X8416D2657E7831A1}{}
{
  It is not expected that readers of this manual will read it in a linear
fashion from cover to cover; some sections contain material that is far too
technical to be absorbed on a first reading. 

 Firstly, installers of the \textsf{ANUPQ} package will need to read Chapter{\nobreakspace}\hyperref[Installing-ANUPQ]{`Installing the ANUPQ Package'}, if they have not already gleaned these details from the \texttt{README} file. 

 Once the \textsf{ANUPQ} package is installed, users of the \textsf{ANUPQ} package will benefit most by first reading Chapter{\nobreakspace}\hyperref[Mathematical Background and Terminology]{`Mathematical Background and Terminology'}, which gives a brief description of the background and terminology used (this
chapter also cites a number of references for further reading), and the
introduction of Chapter{\nobreakspace}\hyperref[Infrastructure]{`Infrastructure'} (skip the remainder of the chapter on a first reading). 

 Then the user/reader should pursue Chapter{\nobreakspace}\hyperref[non-interact]{`Non\texttt{\symbol{45}}interactive ANUPQ functions'} in detail, delving into Chapter{\nobreakspace}\hyperref[ANUPQ Options]{`ANUPQ Options'} as necessary for the options of the functions that are described. The user
will become best acquainted with the \textsf{ANUPQ} package by trying the examples. This chapter describes the
non\texttt{\symbol{45}}interactive functions of the \textsf{ANUPQ} package, i.e.{\nobreakspace}``one\texttt{\symbol{45}}shot'' functions that invoke the \texttt{pq} program in such a way that once \textsf{GAP} has got what it needs, the \texttt{pq} is allowed to exit. It is expected that most of the time, users will only need
these functions. 

 Advanced users will want to explore Chapter{\nobreakspace}\hyperref[Interactive ANUPQ functions]{`Interactive ANUPQ functions'} which describes all the interactive functions of the \textsf{ANUPQ} package; these are functions that extract information via a dialogue with a
running \texttt{pq} process. Occasionally, a user needs the ``next step''; the functions provided in this chapter make use of data from previous steps
retained by the \texttt{pq} program, thus allowing the user to interact with the \texttt{pq} program like one can when one uses the \texttt{pq} program as a stand\texttt{\symbol{45}}alone (see{\nobreakspace}\texttt{guide.dvi} in the \texttt{standalone\texttt{\symbol{45}}doc} directory). 

 After having read Chapters{\nobreakspace}\hyperref[non-interact]{`Non\texttt{\symbol{45}}interactive ANUPQ functions'} and{\nobreakspace}\hyperref[Interactive ANUPQ functions]{`Interactive ANUPQ functions'}, cross\texttt{\symbol{45}}references will have taken the reader into
Chapter{\nobreakspace}\hyperref[ANUPQ Options]{`ANUPQ Options'}; by this stage, the reader need only read the introduction of
Chapter{\nobreakspace}\hyperref[ANUPQ Options]{`ANUPQ Options'}. 

 After the reader has developed some facility with the \textsf{ANUPQ} package, she should explore the examples described in Appendix{\nobreakspace}\hyperref[Examples]{`Examples'}. 

 If you run into trouble using the \textsf{ANUPQ} functions, some troubleshooting hints are given in Section{\nobreakspace}\hyperref[Hints and Warnings regarding the use of Options]{`Hints and Warnings regarding the use of Options'}. If the troubleshooting hints don't help, Section{\nobreakspace}\hyperref[Authors and Acknowledgements]{`Authors and Acknowledgements'} below, gives contact details for the authors of the components of the \textsf{ANUPQ} package. }

 
\section{\textcolor{Chapter }{Authors and Acknowledgements}}\label{Authors and Acknowledgements}
\logpage{[ 1, 3, 0 ]}
\hyperdef{L}{X79D2480A7810A7CC}{}
{
  The C implementation of the ANU \texttt{pq} standalone was developed by Eamonn O'Brien. 

 An interactive interface using iostreams was developed with the assistance of
Werner Nickel by Greg Gamble. 

 The \textsf{GAP} 4 version of this package was adapted from the \textsf{GAP} 3 version by Werner Nickel. 

 A new co\texttt{\symbol{45}}maintainer, Max Horn, joined the team in November,
2011. 

 The authors would like to thank Joachim Neub{\"u}ser for his careful
proof\texttt{\symbol{45}}reading and advice, and for formulating
Chapter{\nobreakspace}\hyperref[Mathematical Background and Terminology]{`Mathematical Background and Terminology'}. 

 We would also like to thank Bettina Eick who by her testing and provision of
examples helped us to eliminate a number of bugs and who provided a number of
valuable suggestions for extensions of the package beyond the \textsf{GAP}{\nobreakspace}3 capabilities. 

 \index{bug reports} If you find a bug, the last section of \textsf{ANUPQ}'s \texttt{README} describes the information we need and where to send us a bug report; please
take the time to read this (i.e.{\nobreakspace}help us to help you). }

 }

 
\chapter{\textcolor{Chapter }{Mathematical Background and Terminology}}\label{Mathematical Background and Terminology}
\logpage{[ 2, 0, 0 ]}
\hyperdef{L}{X7E7F3B617F42EF03}{}
{
  In this chapter we will give a brief description of the mathematical notions
used in the algorithms implemented in the ANU \texttt{pq} program that are made accessible from \textsf{GAP} through this package. For proofs and details we will point to relevant places
in the published literature. Also we will try to give some explanation of
terminology that may help to use the ``low\texttt{\symbol{45}}level'' interactive functions described in Section{\nobreakspace}\hyperref[Low-level Interactive ANUPQ functions based on menu items of the pq program]{`Low\texttt{\symbol{45}}level Interactive ANUPQ functions based on menu items
of the pq program'}. However, users who intend to use these functions are strongly advised to
acquire a thorough understanding of the algorithms from the quoted literature.
There is little or no checking done in these functions and naive use may
result in incorrect results. 

 
\section{\textcolor{Chapter }{Basic notions}}\label{Basic-notions}
\logpage{[ 2, 1, 0 ]}
\hyperdef{L}{X79A052C47C92AF09}{}
{
  Throughout this manual, by $p$\texttt{\symbol{45}}group we always mean \emph{finite} $p$\texttt{\symbol{45}}group. 
\subsection{\textcolor{Chapter }{pc Presentations and Consistency}}\logpage{[ 2, 1, 1 ]}
\hyperdef{L}{X7BD675838609D547}{}
{
 For details, see e.g.{\nobreakspace}\cite{NNN98}. 

 Every finite $p$\texttt{\symbol{45}}group $G$ has a presentation of the form: 
\[ \{a_1,\dots,a_n \mid a_i^p = v_{ii}, 1 \le i \le n, [a_k, a_j] = v_{jk}, 1 \le
j < k \le n \}. \]
 where $v_{jk}$ is a word in the elements $a_{k+1},\dots,a_n$ for $1 \le j \leq k \le n$. 

 \index{power\texttt{\symbol{45}}commutator presentation}\index{pc presentation}\index{pcp} \index{pc generators}\index{collection} This is called a \emph{power\texttt{\symbol{45}}commutator} presentation (or \emph{pc presentation} or \emph{pcp}) of $G$, generators from such a presentation will be referred to as \emph{pc generators}. In terms of such pc generators every element of $G$ can be written in a ``normal form'' $a_1^{e_1}\dots a_n^{e_n}$ with $0 \le e_i < p$. Moreover any given product of the generators can be brought into such a
normal form using the defining relations in the above presentation as rewrite
rules. Any such process is called \emph{collection}. For the discussion of various collection methods see \cite{LGS90} and \cite{VL90a}. 

 \index{consistent}\index{confluent rewriting system}\index{confluent} Every $p$\texttt{\symbol{45}}group of order $p^n$ has such a pcp on $n$ generators and conversely every such presentation defines a $p$\texttt{\symbol{45}}group. However a $p$\texttt{\symbol{45}}group defined by a pcp on $n$ generators can be of smaller order $p^m$ with $m<n$. A pcp on $n$ generators that does in fact define a $p$\texttt{\symbol{45}}group of order $p^n$ is called \emph{consistent} in this manual, in line with most of the literature on the algorithms
occurring here. A consistent pcp determines a \emph{confluent rewriting system} (see{\nobreakspace}\texttt{IsConfluent} (\textbf{Reference: IsConfluent}) of the \textsf{GAP} Reference Manual) for the group it defines and for this reason often (in
particular in the \textsf{GAP} Reference Manual) such a pcp presentation is also called \emph{confluent}. 

 Consistency of a pcp is tantamount to the fact that for any given word in the
generators any two collections will yield the same normal form. 

 \index{consistency conditions} Consistency of a pcp can be checked by a finite set of \emph{consistency conditions}, demanding that collection of the left hand side and of the right hand side
of certain equations, starting with subproducts indicated by bracketing, will
result in the same normal form. There are 3 types of such equations (that will
be referred to in the manual): 
\[ \begin{array}{rclrl} (a^n)a &=& a(a^n) &&{\rm (Type 1)} \\ (b^n)a &=&
b^{(n-1)}(ba), b(a^n) = (ba)a^{(n-1)} &&{\rm (Type 2)} \\ c(ba) &=& (cb)a
&&{\rm (Type 3)} \\ \end{array} \]
 See \cite{VL84} for a description of a sufficient set of consistency conditions in the context
of the $p$\texttt{\symbol{45}}quotient algorithm. }

 
\subsection{\textcolor{Chapter }{Exponent\texttt{\symbol{45}}$p$ Central Series and Weighted pc Presentations}}\logpage{[ 2, 1, 2 ]}
\hyperdef{L}{X7944DB037F2277A6}{}
{
 For details, see \cite{NNN98}. 

 \index{exponent\texttt{\symbol{45}}p central series} The (\emph{descending} or \emph{lower}) (\emph{exponent\texttt{\symbol{45}}})\emph{$p$\texttt{\symbol{45}}central series} of an arbitrary group $G$ is defined by 
\[ P_0(G) := G, P_i(G) := [G, P_{i-1}(G)] P_{i-1}(G)^p. \]
 For a $p$\texttt{\symbol{45}}group $G$ this series terminates with the trivial group. $G$ \index{class}\index{p\texttt{\symbol{45}}class} has \emph{$p$\texttt{\symbol{45}}class} $c$ if $c$ is the smallest integer such that $P_c(G)$ is the trivial group. In this manual, as well as in much of the literature
about the \texttt{pq}\texttt{\symbol{45}} and related algorithms, the $p$\texttt{\symbol{45}}class is often referred to simply by \emph{class}. 

 Let the $p$\texttt{\symbol{45}}group $G$ have a consistent pcp as above. Then the subgroups 
\[ \langle1\rangle < {\langle}a_n\rangle < {\langle}a_n, a_{n-1}\rangle < \dots <
{\langle}a_n,\dots,a_i\rangle < \dots < G \]
 form a central series of $G$. If this refines the $p$\texttt{\symbol{45}}central series, \index{weight function} we can define the \emph{weight function} $w$ for the pc generators by $w(a_i) = k$, if $a_i$ is contained in $P_{k-1}(G)$ but not in $P_k(G)$. 

 \index{weighted pcp} The pair of such a weight function and a pcp allowing it, is called a \emph{weighted pcp}. }

 
\subsection{\textcolor{Chapter }{$p$\texttt{\symbol{45}}Cover, $p$\texttt{\symbol{45}}Multiplicator}}\logpage{[ 2, 1, 3 ]}
\hyperdef{L}{X7C43ACA37D391EBD}{}
{
 For details, see \cite{NNN98}. 

 \index{p\texttt{\symbol{45}}covering group}\index{p\texttt{\symbol{45}}cover} \index{p\texttt{\symbol{45}}multiplicator} \index{p\texttt{\symbol{45}}multiplicator rank} \index{multiplicator rank} Let $d$ be the minimal number of generators of the $p$\texttt{\symbol{45}}group $G$ of $p$\texttt{\symbol{45}}class $c$. Then $G$ is isomorphic to a factor group $F/R$ of a free group $F$ of rank $d$. We denote $[F, R] R^p$ by $R^*$. It can be proved (see e.g.{\nobreakspace}\cite{OBr90}) that the isomorphism type of $G^* := F/R^*$ depends only on $G$. $G^*$ is called the \emph{$p$\texttt{\symbol{45}}covering group} or \emph{$p$\texttt{\symbol{45}}cover} of $G$, and $R/R^*$ the \emph{$p$\texttt{\symbol{45}}multiplicator} of $G$. The $p$\texttt{\symbol{45}}multiplicator is, of course, an elementary abelian $p$\texttt{\symbol{45}}group; its minimal number of generators is called the \emph{($p$\texttt{\symbol{45}})multiplicator rank}. }

 
\subsection{\textcolor{Chapter }{Descendants, Capable, Terminal, Nucleus}}\logpage{[ 2, 1, 4 ]}
\hyperdef{L}{X801A27A08462AFAB}{}
{
 For details, see \cite{New77} and \cite{OBr90}. 

 \index{descendant}\index{immediate descendant}\index{nucleus} \index{capable}\index{terminal} Let again $G$ be a $p$\texttt{\symbol{45}}group of $p$\texttt{\symbol{45}}class $c$ and $d$ the minimal number of generators of $G$. A $p$\texttt{\symbol{45}}group $H$ is a \emph{descendant} of $G$ if the minimal number of generators of $H$ is $d$ and $H/P_c(H)$ is isomorphic to $G$. A descendant $H$ of $G$ is a \emph{proper descendant} if it has $p$\texttt{\symbol{45}}class at least $c+1$. A descendant $H$ of $G$ is an \emph{immediate descendant} if it has $p$\texttt{\symbol{45}}class $c+1$. $G$ is called \emph{capable} if it has immediate descendants; otherwise it is \emph{terminal}. 

 Let $G^* = F/R^*$ again be the $p$\texttt{\symbol{45}}cover of $G$. Then the group $P_c(G^*)$ is called the \emph{nucleus} of $G$. Note that $P_c(G^*)$ is contained in the $p$\texttt{\symbol{45}}multiplicator $R/R^*$. 

 \index{nucleus}\index{allowable subgroup} It is proved (e.g.{\nobreakspace}in \cite{OBr90}) that the immediate descendants of $G$ are obtained as factor groups of the $p$\texttt{\symbol{45}}cover by (proper) supplements of the nucleus in the
(elementary abelian) $p$\texttt{\symbol{45}}multiplicator. These are also called \emph{allowable}. 

 \index{extended automorphism}\index{permutations} It is further proved there that every automorphism $\alpha$ of $F/R$ extends to an automorphism $\alpha^*$ of the $p$\texttt{\symbol{45}}cover $F/R^*$ and that the restriction of $\alpha^*$ to the multiplicator $R/R^*$ is uniquely determined by $\alpha$. Each \emph{extended automorphism} $\alpha^*$ induces a permutation of the allowable subgroups. Thus the extended
automorphisms determine a group $P$ of \emph{permutations} on the set $A$ of allowable subgroups (The group $P$ of permutations will appear in the description of some interactive functions).
Choosing a representative $S$ from each orbit of $P$ on $A$, the set of factor groups $F/S$ contains each (isomorphism type of) immediate descendant of $G$ exactly once. For each immediate descendant, the procedure of computing the $p$\texttt{\symbol{45}}cover, extending the automorphisms and computing the
orbits on allowable subgroups can be repeated. Iteration of this procedure can
in principle be used to determine all descendants of a $p$\texttt{\symbol{45}}group. }

 
\subsection{\textcolor{Chapter }{Laws}}\logpage{[ 2, 1, 5 ]}
\hyperdef{L}{X81FBE7ED79EFF5EF}{}
{
 \index{law}\index{identical relation}\index{exponent law} \index{metabelian law}\index{Engel identity} Let $l(x_1, \dots, x_n)$ be a word in the free generators $x_1, \dots, x_n$ of a free group of rank $n$. Then $l(x_1, \dots, x_n) = 1$ is called a \emph{law} or \emph{identical relation} in a group $G$ if $l(g_1, \dots, g_n) = 1$ for any choice of elements $g_1, \dots, g_n$ in $G$. In particular, $x^e = 1$ is called an \emph{exponent law}, $[[x,y],[u,v]] = 1$ the \emph{metabelian law}, and $[\dots [[x_1,x_2],x_2],\dots, x_2] = 1$ an \emph{Engel identity}. }

 }

 
\section{\textcolor{Chapter }{The p\texttt{\symbol{45}}quotient Algorithm}}\label{The p-quotient Algorithm}
\logpage{[ 2, 2, 0 ]}
\hyperdef{L}{X7C8CE96D80FC3614}{}
{
  For details, see \cite{HN80}, \cite{NO96} and \cite{VL84}. Other descriptions of the algorithm are given in \cite{Sims94}. 

 The \texttt{pq} algorithm successively determines the factor groups of the groups of the $p$\texttt{\symbol{45}}central series of a finitely presented (fp) group $G$. If a bound $b$ for the $p$\texttt{\symbol{45}}class is given, the algorithm will determine those factor
groups up to at most $p$\texttt{\symbol{45}}class $b$. If the $p$\texttt{\symbol{45}}central series terminates with a subgroup $P_k(G)$ with $k < b$, the algorithm will stop with that group. If no such bound is given, it will
try to find the biggest such factor group. 

 $G/P_1(G)$ is the largest elementary abelian $p$\texttt{\symbol{45}}factor group of $G$ and this can be found from the relation matrix of $G$ using matrix diagonalisation modulo $p$. So it suffices to explain how $G/P_{i+1}(G)$ is found from $G$ and $G/P_i(G)$ for some $i \ge 1$. 

 This is done, in principle, in two steps: first the $p$\texttt{\symbol{45}}cover of $G_i := G/P_i(G)$ is determined (which depends only on $G_i$, not on $G$) and then $G/P_{i+1}(G)$ as a factor group of this $p$\texttt{\symbol{45}}cover. 
\subsection{\textcolor{Chapter }{Finding the $p$\texttt{\symbol{45}}cover}}\logpage{[ 2, 2, 1 ]}
\hyperdef{L}{X791EB6F77899CB3D}{}
{
 A very detailed description of the first step is given in \cite{NNN98}, from which we just extract some passages in order to point to some terms
occurring in this manual. 

 \index{labelled pcp}\index{definition!of generator} Let $H$ be a $p$\texttt{\symbol{45}}group and $p^{d(b)}$ be the order of $H/P_b(H)$. So $d := d(1)$ is the minimal number of generators of $H$. A weighted pcp of $H$ will be called \emph{labelled} if for each generator $a_k$, $k > d$ one relation, having this generator as its right hand side, is marked as \emph{definition} of this generator. 

 As described in \cite{NNN98}, a weighted labelled pcp of a $p$\texttt{\symbol{45}}group can be obtained stepping down its $p$\texttt{\symbol{45}}central series. 

 So let us assume that a weighted labelled pcp of $G_i$ is given. A straightforward way of of writing down a (not necessarily
consistent) pcp for its $p$\texttt{\symbol{45}}cover is to add generators, one for each relation which is
not a definition, and modify the right hand side of each such relation by
multiplying it on the right by one of the new generators
\texttt{\symbol{45}}\texttt{\symbol{45}} a different generator for each such
relation. Further relations are then added to make the new generators central
and of order $p$. This procedure is called \emph{adding tails}. A more formal description of it is again given in \cite{NNN98}. 

 \index{tails} It is important to realise that the ``new'' generators will generate an elementary abelian group, that is, in additive
notation, a vector space over the field of $p$ elements. As said, the pcp of the $p$\texttt{\symbol{45}}cover obtained in this way need not be consistent. Since
the pcp of $G_i$ was consistent, applying the consistency conditions to the pcp of the $p$\texttt{\symbol{45}}cover, in case the presentation obtained for $p$\texttt{\symbol{45}}cover is not consistent, will produce a set of equations
between the new generators, that, written additively, are linear equations
over the field of $p$ elements and can hence be used to remove redundant generators until a
consistent pcp is obtained. 

 In reality, to follow this straightforward procedure would be forbiddingly
inefficient except for very small examples. There are many ways of a priori
reducing the number of ``new generators'' to be introduced, using e.g.{\nobreakspace}the weights attached to the
generators, and the main part of \cite{NNN98} is devoted to a detailed discussion with proofs of these possibilities. }

 
\subsection{\textcolor{Chapter }{Imposing the Relations of the fp Group}}\logpage{[ 2, 2, 2 ]}
\hyperdef{L}{X804CF5C97F7BB880}{}
{
 In order to obtain $G/P_{i+1}(G)$ from the pcp of the $p$\texttt{\symbol{45}}cover of $G_i = G/P_i(G)$, the defining relations from the original presentation of $G$ must be imposed. Since $G_i$ is a homomorphic image of $G$, these relations again yield relations between the ``new generators'' in the presentation of the $p$\texttt{\symbol{45}}cover of $G_i$. }

 
\subsection{\textcolor{Chapter }{Imposing Laws}}\logpage{[ 2, 2, 3 ]}
\hyperdef{L}{X7F1A8CCD84462775}{}
{
 While we have so far only considered the computation of the factor groups of a
given fp group by the groups of its descending $p$\texttt{\symbol{45}}central series, the $p$\texttt{\symbol{45}}quotient algorithm allows a very important variant of this
idea: laws can be prescribed that should be fulfilled by the $p$\texttt{\symbol{45}}factor groups computed by the algorithm. The key
observation here is the fact that at each step down the descending $p$\texttt{\symbol{45}}central series it suffices to impose these laws only for a
finite number of words. Again for efficiency of the method it is crucial to
keep the number of such words small, and much of \cite{NO96} and the literature quoted in this paper is devoted to this problem. 

 \index{exponent check} In this form, starting with a free group and imposing an exponent law (also
referred to as an \emph{exponent check}) the \texttt{pq} program has, in fact, found its most noted application in the determination of
(restricted) Burnside groups (as reported in e.g.{\nobreakspace}\cite{HN80}, \cite{NO96} and \cite{VL90b}). 

 Via a \textsf{GAP} program using the ``local'' interactive functions of the \texttt{pq} program made available through this interface also arbitrary laws can be
imposed via the option \texttt{Identities} (see{\nobreakspace}\ref{option Identities}). }

 }

 
\section{\textcolor{Chapter }{The p\texttt{\symbol{45}}group generation Algorithm, Standard Presentation,
Isomorphism Testing}}\label{The p-group generation Algorithm, Standard Presentation, Isomorphism Testing}
\logpage{[ 2, 3, 0 ]}
\hyperdef{L}{X807FB2EC85E6648D}{}
{
  For details, see \cite{New77} and \cite{OBr90}. 

 \index{p\texttt{\symbol{45}}group generation}\index{orbits} The $p$\texttt{\symbol{45}}group generation algorithm determines the immediate
descendants of a given $p$\texttt{\symbol{45}}group $G$ up to isomorphism. From what has been explained in Section{\nobreakspace}\hyperref[Basic-notions]{`Basic notions'}, it is clear that this amounts to the construction of the $p$\texttt{\symbol{45}}cover, the extension of the automorphisms of $G$ to the $p$\texttt{\symbol{45}}cover and the determination of representatives of the
orbits of the action of these automorphisms on the set of supplements of the
nucleus in the $p$\texttt{\symbol{45}}multiplicator. 

 The main practical problem here is the determination of these representatives. \cite{OBr90} describes methods for this and the \texttt{pq} program allows choices according to whether space or time limitations must be
met. 

 As well as the descendants of $G$, the \texttt{pq} program determines their automorphism groups from that of $G$ (see{\nobreakspace}\cite{OBr95}), which is important for an iteration of the process; this has been used by
Eamonn O'Brien, e.g.{\nobreakspace}in the classification of the $2$\texttt{\symbol{45}}groups that are now also part of the \emph{Small Groups} library available through \textsf{GAP}. 

 \index{standard presentation}\index{echelonised matrix} \index{label of standard matrix} A variant of the $p$\texttt{\symbol{45}}group generation algorithm is also used to define a \emph{standard presentation} of a given $p$\texttt{\symbol{45}}group. This is done by constructing an isomorphic copy of
the given group through a chain of descendants and at each step making a
choice of a particular representative for the respective orbit of capable
groups. In a fairly delicate process, subgroups of the $p$\texttt{\symbol{45}}multiplicator are represented by \emph{echelonised matrices} and a first among the \emph{labels for standard matrices} is chosen (this is described in detail in \cite{OBr94}). 

 \index{isomorphism testing}\index{compaction} Finally, the standard presentation provides a way of testing if two given $p$\texttt{\symbol{45}}groups are isomorphic: the standard presentations of the
groups are computed, for practical purposes \emph{compacted} and the results compared for being identical, i.e.{\nobreakspace}the groups
are isomorphic if and only if their standard presentations are identical. }

 }

 
\chapter{\textcolor{Chapter }{Infrastructure}}\label{Infrastructure}
\logpage{[ 3, 0, 0 ]}
\hyperdef{L}{X7917EFDF7AC06F04}{}
{
  Most of the details in this chapter are of a technical nature; the user need
only skim over this chapter on a first reading. Mostly, it is enough to know
that 
\begin{itemize}
\item  you must do a \texttt{LoadPackage("anupq");} before you can expect to use a command defined by the \textsf{ANUPQ} package (details are in Section{\nobreakspace}\hyperref[Loading the ANUPQ Package]{`Loading the ANUPQ Package'}); 
\item  partial results of \textsf{ANUPQ} commands and some other data are stored in the \texttt{ANUPQData} global variable (details are in Section{\nobreakspace}\hyperref[The ANUPQData Record]{`The ANUPQData Record'}); 
\item  doing \texttt{SetInfoLevel(InfoANUPQ, \mbox{\texttt{\mdseries\slshape n}});} for \mbox{\texttt{\mdseries\slshape n}} greater than the default value 1 will give progressively more information of
what is going on ``behind the scenes'' (details are in Section{\nobreakspace}\hyperref[Setting the Verbosity of ANUPQ via Info and InfoANUPQ]{`Setting the Verbosity of ANUPQ via Info and InfoANUPQ'}); 
\item  in Section{\nobreakspace}\hyperref[Utility-Functions]{`Utility Functions'} we describe some utility functions and functions that run examples from the
collection of examples of this package; 
\item  in Section{\nobreakspace}\hyperref[Attributes and a Property for fp and pc p-groups]{`Attributes and a Property for fp and pc p\texttt{\symbol{45}}groups'} we describe the attributes and property \texttt{NuclearRank}, \texttt{MultiplicatorRank} and \texttt{IsCapable}; and 
\item  in Section{\nobreakspace}\hyperref[Hints and Warnings regarding the use of Options]{`Hints and Warnings regarding the use of Options'} we describe some troubleshooting strategies. Also this section explains the
utility of setting \texttt{ANUPQWarnOfOtherOptions := true;} (particularly for novice users) for detecting misspelt options and diagnosing
other option usage problems. 
\end{itemize}
 
\section{\textcolor{Chapter }{Loading the ANUPQ Package}}\label{Loading the ANUPQ Package}
\logpage{[ 3, 1, 0 ]}
\hyperdef{L}{X833D58248067E13B}{}
{
  \index{banner} To use the \textsf{ANUPQ} package, as with any \textsf{GAP} package, it must be requested explicitly. This is done by calling 
\begin{Verbatim}[commandchars=!|D,fontsize=\small,frame=single,label=Example]
  !gapprompt|gap>D !gapinput|LoadPackage( "anupq" );D
  ---------------------------------------------------------------------------
  Loading ANUPQ 3.3.3 (ANU p-Quotient)
  by Greg Gamble (GAP code, Greg.Gamble@uwa.edu.au),
     Werner Nickel (GAP code), and
     Eamonn O'Brien (C code, https://www.math.auckland.ac.nz/~obrien).
  maintained by:
     Max Horn (https://www.quendi.de/math).
  uses ANU pq binary (C code program) version: 1.9
  Homepage: https://gap-packages.github.io/anupq/
  Report issues at https://github.com/gap-packages/anupq/issues
  ---------------------------------------------------------------------------
  true
\end{Verbatim}
 Note that since the \textsf{ANUPQ} package uses the \texttt{AutomorphimGroupPGroup} function of the \textsf{AutPGrp} package and, in any case, often needs other \textsf{AutPGrp} functions when computing descendants, the user must ensure that the \textsf{AutPGrp} package is also installed, at least version 1.5. If the \textsf{AutPGrp} package is not installed, the \textsf{ANUPQ} package will \texttt{fail} to load. 

 Also, if \textsf{GAP} cannot find a working \texttt{pq} binary, the call to \texttt{LoadPackage} will return \texttt{fail}. 

 If you want to load the \textsf{ANUPQ} package by default, you can put the \texttt{LoadPackage} command into your \texttt{gap.ini} file (see Section{\nobreakspace} \textbf{Reference: The gap.ini and gaprc files} in the \textsf{GAP} Reference Manual). By the way, the novice user of the \textsf{ANUPQ} package should probably also append the line 
\begin{Verbatim}[commandchars=!@|,fontsize=\small,frame=single,label=Example]
  ANUPQWarnOfOtherOptions := true;
\end{Verbatim}
 to their \texttt{gap.ini} file, somewhere after the \texttt{LoadPackage( "anupq" );} command (see{\nobreakspace}\texttt{ANUPQWarnOfOtherOptions} (\ref{ANUPQWarnOfOtherOptions})). }

 
\section{\textcolor{Chapter }{The ANUPQData Record}}\label{The ANUPQData Record}
\logpage{[ 3, 2, 0 ]}
\hyperdef{L}{X83DE155A79C38DBE}{}
{
  This section contains fairly technical details which may be skipped on an
initial reading. 

\subsection{\textcolor{Chapter }{ANUPQData}}
\logpage{[ 3, 2, 1 ]}\nobreak
\hyperdef{L}{X7B90E89782BDA6D7}{}
{\noindent\textcolor{FuncColor}{$\triangleright$\enspace\texttt{ANUPQData\index{ANUPQData@\texttt{ANUPQData}}
\label{ANUPQData}
}\hfill{\scriptsize (global variable)}}\\


 is a \textsf{GAP} record in which the essential data for an \textsf{ANUPQ} session within \textsf{GAP} is stored; its fields are: 
\begin{description}
\item[{\texttt{binary}}]  the path of the \texttt{pq} binary; 
\item[{\texttt{tmpdir}}]  the path of the temporary directory used by the \texttt{pq} binary and \textsf{GAP} (i.e.{\nobreakspace}the directory in which all the \texttt{pq}'s temporary files are created) (also see \texttt{ANUPQDirectoryTemporary} (\ref{ANUPQDirectoryTemporary}) below); 
\item[{\texttt{outfile}}]  the full path of the default \texttt{pq} output file; 
\item[{\texttt{SPimages}}]  the full path of the file \texttt{GAP{\textunderscore}library} to which the \texttt{pq} program writes its Standard Presentation images; 
\item[{\texttt{version}}]  the version of the current \texttt{pq} binary; 
\item[{\texttt{ni}}]  a data record used by non\texttt{\symbol{45}}interactive functions (see below
and Chapter{\nobreakspace}\hyperref[non-interact]{`Non\texttt{\symbol{45}}interactive ANUPQ functions'}); 
\item[{\texttt{io}}]  list of data records for \texttt{PqStart} (see below and{\nobreakspace}\texttt{PqStart} (\ref{PqStart})) processes; 
\item[{\texttt{topqlogfile}}]  name of file logged to by \texttt{ToPQLog} (see{\nobreakspace}\texttt{ToPQLog} (\ref{ToPQLog})); and 
\item[{\texttt{logstream}}]  stream of file logged to by \texttt{ToPQLog} (see{\nobreakspace}\texttt{ToPQLog} (\ref{ToPQLog})). 
\end{description}
 Each time an interactive \textsf{ANUPQ} process is initiated via \texttt{PqStart} (see{\nobreakspace}\texttt{PqStart} (\ref{PqStart})), an identifying number \mbox{\texttt{\mdseries\slshape ioIndex}} is generated for the interactive process and a record \texttt{ANUPQData.io[\mbox{\texttt{\mdseries\slshape ioIndex}}]} with some or all of the fields listed below is created. Whenever a
non\texttt{\symbol{45}}interactive function is called (see
Chapter{\nobreakspace}\hyperref[non-interact]{`Non\texttt{\symbol{45}}interactive ANUPQ functions'}), the record \texttt{ANUPQData.ni} is updated with fields that, if bound, have exactly the same purpose as for a \texttt{ANUPQData.io[\mbox{\texttt{\mdseries\slshape ioIndex}}]} record. 
\begin{description}
\item[{\texttt{stream}}]  the IOStream opened for interactive \textsf{ANUPQ} process \mbox{\texttt{\mdseries\slshape ioIndex}} or non\texttt{\symbol{45}}interactive \textsf{ANUPQ} function; 
\item[{\texttt{group}}]  the group given as first argument to \texttt{PqStart}, \texttt{Pq}, \texttt{PqEpimorphism}, \texttt{PqDescendants} or \texttt{PqStandardPresentation} (or any synonymous methods); 
\item[{\texttt{haspcp}}]  is bound and set to \texttt{true} when a pc presentation is first set inside the \texttt{pq} program (e.g.{\nobreakspace}by \texttt{PqPcPresentation} or \texttt{PqRestorePcPresentation} or a higher order function like \texttt{Pq}, \texttt{PqEpimorphism}, \texttt{PqPCover}, \texttt{PqDescendants} or \texttt{PqStandardPresentation} that does a \texttt{PqPcPresentation} operation, but \emph{not} \texttt{PqStart} which only starts up an interactive \textsf{ANUPQ} process); 
\item[{\texttt{gens}}]  a list of the generators of the group \texttt{group} as strings (the same as those passed to the \texttt{pq} program); 
\item[{\texttt{rels}}]  a list of the relators of the group \texttt{group} as strings (the same as those passed to the \texttt{pq} program); 
\item[{\texttt{name}}]  the name of the group whose pc presentation is defined by a call to the \texttt{pq} program (according to the \texttt{pq} program \texttt{\symbol{45}}\texttt{\symbol{45}} unless you have used the \texttt{GroupName} option (see e.g.{\nobreakspace}\texttt{Pq} (\ref{Pq})) or applied the function \texttt{SetName} (see{\nobreakspace}\texttt{SetName} (\textbf{Reference: Name})) to the group, the ``generic'' name \texttt{"[grp]"} is set as a default); 
\item[{\texttt{gpnum}}]  if not a null string, the ``number'' (i.e.{\nobreakspace}the unique label assigned by the \texttt{pq} program) of the last descendant processed; 
\item[{\texttt{class}}]  the largest lower exponent\texttt{\symbol{45}}$p$ central class of a quotient group of the group (usually \texttt{group}) found by a call to the \texttt{pq} program; 
\item[{\texttt{forder}}]  the factored order of the quotient group of largest lower
exponent\texttt{\symbol{45}}$p$ central class found for the group (usually \texttt{group}) by a call to the \texttt{pq} program (this factored order is given as a list  \texttt{[p,n]}, indicating an order of $p^n$); 
\item[{\texttt{pcoverclass}}]  the lower exponent\texttt{\symbol{45}}$p$ central class of the $p$\texttt{\symbol{45}}covering group of a $p$\texttt{\symbol{45}}quotient of the group (usually \texttt{group}) found by a call to the \texttt{pq} program; 
\item[{\texttt{workspace}}]  the workspace set for the \texttt{pq} process (either given as a second argument to \texttt{PqStart}, or set by default to 10000000); 
\item[{\texttt{menu}}]  the current menu of the \texttt{pq} process (the \texttt{pq} program is managed by various menus, the details of which the user shouldn't
normally need to know about \texttt{\symbol{45}}\texttt{\symbol{45}} the \texttt{menu} field remembers which menu the \texttt{pq} process is currently ``in''); 
\item[{\texttt{outfname}}]  is the file to which \texttt{pq} output is directed, which is always \texttt{ANUPQData.outfile}, except when option \texttt{SetupFile} is used with a non\texttt{\symbol{45}}interactive function, in which case \texttt{outfname} is set to \texttt{"PQ{\textunderscore}OUTPUT"}; 
\item[{\texttt{pQuotient}}]  is set to the value returned by \texttt{Pq} (see{\nobreakspace}\texttt{Pq} (\ref{Pq})) (the field \texttt{pQepi} is also set at the same time); 
\item[{\texttt{pQepi}}]  is set to the value returned by \texttt{PqEpimorphism} (see{\nobreakspace}\texttt{PqEpimorphism} (\ref{PqEpimorphism})) (the field \texttt{pQuotient} is also set at the same time); 
\item[{\texttt{pCover}}]  is set to the value returned by \texttt{PqPCover} (see{\nobreakspace}\texttt{PqPCover} (\ref{PqPCover})); 
\item[{\texttt{SP}}]  is set to the value returned by \texttt{PqStandardPresentation} or \texttt{StandardPresentation} (see{\nobreakspace}\texttt{PqStandardPresentation} (\ref{PqStandardPresentation:interactive})) when called interactively, for process \mbox{\texttt{\mdseries\slshape i}} (the field \texttt{SPepi} is also set at the same time); 
\item[{\texttt{SPepi}}]  is set to the value returned by \texttt{EpimorphismPqStandardPresentation} or \texttt{EpimorphismStandardPresentation} (see{\nobreakspace}\texttt{EpimorphismPqStandardPresentation} (\ref{EpimorphismPqStandardPresentation:interactive})) when called interactively, for process \mbox{\texttt{\mdseries\slshape i}} (the field \texttt{SP} is also set at the same time); 
\item[{\texttt{descendants}}]  is set to the value returned by \texttt{PqDescendants} (see{\nobreakspace}\texttt{PqDescendants} (\ref{PqDescendants})); 
\item[{\texttt{treepos}}]  if set by a call to \texttt{PqDescendantsTreeCoclassOne} (see{\nobreakspace}\texttt{PqDescendantsTreeCoclassOne} (\ref{PqDescendantsTreeCoclassOne})), it contains a record with fields \texttt{class}, \texttt{node} and \texttt{ndes} being the information that determines the last descendant with a
non\texttt{\symbol{45}}zero number of descendants processed; 
\item[{\texttt{xgapsheet}}]  if set by a call to \texttt{PqDescendantsTreeCoclassOne} (see{\nobreakspace}\texttt{PqDescendantsTreeCoclassOne} (\ref{PqDescendantsTreeCoclassOne})) during an \textsf{XGAP} session, it contains the \textsf{XGAP} \texttt{Sheet} on which the descendants tree is displayed; and 
\item[{\texttt{nextX}}]  if set by a call to \texttt{PqDescendantsTreeCoclassOne} (see{\nobreakspace}\texttt{PqDescendantsTreeCoclassOne} (\ref{PqDescendantsTreeCoclassOne})) during an \textsf{XGAP} session, it contains a list of integers, the \mbox{\texttt{\mdseries\slshape i}}th entry of which is the \mbox{\texttt{\mdseries\slshape x}}\texttt{\symbol{45}}coordinate of the next node (representing a descendant)
for the \mbox{\texttt{\mdseries\slshape i}}th class. 
\end{description}
 }

 

\subsection{\textcolor{Chapter }{ANUPQDirectoryTemporary}}
\logpage{[ 3, 2, 2 ]}\nobreak
\hyperdef{L}{X7FBB2F457E4BD6AB}{}
{\noindent\textcolor{FuncColor}{$\triangleright$\enspace\texttt{ANUPQDirectoryTemporary({\mdseries\slshape dir})\index{ANUPQDirectoryTemporary@\texttt{ANUPQDirectoryTemporary}}
\label{ANUPQDirectoryTemporary}
}\hfill{\scriptsize (function)}}\\


 calls the UNIX command \texttt{mkdir} to create \mbox{\texttt{\mdseries\slshape dir}}, which must be a string, and if successful a directory object for \mbox{\texttt{\mdseries\slshape dir}} is both assigned to \texttt{ANUPQData.tmpdir} and returned. The field \texttt{ANUPQData.outfile} is also set to be a file in \texttt{ANUPQData.tmpdir}, and on exit from \textsf{GAP} \mbox{\texttt{\mdseries\slshape dir}} is removed. Most users will never need this command; by default, \textsf{GAP} typically chooses a ``random'' subdirectory of \texttt{/tmp} for \texttt{ANUPQData.tmpdir} which may occasionally have limits on what may be written there. \texttt{ANUPQDirectoryTemporary} permits the user to choose a directory (object) where one is not so limited. }

 }

 
\section{\textcolor{Chapter }{Setting the Verbosity of ANUPQ via Info and InfoANUPQ}}\label{Setting the Verbosity of ANUPQ via Info and InfoANUPQ}
\logpage{[ 3, 3, 0 ]}
\hyperdef{L}{X83E5C7CF7D8739DF}{}
{
  

\subsection{\textcolor{Chapter }{InfoANUPQ}}
\logpage{[ 3, 3, 1 ]}\nobreak
\hyperdef{L}{X7CBC9B458497BFF1}{}
{\noindent\textcolor{FuncColor}{$\triangleright$\enspace\texttt{InfoANUPQ\index{InfoANUPQ@\texttt{InfoANUPQ}}
\label{InfoANUPQ}
}\hfill{\scriptsize (info class)}}\\


 The input to and the output from the \texttt{pq} program is, by default, not displayed. However the user may choose to see
some, or all, of this input/output. This is done via the \texttt{Info} mechanism (see Section{\nobreakspace} \textbf{Reference: Info Functions} in the \textsf{GAP} Reference Manual). For this purpose, there is the \mbox{\texttt{\mdseries\slshape InfoClass}} \texttt{InfoANUPQ}. If the \texttt{InfoLevel} of \texttt{InfoANUPQ} is high enough each line of \texttt{pq} input/output is directed to a call to \texttt{Info} and will be displayed for the user to see. By default, the \texttt{InfoLevel} of \texttt{InfoANUPQ} is 1, and it is recommended that you leave it at this level, or higher.
Messages that the user should presumably want to see and output from the \texttt{pq} program influenced by the value of the option \texttt{OutputLevel} (see the options listed in Section{\nobreakspace}\texttt{Pq} (\ref{Pq})), other than timing and memory usage are directed to \texttt{Info} at \texttt{InfoANUPQ} level 1. 

 To turn off \emph{all} \texttt{InfoANUPQ} messaging, set the \texttt{InfoANUPQ} level to 0. 

 There are five other user\texttt{\symbol{45}}intended \texttt{InfoANUPQ} levels: 2, 3, 4, 5 and 6. 
\begin{Verbatim}[commandchars=!@|,fontsize=\small,frame=single,label=Example]
  !gapprompt@gap>| !gapinput@SetInfoLevel(InfoANUPQ, 2);|
\end{Verbatim}
 enables the display of most timing and memory usage data from the \texttt{pq} program, and also the number of identity instances when the \texttt{Identities} option is used. (Some timing and memory usage data, particularly when profuse
in quantity, is \texttt{Info}\texttt{\symbol{45}}ed at \texttt{InfoANUPQ} level 3 instead.) Note that the the \textsf{GAP} functions \texttt{time} and \texttt{Runtime} (see{\nobreakspace}\texttt{Runtime} (\textbf{Reference: Runtime}) in the \textsf{GAP} Reference Manual) count the time spent by \textsf{GAP} and \emph{not} the time spent by the (external) \texttt{pq} program. 
\begin{Verbatim}[commandchars=!@|,fontsize=\small,frame=single,label=Example]
  !gapprompt@gap>| !gapinput@SetInfoLevel(InfoANUPQ, 3);|
\end{Verbatim}
 enables the display of output of the nature of the first two \texttt{InfoANUPQ} that was not directly invoked by the user (e.g.{\nobreakspace}some commands
require \textsf{GAP} to discover something about the current state known to the \texttt{pq} program). The identity instances processed under the \texttt{Identities} option are also displayed at this level. In some cases, the \texttt{pq} program produces a lot of output despite the fact that the \texttt{OutputLevel} (see{\nobreakspace}\ref{option OutputLevel}) is unset or is set to 0; such output is also \texttt{Info}\texttt{\symbol{45}}ed at \texttt{InfoANUPQ} level 3. 
\begin{Verbatim}[commandchars=!@|,fontsize=\small,frame=single,label=Example]
  !gapprompt@gap>| !gapinput@SetInfoLevel(InfoANUPQ, 4);|
\end{Verbatim}
 enables the display of all the commands directed to the \texttt{pq} program, behind a ``\texttt{ToPQ{\textgreater} }'' prompt (so that you can distinguish it from the output from the \texttt{pq} program). See Section{\nobreakspace}\hyperref[Hints and Warnings regarding the use of Options]{`Hints and Warnings regarding the use of Options'} for an example of how this can be a useful troubleshooting tool. 
\begin{Verbatim}[commandchars=!@|,fontsize=\small,frame=single,label=Example]
  !gapprompt@gap>| !gapinput@SetInfoLevel(InfoANUPQ, 5);|
\end{Verbatim}
 enables the display of the \texttt{pq} program's prompts for input. Finally, 
\begin{Verbatim}[commandchars=!@|,fontsize=\small,frame=single,label=Example]
  !gapprompt@gap>| !gapinput@SetInfoLevel(InfoANUPQ, 6);|
\end{Verbatim}
 enables the display of all other output from the \texttt{pq} program, namely the banner and menus. However, the timing data printed when
the \texttt{pq} program exits can never be observed. }

 }

 
\section{\textcolor{Chapter }{Utility Functions}}\label{Utility-Functions}
\logpage{[ 3, 4, 0 ]}
\hyperdef{L}{X810FFB1C8035C8BE}{}
{
  

\subsection{\textcolor{Chapter }{PqLeftNormComm}}
\logpage{[ 3, 4, 1 ]}\nobreak
\hyperdef{L}{X8771393B7F53F534}{}
{\noindent\textcolor{FuncColor}{$\triangleright$\enspace\texttt{PqLeftNormComm({\mdseries\slshape elts})\index{PqLeftNormComm@\texttt{PqLeftNormComm}}
\label{PqLeftNormComm}
}\hfill{\scriptsize (function)}}\\


 returns for a list of elements of some group (e.g.{\nobreakspace}\mbox{\texttt{\mdseries\slshape elts}} may be a list of words in the generators of a free or fp group) the left
normed commutator of \mbox{\texttt{\mdseries\slshape elts}}, e.g.{\nobreakspace}if \mbox{\texttt{\mdseries\slshape w1}}, \mbox{\texttt{\mdseries\slshape w2}}, \mbox{\texttt{\mdseries\slshape w3}} are such elements then \texttt{PqLeftNormComm( [\mbox{\texttt{\mdseries\slshape w1}}, \mbox{\texttt{\mdseries\slshape w2}}, \mbox{\texttt{\mdseries\slshape w3}}] );} is equivalent to \texttt{Comm( Comm( \mbox{\texttt{\mdseries\slshape w1}}, \mbox{\texttt{\mdseries\slshape w2}} ), \mbox{\texttt{\mdseries\slshape w3}} );}. 

 \emph{Note:} \mbox{\texttt{\mdseries\slshape elts}} must contain at least two elements. }

 

\subsection{\textcolor{Chapter }{PqGAPRelators}}
\logpage{[ 3, 4, 2 ]}\nobreak
\hyperdef{L}{X7A567432879510A6}{}
{\noindent\textcolor{FuncColor}{$\triangleright$\enspace\texttt{PqGAPRelators({\mdseries\slshape group, rels})\index{PqGAPRelators@\texttt{PqGAPRelators}}
\label{PqGAPRelators}
}\hfill{\scriptsize (function)}}\\


 returns a list of words that \textsf{GAP} understands, given a list \mbox{\texttt{\mdseries\slshape rels}} of strings in the string representations of the generators of the fp group \mbox{\texttt{\mdseries\slshape group}} prepared as a list of relators for the \texttt{pq} program. 

 \emph{Note:} The \texttt{pq} program does not use \texttt{/} to indicate multiplication by an inverse and uses square brackets to represent
(left normed) commutators. Also, even though the \texttt{pq} program accepts relations, all elements of \mbox{\texttt{\mdseries\slshape rels}} \emph{must} be in relator form, i.e.{\nobreakspace}a relation of form \texttt{\mbox{\texttt{\mdseries\slshape w1}} = \mbox{\texttt{\mdseries\slshape w2}}} must be written as \texttt{\mbox{\texttt{\mdseries\slshape w1}}*(\mbox{\texttt{\mdseries\slshape w2}})\texttt{\symbol{94}}\texttt{\symbol{45}}1}. 

 Here is an example: 
\begin{Verbatim}[commandchars=!@|,fontsize=\small,frame=single,label=Example]
  !gapprompt@gap>| !gapinput@F := FreeGroup("a", "b");|
  <free group on the generators [ a, b ]>
  !gapprompt@gap>| !gapinput@PqGAPRelators(F, [ "a*b^2", "[a,b]^2*a", "[a,b,a,b]^a" ]);|
  [ a*b^2, (a^-1*b^-1*a*b)^2*a, 
    a^-2*b^-1*a^-1*b*(a*b^-1)^2*a^-1*b*a^-1*b^-1*(a*b)^2*a ]
\end{Verbatim}
 }

 

\subsection{\textcolor{Chapter }{PqParseWord}}
\logpage{[ 3, 4, 3 ]}\nobreak
\hyperdef{L}{X7F3C5D1C7EC36EAE}{}
{\noindent\textcolor{FuncColor}{$\triangleright$\enspace\texttt{PqParseWord({\mdseries\slshape word, n})\index{PqParseWord@\texttt{PqParseWord}}
\label{PqParseWord}
}\hfill{\scriptsize (function)}}\\


 parses a \mbox{\texttt{\mdseries\slshape word}}, a string representing a word in the pc generators \texttt{x1,...,x\mbox{\texttt{\mdseries\slshape n}}}, through \textsf{GAP}. This function is provided as a
rough\texttt{\symbol{45}}and\texttt{\symbol{45}}ready check of \mbox{\texttt{\mdseries\slshape word}} for syntax errors. A syntax error will cause the entering of a \texttt{break}\texttt{\symbol{45}}loop, in which the error message may or may not be
meaningful (depending on whether the syntax error gets caught at the \textsf{GAP} or kernel level). 

 \emph{Note:} The reason the generators \emph{must} be \texttt{x1,...,x\mbox{\texttt{\mdseries\slshape n}}} is that these are the pc generator names used by the \texttt{pq} program (as distinct from the generator names for the group provided by the
user to a function like \texttt{Pq} that invokes the \texttt{pq} program). }

 

\subsection{\textcolor{Chapter }{PqExample (no arguments)}}
\logpage{[ 3, 4, 4 ]}\nobreak
\hyperdef{L}{X7BB0EB607F337265}{}
{\noindent\textcolor{FuncColor}{$\triangleright$\enspace\texttt{PqExample({\mdseries\slshape })\index{PqExample@\texttt{PqExample}!no arguments}
\label{PqExample:no arguments}
}\hfill{\scriptsize (function)}}\\
\noindent\textcolor{FuncColor}{$\triangleright$\enspace\texttt{PqExample({\mdseries\slshape example[, PqStart][, Display]})\index{PqExample@\texttt{PqExample}}
\label{PqExample}
}\hfill{\scriptsize (function)}}\\
\noindent\textcolor{FuncColor}{$\triangleright$\enspace\texttt{PqExample({\mdseries\slshape example[, PqStart][, filename]})\index{PqExample@\texttt{PqExample}!with filename}
\label{PqExample:with filename}
}\hfill{\scriptsize (function)}}\\


 With no arguments, or with single argument \texttt{"index"}, or a string \mbox{\texttt{\mdseries\slshape example}} that is not the name of a file in the \texttt{examples} directory, an index of available examples is displayed. 

 With just the one argument \mbox{\texttt{\mdseries\slshape example}} that is the name of a file in the \texttt{examples} directory, the example contained in that file is executed in its simplest
form. Some examples accept options which you may use to modify some of the
options used in the commands of the example. To find out which options an
example accepts, use one of the mechanisms for displaying the example
described below. 

 Some examples have both non\texttt{\symbol{45}}interactive and interactive
forms; those that are non\texttt{\symbol{45}}interactive only have a name
ending in \texttt{\texttt{\symbol{45}}ni}; those that are interactive only have a name ending in \texttt{\texttt{\symbol{45}}i}; examples with names ending in \texttt{.g} also have only one form; all other examples have both
non\texttt{\symbol{45}}interactive and interactive forms and for these giving \texttt{PqStart} as second argument invokes \texttt{PqStart} initially and makes the appropriate adjustments so that the example is
executed or displayed using interactive functions. 

 If \texttt{PqExample} is called with last (second or third) argument \texttt{Display} then the example is displayed without being executed. If the last argument is
a non\texttt{\symbol{45}}empty string \mbox{\texttt{\mdseries\slshape filename}} then the example is also displayed without being executed but is also written
to a file with that name. Passing an empty string as last argument has the
same effect as passing \texttt{Display}. 

 \emph{Note:} The variables used in \texttt{PqExample} are local to the running of \texttt{PqExample}, so there's no danger of having some of your variables
over\texttt{\symbol{45}}written. However, they are not completely lost either.
They are saved to a record \texttt{ANUPQData.examples.vars}, i.e.{\nobreakspace}if \texttt{F} is a variable used in the example then you will be able to access it after \texttt{PqExample} has finished as \texttt{ANUPQData.examples.vars.F}. }

 

\subsection{\textcolor{Chapter }{AllPqExamples}}
\logpage{[ 3, 4, 5 ]}\nobreak
\hyperdef{L}{X823C93FC7B87F5BC}{}
{\noindent\textcolor{FuncColor}{$\triangleright$\enspace\texttt{AllPqExamples({\mdseries\slshape })\index{AllPqExamples@\texttt{AllPqExamples}}
\label{AllPqExamples}
}\hfill{\scriptsize (function)}}\\


 returns a list of all currently available examples in default
UNIX\texttt{\symbol{45}}listing (i.e.{\nobreakspace}alphabetic) order. }

 

\subsection{\textcolor{Chapter }{GrepPqExamples}}
\logpage{[ 3, 4, 6 ]}\nobreak
\hyperdef{L}{X7E3E4B047DC2E323}{}
{\noindent\textcolor{FuncColor}{$\triangleright$\enspace\texttt{GrepPqExamples({\mdseries\slshape string})\index{GrepPqExamples@\texttt{GrepPqExamples}}
\label{GrepPqExamples}
}\hfill{\scriptsize (function)}}\\


 runs the UNIX command \texttt{grep \mbox{\texttt{\mdseries\slshape string}}} over the \textsf{ANUPQ} examples and returns the list of examples for which there is a match. The
actual matches are \texttt{Info}\texttt{\symbol{45}}ed at \texttt{InfoANUPQ} level 2. }

 

\subsection{\textcolor{Chapter }{ToPQLog}}
\logpage{[ 3, 4, 7 ]}\nobreak
\hyperdef{L}{X8104B2BA872EFCCB}{}
{\noindent\textcolor{FuncColor}{$\triangleright$\enspace\texttt{ToPQLog({\mdseries\slshape [filename]})\index{ToPQLog@\texttt{ToPQLog}}
\label{ToPQLog}
}\hfill{\scriptsize (function)}}\\


 With string argument \mbox{\texttt{\mdseries\slshape filename}}, \texttt{ToPQLog} opens the file with name \mbox{\texttt{\mdseries\slshape filename}} for logging; all commands written to the \texttt{pq} binary (that are \texttt{Info}\texttt{\symbol{45}}ed behind a ``\texttt{ToPQ{\textgreater} }'' prompt at \texttt{InfoANUPQ} level 4) are then also written to that file (but without prompts). With no
argument, \texttt{ToPQLog} stops logging to whatever file was being logged to. If a file was already
being logged to, that file is closed and the file with name \mbox{\texttt{\mdseries\slshape filename}} is opened for logging. }

 }

 
\section{\textcolor{Chapter }{Attributes and a Property for fp and pc p\texttt{\symbol{45}}groups}}\label{Attributes and a Property for fp and pc p-groups}
\logpage{[ 3, 5, 0 ]}
\hyperdef{L}{X818175EF85CAA807}{}
{
  

\subsection{\textcolor{Chapter }{NuclearRank}}
\logpage{[ 3, 5, 1 ]}\nobreak
\hyperdef{L}{X87DC922A78EB0DD6}{}
{\noindent\textcolor{FuncColor}{$\triangleright$\enspace\texttt{NuclearRank({\mdseries\slshape G})\index{NuclearRank@\texttt{NuclearRank}}
\label{NuclearRank}
}\hfill{\scriptsize (attribute)}}\\
\noindent\textcolor{FuncColor}{$\triangleright$\enspace\texttt{MultiplicatorRank({\mdseries\slshape G})\index{MultiplicatorRank@\texttt{MultiplicatorRank}}
\label{MultiplicatorRank}
}\hfill{\scriptsize (attribute)}}\\
\noindent\textcolor{FuncColor}{$\triangleright$\enspace\texttt{IsCapable({\mdseries\slshape G})\index{IsCapable@\texttt{IsCapable}}
\label{IsCapable}
}\hfill{\scriptsize (property)}}\\


 return the nuclear rank of \mbox{\texttt{\mdseries\slshape G}}, $p$\texttt{\symbol{45}}multiplicator rank of \mbox{\texttt{\mdseries\slshape G}}, and whether \mbox{\texttt{\mdseries\slshape G}} is capable (i.e.{\nobreakspace}\texttt{true} if it is, or \texttt{false} if it is not), respectively. 

 These attributes and property are set automatically if \mbox{\texttt{\mdseries\slshape G}} is one of the following: 
\begin{itemize}
\item  an fp group returned by \texttt{PqStandardPresentation} or \texttt{StandardPresentation} (see{\nobreakspace}\texttt{PqStandardPresentation} (\ref{PqStandardPresentation})); 
\item  the image (fp group) of the epimorphism returned by an \texttt{EpimorphismPqStandardPresentation} or \texttt{EpimorphismStandardPresentation} call (see{\nobreakspace}\texttt{EpimorphismPqStandardPresentation} (\ref{EpimorphismPqStandardPresentation})); or 
\item  one of the pc groups of the list of descendants returned by \texttt{PqDescendants} (see{\nobreakspace}\texttt{PqDescendants} (\ref{PqDescendants})). 
\end{itemize}
 If \mbox{\texttt{\mdseries\slshape G}} is an fp group or a pc $p$\texttt{\symbol{45}}group and not one of the above and the attribute or
property has not otherwise been set for \mbox{\texttt{\mdseries\slshape G}}, then \texttt{PqStandardPresentation} is called to set all three of \texttt{NuclearRank}, \texttt{MultiplicatorRank} and \texttt{IsCapable}, before returning the value of the attribute or property actually called.
Such a group \mbox{\texttt{\mdseries\slshape G}} must know in advance that it is a $p$\texttt{\symbol{45}}group; this is the case for the groups returned by the
functions \texttt{Pq} and \texttt{PqPCover}, and the image group of the epimorphism returned by \texttt{PqEpimorphism}. Otherwise, if you know the group to be a $p$\texttt{\symbol{45}}group, then this can be set by typing 
\begin{Verbatim}[commandchars=!@|,fontsize=\small,frame=single,label=]
  SetIsPGroup( G, true );
\end{Verbatim}
 or by invoking \texttt{IsPGroup( \mbox{\texttt{\mdseries\slshape G}} )}. Note that for an fp group \mbox{\texttt{\mdseries\slshape G}}, the latter may result in a coset enumeration which might not terminate in a
reasonable time. 

 \emph{Note:} For \mbox{\texttt{\mdseries\slshape G}} such that \texttt{HasNuclearRank(\mbox{\texttt{\mdseries\slshape G}}) = true}, \texttt{IsCapable(\mbox{\texttt{\mdseries\slshape G}})} is equivalent to (the truth or falsity of) \texttt{NuclearRank( \mbox{\texttt{\mdseries\slshape G}} ) = 0}. }

 }

 
\section{\textcolor{Chapter }{Hints and Warnings regarding the use of Options}}\label{Hints and Warnings regarding the use of Options}
\logpage{[ 3, 6, 0 ]}
\hyperdef{L}{X7BA20FA07B166B37}{}
{
  On a first reading we recommend you skip this section and come back to it if
and when you run into trouble. 

 \index{menu item!of pq program} \index{option!of pq program is a menu item}  \emph{Note:} By ``options'' we refer to \textsf{GAP} options. The \texttt{pq} program also uses the term ``option''; to distinguish the two usages of ``option'', in this manual we use the term \emph{menu item} to refer to what the \texttt{pq} program refers to as an ``option''. 

 Options are passed to the \textsf{ANUPQ} interface functions in either of the two usual mechanisms provided by \textsf{GAP}, namely: 
\begin{itemize}
\item  options may be set globally using the function \texttt{PushOptions} (see Chapter{\nobreakspace} \textbf{Reference: Options Stack} in the \textsf{GAP} Reference Manual); or 
\item  options may be appended to the argument list of any function call, separated
by a colon from the argument list (see Chapter{\nobreakspace} \textbf{Reference: Function Calls} in the \textsf{GAP} Reference Manual), in which case they are then passed on recursively to any
subsequent inner function call, which may in turn have options of their own. 
\end{itemize}
 Particularly, when one is using the interactive functions of
Chapter{\nobreakspace}\hyperref[Interactive ANUPQ functions]{`Interactive ANUPQ functions'}, one should, in general, avoid using the global method of passing options. In
fact, it is recommended that prior to calling \texttt{PqStart} the \texttt{OptionsStack} be empty. The essential problem with setting options globally using the
function \texttt{PushOptions} is that options pushed onto \texttt{OptionsStack}, in this way, (generally) remain there until an explicit \texttt{PopOptions()} call is made. 

 In contrast, options passed in the usual way behind a colon following a
function's arguments (see  \textbf{Reference: Function Call With Options} in the \textsf{GAP} Reference Manual) are local, and disappear from \texttt{OptionsStack} after the function has executed successfully. If the function does \emph{not} execute successfully, i.e.{\nobreakspace}it runs into error and the user \texttt{quit}s the resulting \texttt{break} loop (see Section{\nobreakspace} \textbf{Reference: Break Loops} in the Reference Manual) rather than attempting to repair the problem and
typing \texttt{return;} then, unless the error at the kernel level, the \texttt{OptionsStack} is reset. If an error is detected inside the kernel (hopefully, this should
occur only rarely, if at all) then the options of that function will \emph{not} be cleared from \texttt{OptionsStack}; in such cases: 
\begin{Verbatim}[commandchars=!@|,fontsize=\small,frame=single,label=Example]
  !gapprompt@gap>| !gapinput@ResetOptionsStack();|
  #I  Options stack is already empty
\end{Verbatim}
 is usually necessary (see Chapter{\nobreakspace}\texttt{ResetOptionsStack} (\textbf{Reference: ResetOptionsStack}) in the \textsf{GAP} Reference Manual), which recursively calls \texttt{PopOptions()} until \texttt{OptionsStack} is empty, or as in the above case warns you that the \texttt{OptionsStack} is already empty. 

 Note that a function, that is passed options after the colon, will also see
any global options or any options passed down recursively from functions
calling that function, unless those options are over\texttt{\symbol{45}}ridden
by options passed via the function. Also, note that duplication of option
names for different programs may lead to misinterpretations, and
mis\texttt{\symbol{45}}spelled options will not be ``seen''. 

 The non\texttt{\symbol{45}}interactive functions of Chapter{\nobreakspace}\hyperref[non-interact]{`Non\texttt{\symbol{45}}interactive ANUPQ functions'} that have \texttt{Pq} somewhere in their name provide an alternative method of passing options as
additional arguments. This has the advantages that options can be abbreviated
and mis\texttt{\symbol{45}}spelled options will be trapped. \index{troubleshooting tips} 

\subsection{\textcolor{Chapter }{ANUPQWarnOfOtherOptions}}
\logpage{[ 3, 6, 1 ]}\nobreak
\hyperdef{L}{X81F4AAE084C34B4B}{}
{\noindent\textcolor{FuncColor}{$\triangleright$\enspace\texttt{ANUPQWarnOfOtherOptions\index{ANUPQWarnOfOtherOptions@\texttt{ANUPQWarnOfOtherOptions}}
\label{ANUPQWarnOfOtherOptions}
}\hfill{\scriptsize (global variable)}}\\


 is a global variable that is by default \texttt{false}. If it is set to \texttt{true} then any function provided by the \textsf{ANUPQ} function that recognises at least one option, will warn you of ``other'' options, i.e.{\nobreakspace}options that the function does not recognise.
These warnings are emitted at \texttt{InfoWarning} or \texttt{InfoANUPQ} level 1. This is useful for detecting mis\texttt{\symbol{45}}spelled options.
Here is an example using the function \texttt{Pq} (first described in Chapter{\nobreakspace}\hyperref[non-interact]{`Non\texttt{\symbol{45}}interactive ANUPQ functions'}): 
\begin{Verbatim}[commandchars=!@|,fontsize=\small,frame=single,label=Example]
  !gapprompt@gap>| !gapinput@SetInfoLevel(InfoANUPQ, 1);        # Set InfoANUPQ to default level|
  !gapprompt@gap>| !gapinput@ANUPQWarnOfOtherOptions := true;;|
  !gapprompt@gap>| !gapinput@# The following makes entry into break loops very ``quiet'' ...|
  !gapprompt@gap>| !gapinput@OnBreak := function() Where(0); end;;|
  !gapprompt@gap>| !gapinput@F := FreeGroup( "a", "b" );|
  <free group on the generators [ a, b ]>
  !gapprompt@gap>| !gapinput@Pq( F : Prime := 2, Classbound := 1 );|
  #I  ANUPQ Warning: Options: [ "Classbound" ] ignored
  #I  (invalid for generic function: `Pq').
  user interrupt at
  moreOfline := ReadLine( iostream );
  Entering break read-eval-print loop ...
  you can 'quit;' to quit to outer loop, or
  you can 'return;' to continue
\end{Verbatim}
 Here we mistyped \texttt{ClassBound} as \texttt{Classbound}, and after seeing the \texttt{Info}\texttt{\symbol{45}}ed warning that \texttt{Classbound} was ignored, we typed a \mbox{\texttt{\mdseries\slshape control}}\texttt{\symbol{45}}C (that's the ``\texttt{user interrupt at}'' message) which took us into a break loop. Since the \texttt{Pq} command was not able to finish, the options \texttt{Prime} and \texttt{Classbound}, in particular, will still be on the \texttt{OptionsStack}: 
\begin{Verbatim}[commandchars=!@|,fontsize=\small,frame=single,label=Example]
  !gapbrkprompt@brk>| !gapinput@OptionsStack;|
  [ rec( Prime := 2, Classbound := 1 ), 
    rec( Prime := 2, Classbound := 1, PqEpiOrPCover := "pQuotient" ) ]
\end{Verbatim}
 The option \texttt{PqEpiOrPCover} is a behind\texttt{\symbol{45}}the\texttt{\symbol{45}}scenes option that need
not concern the user. On \texttt{quit}ting the \texttt{break}\texttt{\symbol{45}}loop the \texttt{OptionsStack} is reset and a warning telling you this is emitted: 
\begin{Verbatim}[commandchars=!@|,fontsize=\small,frame=single,label=Example]
  !gapbrkprompt@brk>| !gapinput@quit; # to get back to the `gap>' prompt|
  #I  Options stack has been reset
\end{Verbatim}
 Above, we altered \texttt{OnBreak} (see{\nobreakspace}\texttt{OnBreak} (\textbf{Reference: OnBreak}) in the Reference manual) to reduce the back\texttt{\symbol{45}}tracing on
entry into a break loop. We now restore \texttt{OnBreak} to its usual value. 
\begin{Verbatim}[commandchars=!@|,fontsize=\small,frame=single,label=Example]
  !gapprompt@gap>| !gapinput@OnBreak := Where;;|
\end{Verbatim}
 \emph{Notes} 

 In cases where functions recursively call others with options
(e.g.{\nobreakspace}when using \texttt{PqExample} with options), setting \texttt{ANUPQWarnOfOtherOptions := true} may give rise to spurious ``other'' option detections. 

 It is recommended that the novice user set \texttt{ANUPQWarnOfOtherOptions} to \texttt{true} in their \texttt{gap.ini} file (see Section{\nobreakspace}\hyperref[Loading the ANUPQ Package]{`Loading the ANUPQ Package'}). }

 \emph{Other Troubleshooting Strategies} 

 There are some other strategies which may have helped us to see our error
above. The function \texttt{Pq} recognises the option \texttt{OutputLevel} (see{\nobreakspace}\ref{option OutputLevel}); if this option is set to at least 1, the \texttt{pq} program provides information on each class quotient as it is generated: 
\begin{Verbatim}[commandchars=!@|,fontsize=\small,frame=single,label=Example]
  !gapprompt@gap>| !gapinput@ANUPQWarnOfOtherOptions := false;; # Set back to normal|
  !gapprompt@gap>| !gapinput@F := FreeGroup( "a", "b" );;|
  !gapprompt@gap>| !gapinput@Pq( F : Prime := 2, Classbound := 1, OutputLevel := 1 ); |
  #I  Lower exponent-2 central series for [grp]
  #I  Group: [grp] to lower exponent-2 central class 1 has order 2^2
  #I  Group: [grp] to lower exponent-2 central class 2 has order 2^5
  #I  Group: [grp] to lower exponent-2 central class 3 has order 2^10
  #I  Group: [grp] to lower exponent-2 central class 4 has order 2^18
  #I  Group: [grp] to lower exponent-2 central class 5 has order 2^32
  #I  Group: [grp] to lower exponent-2 central class 6 has order 2^55
  #I  Group: [grp] to lower exponent-2 central class 7 has order 2^96
  #I  Group: [grp] to lower exponent-2 central class 8 has order 2^167
  #I  Group: [grp] to lower exponent-2 central class 9 has order 2^294
  #I  Group: [grp] to lower exponent-2 central class 10 has order 2^520
  #I  Group: [grp] to lower exponent-2 central class 11 has order 2^932
  #I  Group: [grp] to lower exponent-2 central class 12 has order 2^1679
  [... output truncated ...]
\end{Verbatim}
 After seeing the information for the class 2 quotient we may have got the idea
that the \texttt{Classbound} option was not recognised and may have realised that this was due to a
mis\texttt{\symbol{45}}spelling. The above will ordinarily cause the available
space to be exhausted, necessitating user\texttt{\symbol{45}}intervention by
typing \mbox{\texttt{\mdseries\slshape control}}\texttt{\symbol{45}}C and \texttt{quit;} (to escape the break loop); otherwise \texttt{Pq} terminates when the class reaches 63 (the default value of \texttt{ClassBound}). 

 If you have some familiarity with ``keyword'' command input to the \texttt{pq} binary, then setting the level of \texttt{InfoANUPQ} to 4 would also have indicated a problem: 
\begin{Verbatim}[commandchars=!@|,fontsize=\small,frame=single,label=Example]
  !gapprompt@gap>| !gapinput@ResetOptionsStack(); # Necessary, if a break-loop was entered above|
  !gapprompt@gap>| !gapinput@SetInfoLevel(InfoANUPQ, 4);|
  !gapprompt@gap>| !gapinput@Pq( F : Prime := 2, Classbound := 1 );|
  #I  ToPQ> 7  #to (Main) p-Quotient Menu
  #I  ToPQ> 1  #define group
  #I  ToPQ> name [grp]
  #I  ToPQ> prime 2
  #I  ToPQ> class 63
  #I  ToPQ> exponent 0
  #I  ToPQ> output 0
  #I  ToPQ> generators { a,b }
  #I  ToPQ> relators   {  };
  [... output truncated ...]
\end{Verbatim}
 Here the line ``\texttt{\#I ToPQ{\textgreater} class 63}'' indicates that a directive to set the classbound to 63 was sent to the \texttt{pq} program. }

 }

 
\chapter{\textcolor{Chapter }{Non\texttt{\symbol{45}}interactive ANUPQ functions}}\label{non-interact}
\logpage{[ 4, 0, 0 ]}
\hyperdef{L}{X7C51F26D839279FF}{}
{
  Here we describe all the non\texttt{\symbol{45}}interactive functions of the \textsf{ANUPQ} package; i.e.{\nobreakspace}``one\texttt{\symbol{45}}shot'' functions that invoke the \texttt{pq} program in such a way that once \textsf{GAP} has got what it needs, the \texttt{pq} program is allowed to exit. It is expected that most of the time users will
only need these functions. The functions interface with three of the four
algorithms (see Chapter{\nobreakspace}\hyperref[Introduction]{`Introduction'}) provided by the ANU \texttt{pq} C program, and are mainly grouped according to the algorithm of the \texttt{pq} program they relate to. 

 In Section{\nobreakspace}\hyperref[Computing p-Quotients]{`Computing p\texttt{\symbol{45}}Quotients'}, we describe the functions that give access to the $p$\texttt{\symbol{45}}quotient algorithm. 

 Section{\nobreakspace}\hyperref[Computing Standard Presentations]{`Computing Standard Presentations'} describe functions that give access to the standard presentation algorithm. 

 Section{\nobreakspace}\hyperref[Testing p-Groups for Isomorphism]{`Testing p\texttt{\symbol{45}}Groups for Isomorphism'} describe functions that implement an isomorphism test for $p$\texttt{\symbol{45}}groups using the standard presentation algorithm. 

 In Section{\nobreakspace}\hyperref[Computing Descendants of a p-Group]{`Computing Descendants of a p\texttt{\symbol{45}}Group'}, we describe functions that give access to the $p$\texttt{\symbol{45}}group generation algorithm. 

 To use any of the functions one must have at some stage previously typed: 
\begin{Verbatim}[commandchars=!@|,fontsize=\small,frame=single,label=Example]
  !gapprompt@gap>| !gapinput@LoadPackage("anupq");|
\end{Verbatim}
 (the response of which we have omitted; see{\nobreakspace}\hyperref[Loading the ANUPQ Package]{`Loading the ANUPQ Package'}). 

 It is strongly recommended that the user try the examples provided. To save
typing there is a \texttt{PqExample} equivalent for each manual example. We also suggest that to start with you may
find the examples more instructive if you set the \texttt{InfoANUPQ} level to 2 (see{\nobreakspace}\texttt{InfoANUPQ} (\ref{InfoANUPQ})). 
\section{\textcolor{Chapter }{Computing p\texttt{\symbol{45}}Quotients}}\label{Computing p-Quotients}
\logpage{[ 4, 1, 0 ]}
\hyperdef{L}{X80406BD47EB186E9}{}
{
  

\subsection{\textcolor{Chapter }{Pq}}
\logpage{[ 4, 1, 1 ]}\nobreak
\hyperdef{L}{X86DD32D5803CF2C8}{}
{\noindent\textcolor{FuncColor}{$\triangleright$\enspace\texttt{Pq({\mdseries\slshape F: options})\index{Pq@\texttt{Pq}}
\label{Pq}
}\hfill{\scriptsize (function)}}\\


 returns for the fp or pc group \mbox{\texttt{\mdseries\slshape F}}, the $p$\texttt{\symbol{45}}quotient of \mbox{\texttt{\mdseries\slshape F}} specified by \mbox{\texttt{\mdseries\slshape options}}, as a pc group. Following the colon, \mbox{\texttt{\mdseries\slshape options}} is a selection of the options from the following list, separated by commas
like record components (see Section{\nobreakspace} \textbf{Reference: Function Call With Options} in the \textsf{GAP} Reference Manual). As a minimum the user \emph{must} supply a value for the \texttt{Prime} option. Below we list the options recognised by \texttt{Pq} (see Chapter{\nobreakspace}\hyperref[ANUPQ Options]{`ANUPQ Options'} for detailed descriptions). 
\begin{itemize}
\item  \texttt{Prime := \mbox{\texttt{\mdseries\slshape p}}}\index{option Prime} 
\item  \texttt{ClassBound := \mbox{\texttt{\mdseries\slshape n}}}\index{option ClassBound} 
\item  \texttt{Exponent := \mbox{\texttt{\mdseries\slshape n}}}\index{option Exponent} 
\item  \texttt{Relators := \mbox{\texttt{\mdseries\slshape rels}}}\index{option Relators} 
\item  \texttt{Metabelian}\index{option Metabelian} 
\item  \texttt{Identities := \mbox{\texttt{\mdseries\slshape funcs}}}\index{option Identities} 
\item  \texttt{GroupName := \mbox{\texttt{\mdseries\slshape name}}}\index{option GroupName} 
\item  \texttt{OutputLevel := \mbox{\texttt{\mdseries\slshape n}}}\index{option OutputLevel} 
\item  \texttt{SetupFile := \mbox{\texttt{\mdseries\slshape filename}}}\index{option SetupFile} 
\item  \texttt{PqWorkspace := \mbox{\texttt{\mdseries\slshape workspace}}}\index{option PqWorkspace} 
\end{itemize}
 \emph{Notes:} \texttt{Pq} may also be called with no arguments or one integer argument, in which case it
is being used interactively (see{\nobreakspace}\texttt{Pq} (\ref{Pq:interactive})); the same options may be used, except that \texttt{SetupFile} and \texttt{PqWorkspace} are ignored by the interactive \texttt{Pq} function. 

 See Section{\nobreakspace}\hyperref[Attributes and a Property for fp and pc p-groups]{`Attributes and a Property for fp and pc p\texttt{\symbol{45}}groups'} for the attributes and property \texttt{NuclearRank}, \texttt{MultiplicatorRank} and \texttt{IsCapable} which may be applied to the group returned by \texttt{Pq}. 

 See also \texttt{PqEpimorphism} (\texttt{PqEpimorphism} (\ref{PqEpimorphism})). 

 We now give a few examples of the use of \texttt{Pq}. Except for the addition of a few comments and the
non\texttt{\symbol{45}}suppression of output (by not using duplicated
semicolons) the next 3 examples may be run by typing: \texttt{PqExample( "Pq" );} (see{\nobreakspace}\texttt{PqExample} (\ref{PqExample})). 
\begin{Verbatim}[commandchars=!@|,fontsize=\small,frame=single,label=Example]
  !gapprompt@gap>| !gapinput@LoadPackage("anupq");; # does nothing if ANUPQ is already loaded|
  !gapprompt@gap>| !gapinput@# First we get a p-quotient of a free group of rank 2|
  !gapprompt@gap>| !gapinput@F := FreeGroup("a", "b");; a := F.1;; b := F.2;;|
  !gapprompt@gap>| !gapinput@Pq( F : Prime := 2, ClassBound := 3 ); |
  <pc group of size 1024 with 10 generators>
  !gapprompt@gap>| !gapinput@# Now let us get a p-quotient of an fp group|
  !gapprompt@gap>| !gapinput@G := F / [a^4, b^4];|
  <fp group on the generators [ a, b ]>
  !gapprompt@gap>| !gapinput@Pq( G : Prime := 2, ClassBound := 3 ); |
  <pc group of size 256 with 8 generators>
  !gapprompt@gap>| !gapinput@# Now let's get a different p-quotient of the same group|
  !gapprompt@gap>| !gapinput@Pq( G : Prime := 2, ClassBound := 3, Exponent := 4 ); |
  <pc group of size 128 with 7 generators>
  !gapprompt@gap>| !gapinput@# Now we'll get a p-quotient of another fp group|
  !gapprompt@gap>| !gapinput@# which we will redo using the `Relators' option|
  !gapprompt@gap>| !gapinput@R := [ a^25, Comm(Comm(b, a), a), b^5 ];|
  [ a^25, a^-1*b^-1*a*b*a^-1*b^-1*a^-1*b*a^2, b^5 ]
  !gapprompt@gap>| !gapinput@H := F / R;|
  <fp group on the generators [ a, b ]>
  !gapprompt@gap>| !gapinput@Pq( H : Prime := 5, ClassBound := 5, Metabelian );|
  <pc group of size 78125 with 7 generators>
\end{Verbatim}
 \index{option Relators!example of usage} Now we redo the last example to show how one may use the \texttt{Relators} option. Observe that \texttt{Comm(Comm(b, a), a)} is a left normed commutator which must be written in square bracket notation
for the \texttt{pq} program and embedded in a pair of double quotes. The function \texttt{PqGAPRelators} (see{\nobreakspace}\texttt{PqGAPRelators} (\ref{PqGAPRelators})) can be used to translate a list of strings prepared for the \texttt{Relators} option into \textsf{GAP} format. Below we use it. Observe that the value of \texttt{R} is the same as before. 
\begin{Verbatim}[commandchars=!@|,fontsize=\small,frame=single,label=Example]
  !gapprompt@gap>| !gapinput@F := FreeGroup("a", "b");;|
  !gapprompt@gap>| !gapinput@# `F' was defined for `Relators'. We use the same strings that GAP uses|
  !gapprompt@gap>| !gapinput@# for printing the free group generators. It is *not* necessary to|
  !gapprompt@gap>| !gapinput@# predefine: a := F.1; etc. (as it was above).|
  !gapprompt@gap>| !gapinput@rels := [ "a^25", "[b, a, a]", "b^5" ];|
  [ "a^25", "[b, a, a]", "b^5" ]
  !gapprompt@gap>| !gapinput@R := PqGAPRelators(F, rels);|
  [ a^25, a^-1*b^-1*a*b*a^-1*b^-1*a^-1*b*a^2, b^5 ]
  !gapprompt@gap>| !gapinput@H := F / R;|
  <fp group on the generators [ a, b ]>
  !gapprompt@gap>| !gapinput@Pq( H : Prime := 5, ClassBound := 5, Metabelian, |
  !gapprompt@>| !gapinput@           Relators := rels );|
  <pc group of size 78125 with 7 generators>
\end{Verbatim}
 In fact, above we could have just passed \texttt{F} (rather than \texttt{H}), i.e.{\nobreakspace}we could have done: 
\begin{Verbatim}[commandchars=!@|,fontsize=\small,frame=single,label=Example]
  !gapprompt@gap>| !gapinput@F := FreeGroup("a", "b");;|
  !gapprompt@gap>| !gapinput@rels := [ "a^25", "[b, a, a]", "b^5" ];|
  [ "a^25", "[b, a, a]", "b^5" ]
  !gapprompt@gap>| !gapinput@Pq( F : Prime := 5, ClassBound := 5, Metabelian, |
  !gapprompt@>| !gapinput@           Relators := rels );|
  <pc group of size 78125 with 7 generators>
\end{Verbatim}
 The non\texttt{\symbol{45}}interactive \texttt{Pq} function also allows the options to be passed in two other ways; these
alternatives have been included for those familiar with the \textsf{GAP}{\nobreakspace}3 version of the \textsf{ANUPQ} package; the preferred method of passing options is the one already described.
Firstly, they may be passed in a record as a second argument; note that any
boolean options must be set explicitly e.g. 
\begin{Verbatim}[commandchars=!@|,fontsize=\small,frame=single,label=Example]
  !gapprompt@gap>| !gapinput@Pq( H, rec( Prime := 5, ClassBound := 5, Metabelian := true ) );|
  <pc group of size 78125 with 7 generators>
\end{Verbatim}
 It is also possible to pass them as extra arguments, where each option name
appears as a string followed immediately by its value (if not a boolean
option) e.g. 
\begin{Verbatim}[commandchars=!@|,fontsize=\small,frame=single,label=Example]
  !gapprompt@gap>| !gapinput@Pq( H, "Prime", 5, "ClassBound", 5, "Metabelian" );|
  <pc group of size 78125 with 7 generators>
\end{Verbatim}
 The preceding two examples can be run from \textsf{GAP} via \texttt{PqExample( "Pq\texttt{\symbol{45}}ni" );} (see{\nobreakspace}\texttt{PqExample} (\ref{PqExample})). 

 This method of passing options permits abbreviation; the only restriction is
that the abbreviation must be unique. So \texttt{"Pr"} may be used for \texttt{"Prime"}, \texttt{"Class"} or even just \texttt{"C"} for \texttt{"ClassBound"}, etc. 

 \index{option Identities!example of usage} The following example illustrates the use of the option \texttt{Identities}. We compute the largest finite Burnside group of exponent $5$ that also satisfies the $3$\texttt{\symbol{45}}Engel identity. Each identity is defined by a function
whose arguments correspond to the variables of the identity. The return value
of each of those functions is the identity evaluated on the arguments of the
function. 
\begin{Verbatim}[commandchars=!@|,fontsize=\small,frame=single,label=Example]
  !gapprompt@gap>| !gapinput@F := FreeGroup(2);|
  <free group on the generators [ f1, f2 ]>
  !gapprompt@gap>| !gapinput@Burnside5 := x->x^5;|
  function( x ) ... end
  !gapprompt@gap>| !gapinput@Engel3 := function( x,y ) return PqLeftNormComm( [x,y,y,y] ); end;|
  function( x, y ) ... end
  !gapprompt@gap>| !gapinput@Pq( F : Prime := 5, Identities := [ Burnside5, Engel3 ] );|
  #I  Class 1 with 2 generators.
  #I  Class 2 with 3 generators.
  #I  Class 3 with 5 generators.
  #I  Class 3 with 5 generators.
  <pc group of size 3125 with 5 generators>
\end{Verbatim}
 The above example can be run from \textsf{GAP} via \texttt{PqExample(
"B5\texttt{\symbol{45}}5\texttt{\symbol{45}}Engel3\texttt{\symbol{45}}Id" );} (see{\nobreakspace}\texttt{PqExample} (\ref{PqExample})). }

 

\subsection{\textcolor{Chapter }{PqEpimorphism}}
\logpage{[ 4, 1, 2 ]}\nobreak
\hyperdef{L}{X7ADE5DDB87B99CC0}{}
{\noindent\textcolor{FuncColor}{$\triangleright$\enspace\texttt{PqEpimorphism({\mdseries\slshape F: options})\index{PqEpimorphism@\texttt{PqEpimorphism}}
\label{PqEpimorphism}
}\hfill{\scriptsize (function)}}\\


 returns for the fp or pc group \mbox{\texttt{\mdseries\slshape F}} an epimorphism from \mbox{\texttt{\mdseries\slshape F}} onto the $p$\texttt{\symbol{45}}quotient of \mbox{\texttt{\mdseries\slshape F}} specified by \mbox{\texttt{\mdseries\slshape options}}; the possible options \mbox{\texttt{\mdseries\slshape options}} and \emph{required} option (\texttt{"Prime"}) are as for \texttt{Pq} (see{\nobreakspace}\texttt{Pq} (\ref{Pq})). \texttt{PqEpimorphism} only differs from \texttt{Pq} in what it outputs; everything about what must/may be passed as input to \texttt{PqEpimorphism} is the same as for \texttt{Pq}. The same alternative methods of passing options to the
non\texttt{\symbol{45}}interactive \texttt{Pq} function are available to the non\texttt{\symbol{45}}interactive version of \texttt{PqEpimorphism}. 

 \emph{Notes:} \texttt{PqEpimorphism} may also be called with no arguments or one integer argument, in which case it
is being used interactively (see{\nobreakspace}\texttt{PqEpimorphism} (\ref{PqEpimorphism:interactive})), and the options \texttt{SetupFile} and \texttt{PqWorkspace} are ignored by the interactive \texttt{PqEpimorphism} function. 

 See Section{\nobreakspace}\hyperref[Attributes and a Property for fp and pc p-groups]{`Attributes and a Property for fp and pc p\texttt{\symbol{45}}groups'} for the attributes and property \texttt{NuclearRank}, \texttt{MultiplicatorRank} and \texttt{IsCapable} which may be applied to the image group of the epimorphism returned by \texttt{PqEpimorphism}. 
\begin{Verbatim}[commandchars=!@|,fontsize=\small,frame=single,label=Example]
  !gapprompt@gap>| !gapinput@F := FreeGroup (2, "F");|
  <free group on the generators [ F1, F2 ]>
  !gapprompt@gap>| !gapinput@phi := PqEpimorphism( F : Prime := 5, ClassBound := 2 );|
  [ F1, F2 ] -> [ f1, f2 ]
  !gapprompt@gap>| !gapinput@Image( phi );|
  <pc group of size 3125 with 5 generators>
\end{Verbatim}
 Typing: \texttt{PqExample( "PqEpimorphism" );} runs the above example in \textsf{GAP} (see{\nobreakspace}\texttt{PqExample} (\ref{PqExample})). }

 

\subsection{\textcolor{Chapter }{PqPCover}}
\logpage{[ 4, 1, 3 ]}\nobreak
\hyperdef{L}{X81C3CE1E850DA252}{}
{\noindent\textcolor{FuncColor}{$\triangleright$\enspace\texttt{PqPCover({\mdseries\slshape F: options})\index{PqPCover@\texttt{PqPCover}}
\label{PqPCover}
}\hfill{\scriptsize (function)}}\\


 returns for the fp or pc group \mbox{\texttt{\mdseries\slshape F}}, the $p$\texttt{\symbol{45}}covering group of the $p$\texttt{\symbol{45}}quotient of \mbox{\texttt{\mdseries\slshape F}} specified by \mbox{\texttt{\mdseries\slshape options}}, as a pc group, i.e.{\nobreakspace}the $p$\texttt{\symbol{45}}covering group of the $p$\texttt{\symbol{45}}quotient \texttt{Pq( \mbox{\texttt{\mdseries\slshape F}} : \mbox{\texttt{\mdseries\slshape options}} )}. Thus the options that \texttt{PqPCover} accepts are exactly those expected for \texttt{Pq} (and hence as a minimum the user \emph{must} supply a value for the \texttt{Prime} option; see{\nobreakspace}\texttt{Pq} (\ref{Pq}) for more details), except in the following special case. 

 If \mbox{\texttt{\mdseries\slshape F}} is already a $p$\texttt{\symbol{45}}group, in the sense that \texttt{IsPGroup(\mbox{\texttt{\mdseries\slshape F}})} is \texttt{true}, then 
\begin{description}
\item[{\texttt{Prime}}]  defaults to \texttt{PrimePGroup(\mbox{\texttt{\mdseries\slshape F}})}, if not supplied and \texttt{HasPrimePGroup(\mbox{\texttt{\mdseries\slshape F}}) = true}; and 
\item[{\texttt{ClassBound}}]  defaults to \texttt{PClassPGroup(\mbox{\texttt{\mdseries\slshape F}})} if \texttt{HasPClassPGroup(\mbox{\texttt{\mdseries\slshape F}}) = true} if not supplied, or to the usual default of 63, otherwise. 
\end{description}
 The same alternative methods of passing options to the
non\texttt{\symbol{45}}interactive \texttt{Pq} function are available to the non\texttt{\symbol{45}}interactive version of \texttt{PqPCover}. 

 We now give a few examples of the use of \texttt{PqPCover}. These examples are just a subset of the ones we gave for \texttt{Pq} (see{\nobreakspace}\texttt{Pq} (\ref{Pq})), except that in each instance the command \texttt{Pq} has been replaced with \texttt{PqPCover}. Essentially the same examples may be run by typing: \texttt{PqExample( "PqPCover" );} (see{\nobreakspace}\texttt{PqExample} (\ref{PqExample})). 
\begin{Verbatim}[commandchars=!@|,fontsize=\small,frame=single,label=Example]
  !gapprompt@gap>| !gapinput@F := FreeGroup("a", "b");; a := F.1;; b := F.2;;|
  !gapprompt@gap>| !gapinput@PqPCover( F : Prime := 2, ClassBound := 3 );|
  <pc group of size 262144 with 18 generators>
  !gapprompt@gap>| !gapinput@|
  !gapprompt@gap>| !gapinput@# Now let's get a p-cover of a p-quotient of an fp group|
  !gapprompt@gap>| !gapinput@G := F / [a^4, b^4];|
  <fp group on the generators [ a, b ]>
  !gapprompt@gap>| !gapinput@PqPCover( G : Prime := 2, ClassBound := 3 );|
  <pc group of size 16384 with 14 generators>
  !gapprompt@gap>| !gapinput@|
  !gapprompt@gap>| !gapinput@# Now let's get a p-cover of a different p-quotient of the same group|
  !gapprompt@gap>| !gapinput@PqPCover( G : Prime := 2, ClassBound := 3, Exponent := 4 );|
  <pc group of size 8192 with 13 generators>
  !gapprompt@gap>| !gapinput@|
  !gapprompt@gap>| !gapinput@# Now we'll get a p-cover of a p-quotient of another fp group|
  !gapprompt@gap>| !gapinput@# which we will redo using the `Relators' option|
  !gapprompt@gap>| !gapinput@R := [ a^25, Comm(Comm(b, a), a), b^5 ];|
  [ a^25, a^-1*b^-1*a*b*a^-1*b^-1*a^-1*b*a^2, b^5 ]
  !gapprompt@gap>| !gapinput@H := F / R;|
  <fp group on the generators [ a, b ]>
  !gapprompt@gap>| !gapinput@PqPCover( H : Prime := 5, ClassBound := 5, Metabelian );|
  <pc group of size 48828125 with 11 generators>
  !gapprompt@gap>| !gapinput@|
  !gapprompt@gap>| !gapinput@# Now we redo the previous example using the `Relators' option|
  !gapprompt@gap>| !gapinput@F := FreeGroup("a", "b");;|
  !gapprompt@gap>| !gapinput@rels := [ "a^25", "[b, a, a]", "b^5" ];|
  [ "a^25", "[b, a, a]", "b^5" ]
  !gapprompt@gap>| !gapinput@PqPCover( F : Prime := 5, ClassBound := 5, Metabelian, |
  !gapprompt@>| !gapinput@                 Relators := rels );|
  <pc group of size 48828125 with 11 generators>
\end{Verbatim}
 }

 }

 
\section{\textcolor{Chapter }{Computing Standard Presentations}}\label{Computing Standard Presentations}
\logpage{[ 4, 2, 0 ]}
\hyperdef{L}{X7F05F8DC79733A8E}{}
{
  \index{automorphisms!of $p$\texttt{\symbol{45}}groups} 

\subsection{\textcolor{Chapter }{PqStandardPresentation}}
\logpage{[ 4, 2, 1 ]}\nobreak
\hyperdef{L}{X86C9575D7CB65FBB}{}
{\noindent\textcolor{FuncColor}{$\triangleright$\enspace\texttt{PqStandardPresentation({\mdseries\slshape F: options})\index{PqStandardPresentation@\texttt{PqStandardPresentation}}
\label{PqStandardPresentation}
}\hfill{\scriptsize (function)}}\\
\noindent\textcolor{FuncColor}{$\triangleright$\enspace\texttt{StandardPresentation({\mdseries\slshape F: options})\index{StandardPresentation@\texttt{StandardPresentation}}
\label{StandardPresentation}
}\hfill{\scriptsize (method)}}\\


 return the \mbox{\texttt{\mdseries\slshape p}}\texttt{\symbol{45}}quotient specified by \mbox{\texttt{\mdseries\slshape options}} of the fp or pc $p$\texttt{\symbol{45}}group \mbox{\texttt{\mdseries\slshape F}}, as an \emph{fp group} which has a standard presentation. Here \mbox{\texttt{\mdseries\slshape options}} is a selection of the options from the following list (see
Chapter{\nobreakspace}\hyperref[ANUPQ Options]{`ANUPQ Options'} for detailed descriptions). Section{\nobreakspace}\hyperref[Hints and Warnings regarding the use of Options]{`Hints and Warnings regarding the use of Options'} gives some important hints and warnings regarding option usage, and
Section{\nobreakspace} \textbf{Reference: Function Call With Options} in the \textsf{GAP} Reference Manual describes their ``record''\texttt{\symbol{45}}like syntax. 
\begin{itemize}
\item  \texttt{Prime := \mbox{\texttt{\mdseries\slshape p}}}\index{option Prime} 
\item  \texttt{pQuotient := \mbox{\texttt{\mdseries\slshape Q}}}\index{option pQuotient} 
\item  \texttt{ClassBound := \mbox{\texttt{\mdseries\slshape n}}}\index{option ClassBound} 
\item  \texttt{Exponent := \mbox{\texttt{\mdseries\slshape n}}}\index{option Exponent} 
\item  \texttt{Metabelian}\index{option Metabelian} 
\item  \texttt{GroupName := \mbox{\texttt{\mdseries\slshape name}}}\index{option GroupName} 
\item  \texttt{OutputLevel := \mbox{\texttt{\mdseries\slshape n}}}\index{option OutputLevel} 
\item  \texttt{StandardPresentationFile := \mbox{\texttt{\mdseries\slshape filename}}}\index{option StandardPresentationFile} 
\item  \texttt{SetupFile := \mbox{\texttt{\mdseries\slshape filename}}}\index{option SetupFile} 
\item  \texttt{PqWorkspace := \mbox{\texttt{\mdseries\slshape workspace}}}\index{option PqWorkspace} 
\end{itemize}
 Unless \mbox{\texttt{\mdseries\slshape F}} is a pc \mbox{\texttt{\mdseries\slshape p}}\texttt{\symbol{45}}group, the user \emph{must} supply either the option \texttt{Prime} or the option \texttt{pQuotient} (if both \texttt{Prime} and \texttt{pQuotient} are supplied, the prime \mbox{\texttt{\mdseries\slshape p}} is determined by applying \texttt{PrimePGroup} (see{\nobreakspace}\texttt{PrimePGroup} (\textbf{Reference: PrimePGroup}) in the Reference Manual) to the value of \texttt{pQuotient}). 

 The options for \texttt{PqStandardPresentation} may also be passed in the two other alternative ways described for \texttt{Pq} (see{\nobreakspace}\texttt{Pq} (\ref{Pq})). \texttt{StandardPresentation} does not provide these alternative ways of passing options. 

 \emph{Notes:} In contrast to the function \texttt{Pq} (see{\nobreakspace}\texttt{Pq} (\ref{Pq})) which returns a pc group, \texttt{PqStandardPresentation} or \texttt{StandardPresentation} returns an fp group. This is because the output is mainly used for isomorphism
testing for which an fp group is enough. However, the presentation is a
polycyclic presentation and if you need to do any further computation with
this group (e.g.{\nobreakspace}to find the order) you can use the function \texttt{PcGroupFpGroup} (see{\nobreakspace}\texttt{PcGroupFpGroup} (\textbf{Reference: PcGroupFpGroup}) in the \textsf{GAP} Reference Manual) to form a pc group. 

 If the user does not supply a \mbox{\texttt{\mdseries\slshape p}}\texttt{\symbol{45}}quotient \mbox{\texttt{\mdseries\slshape Q}} via the \texttt{pQuotient} option and the prime \mbox{\texttt{\mdseries\slshape p}} is either supplied or \mbox{\texttt{\mdseries\slshape F}} is a pc \mbox{\texttt{\mdseries\slshape p}}\texttt{\symbol{45}}group, then a \mbox{\texttt{\mdseries\slshape p}}\texttt{\symbol{45}}quotient \mbox{\texttt{\mdseries\slshape Q}} is computed. If the user does supply a \mbox{\texttt{\mdseries\slshape p}}\texttt{\symbol{45}}quotient \mbox{\texttt{\mdseries\slshape Q}} via the \texttt{pQuotient} option, the package \textsf{AutPGrp} is called to compute the automorphism group of \mbox{\texttt{\mdseries\slshape Q}}; an error will occur that asks the user to install the package \textsf{AutPGrp} if the automorphism group cannot be computed. 

 The attributes and property \texttt{NuclearRank}, \texttt{MultiplicatorRank} and \texttt{IsCapable} are set for the group returned by \texttt{PqStandardPresentation} or \texttt{StandardPresentation} (see Section{\nobreakspace}\hyperref[Attributes and a Property for fp and pc p-groups]{`Attributes and a Property for fp and pc p\texttt{\symbol{45}}groups'}). 

 We illustrate the method with the following examples. 
\begin{Verbatim}[commandchars=!@|,fontsize=\small,frame=single,label=Example]
  !gapprompt@gap>| !gapinput@F := FreeGroup( "a", "b" );; a := F.1;; b := F.2;;|
  !gapprompt@gap>| !gapinput@G := F / [a^25, Comm(Comm(b, a), a), b^5];|
  <fp group on the generators [ a, b ]>
  !gapprompt@gap>| !gapinput@S := StandardPresentation( G : Prime := 5, ClassBound := 10 );|
  <fp group on the generators [ f1, f2, f3, f4, f5, f6, f7, f8, f9, f10, f11, 
    f12, f13, f14, f15, f16, f17, f18, f19, f20, f21, f22, f23, f24, f25, f26 ]>
  !gapprompt@gap>| !gapinput@IsPcGroup( S );|
  false
  !gapprompt@gap>| !gapinput@# if we need to compute with S we should convert it to a pc group|
  !gapprompt@gap>| !gapinput@Spc := PcGroupFpGroup( S );|
  <pc group of size 1490116119384765625 with 26 generators>
  !gapprompt@gap>| !gapinput@|
  !gapprompt@gap>| !gapinput@H := F / [ a^625, Comm(Comm(Comm(Comm(b, a), a), a), a)/Comm(b, a)^5,|
  !gapprompt@>| !gapinput@              Comm(Comm(b, a), b), b^625 ];;|
  !gapprompt@gap>| !gapinput@StandardPresentation( H : Prime := 5, ClassBound := 15, Metabelian );|
  <fp group on the generators [ f1, f2, f3, f4, f5, f6, f7, f8, f9, f10, f11, 
    f12, f13, f14, f15, f16, f17, f18, f19, f20 ]>
  !gapprompt@gap>| !gapinput@|
  !gapprompt@gap>| !gapinput@F4 := FreeGroup( "a", "b", "c", "d" );;|
  !gapprompt@gap>| !gapinput@a := F4.1;; b := F4.2;; c := F4.3;; d := F4.4;;|
  !gapprompt@gap>| !gapinput@G4 := F4 / [ b^4, b^2 / Comm(Comm (b, a), a), d^16,|
  !gapprompt@>| !gapinput@                a^16 / (c * d), b^8 / (d * c^4) ];|
  <fp group on the generators [ a, b, c, d ]>
  !gapprompt@gap>| !gapinput@K := Pq( G4 : Prime := 2, ClassBound := 1 );|
  <pc group of size 4 with 2 generators>
  !gapprompt@gap>| !gapinput@StandardPresentation( G4 : pQuotient := K, ClassBound := 14 );|
  <fp group with 53 generators>
\end{Verbatim}
 Typing: \texttt{PqExample( "StandardPresentation" );} runs the above example in \textsf{GAP} (see{\nobreakspace}\texttt{PqExample} (\ref{PqExample})). }

 

\subsection{\textcolor{Chapter }{EpimorphismPqStandardPresentation}}
\logpage{[ 4, 2, 2 ]}\nobreak
\hyperdef{L}{X828C06D083C0D089}{}
{\noindent\textcolor{FuncColor}{$\triangleright$\enspace\texttt{EpimorphismPqStandardPresentation({\mdseries\slshape F: options})\index{EpimorphismPqStandardPresentation@\texttt{EpimorphismPqStandardPresentation}}
\label{EpimorphismPqStandardPresentation}
}\hfill{\scriptsize (function)}}\\
\noindent\textcolor{FuncColor}{$\triangleright$\enspace\texttt{EpimorphismStandardPresentation({\mdseries\slshape F: options})\index{EpimorphismStandardPresentation@\texttt{EpimorphismStandardPresentation}}
\label{EpimorphismStandardPresentation}
}\hfill{\scriptsize (method)}}\\


 Each of the above functions accepts the same arguments and options as the
function \texttt{StandardPresentation} (see{\nobreakspace}\texttt{StandardPresentation} (\ref{StandardPresentation})) and returns an epimorphism from the fp or pc group \mbox{\texttt{\mdseries\slshape F}} onto the finitely presented group given by a standard presentation,
i.e.{\nobreakspace}if \mbox{\texttt{\mdseries\slshape S}} is the standard presentation computed for the $p$\texttt{\symbol{45}}quotient of \mbox{\texttt{\mdseries\slshape F}} by \texttt{StandardPresentation} then \texttt{EpimorphismStandardPresentation} returns the epimorphism from \mbox{\texttt{\mdseries\slshape F}} to the group with presentation \mbox{\texttt{\mdseries\slshape S}}. 

 \emph{Note:} The attributes and property \texttt{NuclearRank}, \texttt{MultiplicatorRank} and \texttt{IsCapable} are set for the image group of the epimorphism returned by \texttt{EpimorphismPqStandardPresentation} or \texttt{EpimorphismStandardPresentation} (see Section{\nobreakspace}\hyperref[Attributes and a Property for fp and pc p-groups]{`Attributes and a Property for fp and pc p\texttt{\symbol{45}}groups'}). 

 We illustrate the function with the following example. 
\begin{Verbatim}[commandchars=!@|,fontsize=\small,frame=single,label=Example]
  !gapprompt@gap>| !gapinput@F := FreeGroup(6, "F");|
  <free group on the generators [ F1, F2, F3, F4, F5, F6 ]>
  !gapprompt@gap>| !gapinput@# For printing GAP uses the symbols F1, ... for the generators of F|
  !gapprompt@gap>| !gapinput@x := F.1;; y := F.2;; z := F.3;; w := F.4;; a := F.5;; b := F.6;;|
  !gapprompt@gap>| !gapinput@R := [x^3 / w, y^3 / w * a^2 * b^2, w^3 / b,|
  !gapprompt@>| !gapinput@         Comm (y, x) / z, Comm (z, x), Comm (z, y) / a, z^3 ];;|
  !gapprompt@gap>| !gapinput@Q := F / R;|
  <fp group on the generators [ F1, F2, F3, F4, F5, F6 ]>
  !gapprompt@gap>| !gapinput@# For printing GAP also uses the symbols F1, ... for the generators of Q|
  !gapprompt@gap>| !gapinput@# (the same as used for F) ... but the gen'rs of Q and F are different:|
  !gapprompt@gap>| !gapinput@GeneratorsOfGroup(F) = GeneratorsOfGroup(Q);|
  false
  !gapprompt@gap>| !gapinput@G := Pq( Q : Prime := 3, ClassBound := 3 );|
  <pc group of size 729 with 6 generators>
  !gapprompt@gap>| !gapinput@phi := EpimorphismStandardPresentation( Q : Prime := 3,|
  !gapprompt@>| !gapinput@                                               ClassBound := 3 );|
  [ F1, F2, F3, F4, F5, F6 ] -> [ f1*f2^2*f3*f4^2*f5^2, f1*f2*f3*f5, f3^2, 
    f4*f6^2, f5, f6 ]
  !gapprompt@gap>| !gapinput@Source(phi); # This is the group Q (GAP uses F1, ... for gen'r symbols)|
  <fp group of size infinity on the generators [ F1, F2, F3, F4, F5, F6 ]>
  !gapprompt@gap>| !gapinput@Range(phi);  # This is the group G (GAP uses f1, ... for gen'r symbols)|
  <fp group on the generators [ f1, f2, f3, f4, f5, f6 ]>
  !gapprompt@gap>| !gapinput@AssignGeneratorVariables(G);|
  #I  Assigned the global variables [ f1, f2, f3, f4, f5, f6 ]
  !gapprompt@gap>| !gapinput@# Just to see that the images of [F1, ..., F6] do generate G|
  !gapprompt@gap>| !gapinput@Group([ f1*f2^2*f3, f1*f2*f3*f4*f5^2*f6^2, f3^2, f4, f5, f6 ]) = G;|
  true
  !gapprompt@gap>| !gapinput@Size( Image(phi) );|
  729
\end{Verbatim}
 Typing: \texttt{PqExample( "EpimorphismStandardPresentation" );} runs the above example in \textsf{GAP} (see{\nobreakspace}\texttt{PqExample} (\ref{PqExample})). Note that \texttt{AssignGeneratorVariables} (see{\nobreakspace}\texttt{AssignGeneratorVariables} (\textbf{Reference: AssignGeneratorVariables})) has only been available since \textsf{GAP}{\nobreakspace}4.3. }

 }

 
\section{\textcolor{Chapter }{Testing p\texttt{\symbol{45}}Groups for Isomorphism}}\label{Testing p-Groups for Isomorphism}
\logpage{[ 4, 3, 0 ]}
\hyperdef{L}{X8030FAE586423615}{}
{
  

\subsection{\textcolor{Chapter }{IsPqIsomorphicPGroup}}
\logpage{[ 4, 3, 1 ]}\nobreak
\hyperdef{L}{X854389FC780BB178}{}
{\noindent\textcolor{FuncColor}{$\triangleright$\enspace\texttt{IsPqIsomorphicPGroup({\mdseries\slshape G, H})\index{IsPqIsomorphicPGroup@\texttt{IsPqIsomorphicPGroup}}
\label{IsPqIsomorphicPGroup}
}\hfill{\scriptsize (function)}}\\
\noindent\textcolor{FuncColor}{$\triangleright$\enspace\texttt{IsIsomorphicPGroup({\mdseries\slshape G, H})\index{IsIsomorphicPGroup@\texttt{IsIsomorphicPGroup}}
\label{IsIsomorphicPGroup}
}\hfill{\scriptsize (method)}}\\


 each return true if \mbox{\texttt{\mdseries\slshape G}} is isomorphic to \mbox{\texttt{\mdseries\slshape H}}, where both \mbox{\texttt{\mdseries\slshape G}} and \mbox{\texttt{\mdseries\slshape H}} must be pc groups of prime power order. These functions compute and compare in \textsf{GAP} the fp groups given by standard presentations for \mbox{\texttt{\mdseries\slshape G}} and \mbox{\texttt{\mdseries\slshape H}} (see \texttt{StandardPresentation} (\ref{StandardPresentation})). 
\begin{Verbatim}[commandchars=!@|,fontsize=\small,frame=single,label=Example]
  !gapprompt@gap>| !gapinput@G := Group( (1,2,3,4), (1,3) );|
  Group([ (1,2,3,4), (1,3) ])
  !gapprompt@gap>| !gapinput@P1 := Image( IsomorphismPcGroup( G ) );|
  Group([ f1, f2, f3 ])
  !gapprompt@gap>| !gapinput@P2 := ElementaryAbelianGroup( 8 );|
  <pc group of size 8 with 3 generators>
  !gapprompt@gap>| !gapinput@IsIsomorphicPGroup( P1, P2 );|
  false
  !gapprompt@gap>| !gapinput@P3 := QuaternionGroup( 8 );|
  <pc group of size 8 with 3 generators>
  !gapprompt@gap>| !gapinput@IsIsomorphicPGroup( P1, P3 );|
  false
  !gapprompt@gap>| !gapinput@P4 := DihedralGroup( 8 );|
  <pc group of size 8 with 3 generators>
  !gapprompt@gap>| !gapinput@IsIsomorphicPGroup( P1, P4 );|
  true
\end{Verbatim}
 Typing: \texttt{PqExample( "IsIsomorphicPGroup" );} runs the above example in \textsf{GAP} (see{\nobreakspace}\texttt{PqExample} (\ref{PqExample})). }

 }

 
\section{\textcolor{Chapter }{Computing Descendants of a p\texttt{\symbol{45}}Group}}\label{Computing Descendants of a p-Group}
\logpage{[ 4, 4, 0 ]}
\hyperdef{L}{X8450C72580D91245}{}
{
  

\subsection{\textcolor{Chapter }{PqDescendants}}
\logpage{[ 4, 4, 1 ]}\nobreak
\hyperdef{L}{X80985CC479CD9FA3}{}
{\noindent\textcolor{FuncColor}{$\triangleright$\enspace\texttt{PqDescendants({\mdseries\slshape G: options})\index{PqDescendants@\texttt{PqDescendants}}
\label{PqDescendants}
}\hfill{\scriptsize (function)}}\\


 returns, for the pc group \mbox{\texttt{\mdseries\slshape G}} which must be of prime power order with a confluent pc presentation
(see{\nobreakspace}\texttt{IsConfluent} (\textbf{Reference: IsConfluent for pc groups}) in the \textsf{GAP} Reference Manual), a list of proper descendants (pc groups) of \mbox{\texttt{\mdseries\slshape G}}. Following the colon \mbox{\texttt{\mdseries\slshape options}} a selection of the options listed below should be given, separated by commas
like record components (see Section{\nobreakspace} \textbf{Reference: Function Call With Options} in the \textsf{GAP} Reference Manual). See Chapter{\nobreakspace}\hyperref[ANUPQ Options]{`ANUPQ Options'} for detailed descriptions of the options. 

 The automorphism group of each descendant \mbox{\texttt{\mdseries\slshape D}} is also computed via a call to the \texttt{AutomorphismGroupPGroup} function of the \textsf{AutPGrp} package.      
\begin{itemize}
\item  \texttt{ClassBound := \mbox{\texttt{\mdseries\slshape n}}}\index{option ClassBound} 
\item  \texttt{Relators := \mbox{\texttt{\mdseries\slshape rels}}}\index{option Relators} 
\item  \texttt{OrderBound := \mbox{\texttt{\mdseries\slshape n}}}\index{option OrderBound} 
\item  \texttt{StepSize := \mbox{\texttt{\mdseries\slshape n}}}, \texttt{StepSize := \mbox{\texttt{\mdseries\slshape list}}} \index{option StepSize} 
\item  \texttt{RankInitialSegmentSubgroups := \mbox{\texttt{\mdseries\slshape n}}}\index{option RankInitialSegmentSubgroups} 
\item  \texttt{SpaceEfficient}\index{option SpaceEfficient} 
\item  \texttt{CapableDescendants}\index{option CapableDescendants} 
\item  \texttt{AllDescendants := false}\index{option AllDescendants} 
\item  \texttt{Exponent := \mbox{\texttt{\mdseries\slshape n}}}\index{option Exponent} 
\item  \texttt{Metabelian}\index{option Metabelian} 
\item  \texttt{GroupName := \mbox{\texttt{\mdseries\slshape name}}}\index{option GroupName} 
\item  \texttt{SubList := \mbox{\texttt{\mdseries\slshape sub}}}\index{option SubList} 
\item  \texttt{BasicAlgorithm}\index{option BasicAlgorithm} 
\item  \texttt{CustomiseOutput := \mbox{\texttt{\mdseries\slshape rec}}}\index{option CustomiseOutput} 
\item  \texttt{SetupFile := \mbox{\texttt{\mdseries\slshape filename}}}\index{option SetupFile} 
\item  \texttt{PqWorkspace := \mbox{\texttt{\mdseries\slshape workspace}}}\index{option PqWorkspace} 
\end{itemize}
 \emph{Notes:} The function \texttt{PqDescendants} uses the automorphism group of \mbox{\texttt{\mdseries\slshape G}} which it computes via the package \textsf{AutPGrp}. If this package is not installed an error may be raised. If the automorphism
group of \mbox{\texttt{\mdseries\slshape G}} is insoluble, the \texttt{pq} program will call \textsf{GAP} together with the \textsf{AutPGrp} package for certain orbit\texttt{\symbol{45}}stabilizer calculations. (So, in
any case, one should ensure the \textsf{AutPGrp} package is installed.) 

 The attributes and property \texttt{NuclearRank}, \texttt{MultiplicatorRank} and \texttt{IsCapable} are set for each group of the list returned by \texttt{PqDescendants} (see Section{\nobreakspace}\hyperref[Attributes and a Property for fp and pc p-groups]{`Attributes and a Property for fp and pc p\texttt{\symbol{45}}groups'}). 

 The options \mbox{\texttt{\mdseries\slshape options}} for \texttt{PqDescendants} may be passed in an alternative manner to that already described, namely you
can pass \texttt{PqDescendants} a record as an argument, which contains as entries some (or all) of the above
mentioned. Those parameters which do not occur in the record are set to their
default values. 

 Note that you cannot set both \texttt{OrderBound} and \texttt{StepSize}. 

 In the first example we compute all proper descendants of the Klein four group
which have exponent\texttt{\symbol{45}}2 class at most 5 and order at most $2^6$. 
\begin{Verbatim}[commandchars=!@|,fontsize=\small,frame=single,label=Example]
  !gapprompt@gap>| !gapinput@F := FreeGroup( "a", "b" );; a := F.1;; b := F.2;;|
  !gapprompt@gap>| !gapinput@G := PcGroupFpGroup( F / [ a^2, b^2, Comm(b, a) ] );|
  <pc group of size 4 with 2 generators>
  !gapprompt@gap>| !gapinput@des := PqDescendants( G : OrderBound := 6, ClassBound := 5 );;|
  !gapprompt@gap>| !gapinput@Length(des);|
  83
  !gapprompt@gap>| !gapinput@List(des, Size); |
  [ 8, 8, 8, 16, 16, 16, 32, 16, 16, 16, 16, 16, 32, 32, 64, 64, 32, 32, 32, 
    32, 32, 32, 32, 64, 64, 64, 64, 64, 64, 64, 64, 64, 64, 64, 32, 32, 32, 32, 
    64, 64, 64, 64, 64, 64, 64, 64, 64, 64, 64, 32, 32, 32, 32, 32, 64, 64, 64, 
    64, 64, 64, 64, 64, 64, 64, 64, 64, 64, 64, 64, 64, 64, 64, 64, 64, 64, 64, 
    64, 64, 64, 64, 64, 64, 64 ]
  !gapprompt@gap>| !gapinput@List(des, d -> Length( PCentralSeries( d, 2 ) ) - 1 );|
  [ 2, 2, 2, 2, 2, 2, 2, 3, 3, 3, 3, 3, 3, 3, 3, 3, 3, 3, 3, 3, 3, 3, 3, 3, 3, 
    3, 3, 3, 3, 3, 3, 3, 3, 3, 3, 3, 3, 3, 3, 3, 3, 3, 3, 3, 3, 3, 3, 3, 3, 4, 
    4, 4, 4, 4, 4, 4, 4, 4, 4, 4, 4, 4, 4, 4, 4, 4, 4, 4, 4, 4, 4, 4, 4, 4, 4, 
    4, 4, 4, 5, 5, 5, 5, 5 ]
\end{Verbatim}
 Below, we compute all capable descendants of order 27 of the elementary
abelian group of order 9. 
\begin{Verbatim}[commandchars=!@|,fontsize=\small,frame=single,label=Example]
  !gapprompt@gap>| !gapinput@F := FreeGroup( 2, "g" );|
  <free group on the generators [ g1, g2 ]>
  !gapprompt@gap>| !gapinput@G := PcGroupFpGroup( F / [ F.1^3, F.2^3, Comm(F.1, F.2) ] );|
  <pc group of size 9 with 2 generators>
  !gapprompt@gap>| !gapinput@des := PqDescendants( G : OrderBound := 3, ClassBound := 2,|
  !gapprompt@>| !gapinput@                             CapableDescendants );|
  [ <pc group of size 27 with 3 generators>, 
    <pc group of size 27 with 3 generators> ]
  !gapprompt@gap>| !gapinput@List(des, d -> Length( PCentralSeries( d, 3 ) ) - 1 );|
  [ 2, 2 ]
  !gapprompt@gap>| !gapinput@# For comparison let us now compute all proper descendants|
  !gapprompt@gap>| !gapinput@PqDescendants( G : OrderBound := 3, ClassBound := 2);|
  [ <pc group of size 27 with 3 generators>, 
    <pc group of size 27 with 3 generators>, 
    <pc group of size 27 with 3 generators> ]
\end{Verbatim}
 In the third example, we compute all proper capable descendants of the
elementary abelian group of order $5^2$ which have exponent\texttt{\symbol{45}}$5$ class at most $3$, exponent $5$, and are metabelian. 
\begin{Verbatim}[commandchars=!@|,fontsize=\small,frame=single,label=Example]
  !gapprompt@gap>| !gapinput@F := FreeGroup( 2, "g" );;|
  !gapprompt@gap>| !gapinput@G := PcGroupFpGroup( F / [ F.1^5, F.2^5, Comm(F.2, F.1) ] );|
  <pc group of size 25 with 2 generators>
  !gapprompt@gap>| !gapinput@des := PqDescendants( G : Metabelian, ClassBound := 3,|
  !gapprompt@>| !gapinput@                             Exponent := 5, CapableDescendants );|
  [ <pc group of size 125 with 3 generators>, 
    <pc group of size 625 with 4 generators>, 
    <pc group of size 3125 with 5 generators> ]
  !gapprompt@gap>| !gapinput@List(des, d -> Length( PCentralSeries( d, 5 ) ) - 1 );|
  [ 2, 3, 3 ]
  !gapprompt@gap>| !gapinput@List(des, d -> Length( DerivedSeries( d ) ) );|
  [ 3, 3, 3 ]
  !gapprompt@gap>| !gapinput@List(des, d -> Maximum( List( d, Order ) ) );|
  [ 5, 5, 5 ]
\end{Verbatim}
 The examples \texttt{"PqDescendants\texttt{\symbol{45}}1"}, \texttt{"PqDescendants\texttt{\symbol{45}}2"} and \texttt{"PqDescendants\texttt{\symbol{45}}3"} (in order) are essentially the same as the above three examples
(see{\nobreakspace}\texttt{PqExample} (\ref{PqExample})). }

 

\subsection{\textcolor{Chapter }{PqSupplementInnerAutomorphisms}}
\logpage{[ 4, 4, 2 ]}\nobreak
\hyperdef{L}{X8364566A80A8C24B}{}
{\noindent\textcolor{FuncColor}{$\triangleright$\enspace\texttt{PqSupplementInnerAutomorphisms({\mdseries\slshape D})\index{PqSupplementInnerAutomorphisms@\texttt{PqSupplementInnerAutomorphisms}}
\label{PqSupplementInnerAutomorphisms}
}\hfill{\scriptsize (function)}}\\


 returns a generating set for a supplement to the inner automorphisms of \mbox{\texttt{\mdseries\slshape D}}, in the form of a record with fields \texttt{agAutos}, \texttt{agOrder} and \texttt{glAutos}, as provided by the \texttt{pq} program. One should be very careful in using these automorphisms for a
descendant calculation.       

 \emph{Note:} In principle there must be a way to use those automorphisms in order to
compute descendants but there does not seem to be a way to hand back these
automorphisms properly to the \texttt{pq} program. 
\begin{Verbatim}[commandchars=!@|,fontsize=\small,frame=single,label=Example]
  !gapprompt@gap>| !gapinput@Q := Pq( FreeGroup(2) : Prime := 3, ClassBound := 1 );|
  <pc group of size 9 with 2 generators>
  !gapprompt@gap>| !gapinput@des := PqDescendants( Q : StepSize := 1 );|
  [ <pc group of size 27 with 3 generators>, 
    <pc group of size 27 with 3 generators>, 
    <pc group of size 27 with 3 generators> ]
  !gapprompt@gap>| !gapinput@S := PqSupplementInnerAutomorphisms( des[3] );|
  rec( agAutos := [  ], agOrder := [ 3, 2, 2, 2 ], 
    glAutos := [ Pcgs([ f1, f2, f3 ]) -> [ f1*f2^2, f2, f3 ], 
        Pcgs([ f1, f2, f3 ]) -> [ f1^2, f2, f3^2 ], 
        Pcgs([ f1, f2, f3 ]) -> [ f1^2, f2, f3^2 ] ] )
  !gapprompt@gap>| !gapinput@A := AutomorphismGroupPGroup( des[3] );|
  rec( 
    agAutos := [ Pcgs([ f1, f2, f3 ]) -> [ f1^2, f2, f3^2 ], 
        Pcgs([ f1, f2, f3 ]) -> [ f1*f2^2, f2, f3 ], 
        Pcgs([ f1, f2, f3 ]) -> [ f1*f3, f2, f3 ], 
        Pcgs([ f1, f2, f3 ]) -> [ f1, f2*f3, f3 ] ], agOrder := [ 2, 3, 3, 3 ], 
    glAutos := [  ], glOper := [  ], glOrder := 1, 
    group := <pc group of size 27 with 3 generators>, 
    one := IdentityMapping( <pc group of size 27 with 3 generators> ), 
    size := 54 )
\end{Verbatim}
 Typing: \texttt{PqExample( "PqSupplementInnerAutomorphisms" );} runs the above example in \textsf{GAP} (see{\nobreakspace}\texttt{PqExample} (\ref{PqExample})). 

 Note that by also including \texttt{PqStart} as a second argument to \texttt{PqExample} one can see how it is possible, with the aid of \texttt{PqSetPQuotientToGroup} (see{\nobreakspace}\texttt{PqSetPQuotientToGroup} (\ref{PqSetPQuotientToGroup})), to do the equivalent computations with the interactive versions of \texttt{Pq} and \texttt{PqDescendants} and a single \texttt{pq} process (recall \texttt{pq} is the name of the external C program). }

 

\subsection{\textcolor{Chapter }{PqList}}
\logpage{[ 4, 4, 3 ]}\nobreak
\hyperdef{L}{X79AD54E987FCACBA}{}
{\noindent\textcolor{FuncColor}{$\triangleright$\enspace\texttt{PqList({\mdseries\slshape filename: [SubList := sub]})\index{PqList@\texttt{PqList}}
\label{PqList}
}\hfill{\scriptsize (function)}}\\


 reads a file with name \mbox{\texttt{\mdseries\slshape filename}} (a string) and returns the list \mbox{\texttt{\mdseries\slshape L}} of pc groups (or with option \texttt{SubList} a sublist of \mbox{\texttt{\mdseries\slshape L}} or a single pc group in \mbox{\texttt{\mdseries\slshape L}}) defined in that file. If the option \texttt{SubList} is passed and has the value \mbox{\texttt{\mdseries\slshape sub}}, then it has the same meaning as for \texttt{PqDescendants}, i.e.{\nobreakspace}if \mbox{\texttt{\mdseries\slshape sub}} is an integer then \texttt{PqList} returns \texttt{\mbox{\texttt{\mdseries\slshape L}}[\mbox{\texttt{\mdseries\slshape sub}}]}; otherwise, if \mbox{\texttt{\mdseries\slshape sub}} is a list of integers \texttt{PqList} returns \texttt{Sublist(\mbox{\texttt{\mdseries\slshape L}}, \mbox{\texttt{\mdseries\slshape sub}} )}. 

 Both \texttt{PqList} and \texttt{SavePqList} (see \texttt{SavePqList} (\ref{SavePqList})) can be used to save and restore a list of descendants (see \texttt{PqDescendants} (\ref{PqDescendants})). }

 

\subsection{\textcolor{Chapter }{SavePqList}}
\logpage{[ 4, 4, 4 ]}\nobreak
\hyperdef{L}{X85ADB9E6870EC3D1}{}
{\noindent\textcolor{FuncColor}{$\triangleright$\enspace\texttt{SavePqList({\mdseries\slshape filename, list})\index{SavePqList@\texttt{SavePqList}}
\label{SavePqList}
}\hfill{\scriptsize (function)}}\\


 writes a list of descendants \mbox{\texttt{\mdseries\slshape list}} to a file with name \mbox{\texttt{\mdseries\slshape filename}} (a string). 

 \texttt{SavePqList} and \texttt{PqList} (see \texttt{PqList} (\ref{PqList})) can be used to save and restore, respectively, the results of \texttt{PqDescendants} (see \texttt{PqDescendants} (\ref{PqDescendants})). }

 }

 }

 
\chapter{\textcolor{Chapter }{Interactive ANUPQ functions}}\label{Interactive ANUPQ functions}
\logpage{[ 5, 0, 0 ]}
\hyperdef{L}{X842C13C27B941744}{}
{
  Here we describe the interactive functions defined by the \textsf{ANUPQ} package, i.e.{\nobreakspace}the functions that manipulate and initiate
interactive \textsf{ANUPQ} processes. These are functions that extract information via a dialogue with a
running \texttt{pq} process (process used in the UNIX sense). Occasionally, a user needs the ``next step''; the functions provided in this chapter make use of data from previous steps
retained by the \texttt{pq} program, thus allowing the user to interact with the \texttt{pq} program like one can when one uses the \texttt{pq} program as a stand\texttt{\symbol{45}}alone (see{\nobreakspace}\texttt{guide.dvi} in the \texttt{standalone\texttt{\symbol{45}}doc} directory). 

 An interactive \textsf{ANUPQ} process is initiated by \texttt{PqStart} and terminated via \texttt{PqQuit}; these functions are described in ection{\nobreakspace}\hyperref[Starting and Stopping Interactive ANUPQ Processes]{`Starting and Stopping Interactive ANUPQ Processes'}. 

 Each interactive \textsf{ANUPQ} function that manipulates an already started interactive \textsf{ANUPQ} process, has a form where the first argument is the integer \mbox{\texttt{\mdseries\slshape i}} returned by the initiating \texttt{PqStart} command, and a second form with one argument fewer (where the integer \mbox{\texttt{\mdseries\slshape i}} is discovered by a default mechanism, namely by determining the least integer \mbox{\texttt{\mdseries\slshape i}} for which there is a currently active interactive \textsf{ANUPQ} process). We will thus commonly say that ``for the \mbox{\texttt{\mdseries\slshape i}}th (or default) interactive \textsf{ANUPQ} process'' a certain function performs a given action. In each case, it is an error if \mbox{\texttt{\mdseries\slshape i}} is not the index of an active interactive process, or there are no current
active interactive processes. 

 \emph{Notes}: The global method of passing options (via \texttt{PushOptions}), should not be used with any of the interactive functions. In fact, the \texttt{OptionsStack} should be empty at the time any of the interactive functions is called. 

 On \texttt{quit}ting \textsf{GAP}, \texttt{PqQuitAll();} is executed, which terminates all active interactive \textsf{ANUPQ} processes. If \textsf{GAP} is killed without \texttt{quit}ting, before all interactive \textsf{ANUPQ} processes are terminated, \emph{zombie} processes (still living \emph{child} processes whose \emph{parents} have died), may result. Since zombie processes do consume resources, in such
an event, the responsible computer user should seek out and terminate those
zombie processes (e.g.{\nobreakspace}on Linux: \texttt{ps xw | grep pq} gives you information on the \texttt{pq} processes corresponding to any interactive \textsf{ANUPQ} processes started in a \textsf{GAP} session; you can then do \texttt{kill \mbox{\texttt{\mdseries\slshape N}}} for each number \mbox{\texttt{\mdseries\slshape N}} appearing in the first column of this output). 
\section{\textcolor{Chapter }{Starting and Stopping Interactive ANUPQ Processes}}\label{Starting and Stopping Interactive ANUPQ Processes}
\logpage{[ 5, 1, 0 ]}
\hyperdef{L}{X7935CDAD7936CA0A}{}
{
  

\subsection{\textcolor{Chapter }{PqStart (with group and workspace size)}}
\logpage{[ 5, 1, 1 ]}\nobreak
\hyperdef{L}{X83B2EC237F37623C}{}
{\noindent\textcolor{FuncColor}{$\triangleright$\enspace\texttt{PqStart({\mdseries\slshape G, workspace: options})\index{PqStart@\texttt{PqStart}!with group and workspace size}
\label{PqStart:with group and workspace size}
}\hfill{\scriptsize (function)}}\\
\noindent\textcolor{FuncColor}{$\triangleright$\enspace\texttt{PqStart({\mdseries\slshape G: options})\index{PqStart@\texttt{PqStart}!with group}
\label{PqStart:with group}
}\hfill{\scriptsize (function)}}\\
\noindent\textcolor{FuncColor}{$\triangleright$\enspace\texttt{PqStart({\mdseries\slshape workspace: options})\index{PqStart@\texttt{PqStart}!with workspace size}
\label{PqStart:with workspace size}
}\hfill{\scriptsize (function)}}\\
\noindent\textcolor{FuncColor}{$\triangleright$\enspace\texttt{PqStart({\mdseries\slshape : options})\index{PqStart@\texttt{PqStart}}
\label{PqStart}
}\hfill{\scriptsize (function)}}\\


 activate an iostream for an interactive \textsf{ANUPQ} process (i.e. \texttt{PqStart} starts up a \texttt{pq} process and opens a \textsf{GAP} iostream to ``talk'' to that process) and returns an integer \mbox{\texttt{\mdseries\slshape i}} that can be used to identify that process. The argument \mbox{\texttt{\mdseries\slshape G}} should be an \emph{fp group} or \emph{pc group} that the user intends to manipulate using interactive \textsf{ANUPQ} functions. If the function is called without specifying \mbox{\texttt{\mdseries\slshape G}}, a group can be read in by using the function \texttt{PqRestorePcPresentation} (see{\nobreakspace}\texttt{PqRestorePcPresentation} (\ref{PqRestorePcPresentation})). If \texttt{PqStart} is given an integer argument \mbox{\texttt{\mdseries\slshape workspace}}, then the \texttt{pq} program is started up with a workspace (an integer array) of size \mbox{\texttt{\mdseries\slshape workspace}} (i.e. $4 \times \mbox{\texttt{\mdseries\slshape workspace}}$ bytes in a 32\texttt{\symbol{45}}bit environment); otherwise, the \texttt{pq} program sets a default workspace of $10000000$. 

 The only \mbox{\texttt{\mdseries\slshape options}} currently recognised by \texttt{PqStart} are \texttt{Prime}, \texttt{Exponent} and \texttt{Relators} (see Chapter{\nobreakspace}\hyperref[ANUPQ Options]{`ANUPQ Options'} for detailed descriptions of these options) and if provided they are
essentially global for the interactive \textsf{ANUPQ} process, except that any interactive function interacting with the process and
passing new values for these options will over\texttt{\symbol{45}}ride the
global values. }

 

\subsection{\textcolor{Chapter }{PqQuit}}
\logpage{[ 5, 1, 2 ]}\nobreak
\hyperdef{L}{X79DB761185BAB9C8}{}
{\noindent\textcolor{FuncColor}{$\triangleright$\enspace\texttt{PqQuit({\mdseries\slshape i})\index{PqQuit@\texttt{PqQuit}}
\label{PqQuit}
}\hfill{\scriptsize (function)}}\\
\noindent\textcolor{FuncColor}{$\triangleright$\enspace\texttt{PqQuit({\mdseries\slshape })\index{PqQuit@\texttt{PqQuit}!for default process}
\label{PqQuit:for default process}
}\hfill{\scriptsize (function)}}\\


 closes the stream of the \mbox{\texttt{\mdseries\slshape i}}th or default interactive \textsf{ANUPQ} process and unbinds its \texttt{ANUPQData.io} record. 

 \emph{Note:} It can happen that the \texttt{pq} process, and hence the \textsf{GAP} iostream assigned to communicate with it, can die, e.g.{\nobreakspace}by the
user typing a \texttt{Ctrl\texttt{\symbol{45}}C} while the \texttt{pq} process is engaged in a long calculation. \texttt{IsPqProcessAlive} (see{\nobreakspace}\texttt{IsPqProcessAlive} (\ref{IsPqProcessAlive})) is provided to check the status of the \textsf{GAP} iostream (and hence the status of the \texttt{pq} process it was communicating with). }

 

\subsection{\textcolor{Chapter }{PqQuitAll}}
\logpage{[ 5, 1, 3 ]}\nobreak
\hyperdef{L}{X7FF8F2657B1B008E}{}
{\noindent\textcolor{FuncColor}{$\triangleright$\enspace\texttt{PqQuitAll({\mdseries\slshape })\index{PqQuitAll@\texttt{PqQuitAll}}
\label{PqQuitAll}
}\hfill{\scriptsize (function)}}\\


 is provided as a convenience, to terminate all active interactive \textsf{ANUPQ} processes with a single command. It is equivalent to executing \texttt{PqQuit(\mbox{\texttt{\mdseries\slshape i}})} for all active interactive \textsf{ANUPQ} processes \mbox{\texttt{\mdseries\slshape i}} (see{\nobreakspace}\texttt{PqQuit} (\ref{PqQuit})). }

 }

 
\section{\textcolor{Chapter }{Interactive ANUPQ Process Utility Functions}}\label{Interactive ANUPQ Process Utility Functions}
\logpage{[ 5, 2, 0 ]}
\hyperdef{L}{X7EE8000A800CEE2D}{}
{
  

\subsection{\textcolor{Chapter }{PqProcessIndex}}
\logpage{[ 5, 2, 1 ]}\nobreak
\hyperdef{L}{X8558B20A80B999AE}{}
{\noindent\textcolor{FuncColor}{$\triangleright$\enspace\texttt{PqProcessIndex({\mdseries\slshape i})\index{PqProcessIndex@\texttt{PqProcessIndex}}
\label{PqProcessIndex}
}\hfill{\scriptsize (function)}}\\
\noindent\textcolor{FuncColor}{$\triangleright$\enspace\texttt{PqProcessIndex({\mdseries\slshape })\index{PqProcessIndex@\texttt{PqProcessIndex}!for default process}
\label{PqProcessIndex:for default process}
}\hfill{\scriptsize (function)}}\\


 With argument \mbox{\texttt{\mdseries\slshape i}}, which must be a positive integer, \texttt{PqProcessIndex} returns \mbox{\texttt{\mdseries\slshape i}} if it corresponds to an active interactive process, or raises an error. With
no arguments it returns the default active interactive process or returns \texttt{fail} and emits a warning message to \texttt{Info} at \texttt{InfoANUPQ} or \texttt{InfoWarning} level 1. 

 \emph{Note:} Essentially, an interactive \textsf{ANUPQ} process \mbox{\texttt{\mdseries\slshape i}} is ``active'' if \texttt{ANUPQData.io[\mbox{\texttt{\mdseries\slshape i}}]} is bound (i.e.{\nobreakspace}we still have some data telling us about it).
Also see{\nobreakspace}\texttt{PqStart} (\ref{PqStart}). }

 

\subsection{\textcolor{Chapter }{PqProcessIndices}}
\logpage{[ 5, 2, 2 ]}\nobreak
\hyperdef{L}{X7E4A56D67C865291}{}
{\noindent\textcolor{FuncColor}{$\triangleright$\enspace\texttt{PqProcessIndices({\mdseries\slshape })\index{PqProcessIndices@\texttt{PqProcessIndices}}
\label{PqProcessIndices}
}\hfill{\scriptsize (function)}}\\


 returns the list of integer indices of all active interactive \textsf{ANUPQ} processes (see{\nobreakspace}\texttt{PqProcessIndex} (\ref{PqProcessIndex}) for the meaning of ``active''). }

 

\subsection{\textcolor{Chapter }{IsPqProcessAlive}}
\logpage{[ 5, 2, 3 ]}\nobreak
\hyperdef{L}{X7C103DF78435AEC7}{}
{\noindent\textcolor{FuncColor}{$\triangleright$\enspace\texttt{IsPqProcessAlive({\mdseries\slshape i})\index{IsPqProcessAlive@\texttt{IsPqProcessAlive}}
\label{IsPqProcessAlive}
}\hfill{\scriptsize (function)}}\\
\noindent\textcolor{FuncColor}{$\triangleright$\enspace\texttt{IsPqProcessAlive({\mdseries\slshape })\index{IsPqProcessAlive@\texttt{IsPqProcessAlive}!for default process}
\label{IsPqProcessAlive:for default process}
}\hfill{\scriptsize (function)}}\\


 return \texttt{true} if the \textsf{GAP} iostream of the \mbox{\texttt{\mdseries\slshape i}}th (or default) interactive \textsf{ANUPQ} process started by \texttt{PqStart} is alive (i.e.{\nobreakspace}can still be written to), or \texttt{false}, otherwise. (See the notes for{\nobreakspace}\texttt{PqStart} (\ref{PqStart}) and{\nobreakspace}\texttt{PqQuit} (\ref{PqQuit}).) 

 \index{interruption} If the user does not yet have a \texttt{gap{\textgreater}} prompt then usually the \texttt{pq} program is still away doing something and an \textsf{ANUPQ} interface function is still waiting for a reply. Typing a \texttt{Ctrl\texttt{\symbol{45}}C} (i.e.{\nobreakspace}holding down the \texttt{Ctrl} key and typing \texttt{c}) will stop the waiting and send \textsf{GAP} into a \texttt{break}\texttt{\symbol{45}}loop, from which one has no option but to \texttt{quit;}. The typing of \texttt{Ctrl\texttt{\symbol{45}}C}, in such a circumstance, usually causes the stream of the interactive \textsf{ANUPQ} process to die; to check this we provide \texttt{IsPqProcessAlive} (see{\nobreakspace}\texttt{IsPqProcessAlive}). 

 The \textsf{GAP} iostream of an interactive \textsf{ANUPQ} process will also die if the \texttt{pq} program has a segmentation fault. We do hope that this never happens to you,
but if it does and the failure is reproducible, then it's a bug and we'd like
to know about it. Please read the \texttt{README} that comes with the \textsf{ANUPQ} package to find out what to include in a bug report and who to email it to. }

 }

 
\section{\textcolor{Chapter }{Interactive Versions of Non\texttt{\symbol{45}}interactive ANUPQ Functions}}\label{Interactive Versions of Non-interactive ANUPQ Functions}
\logpage{[ 5, 3, 0 ]}
\hyperdef{L}{X823D2BBF7C4515D4}{}
{
  

\subsection{\textcolor{Chapter }{Pq (interactive)}}
\logpage{[ 5, 3, 1 ]}\nobreak
\hyperdef{L}{X85408C4C790439E7}{}
{\noindent\textcolor{FuncColor}{$\triangleright$\enspace\texttt{Pq({\mdseries\slshape i: options})\index{Pq@\texttt{Pq}!interactive}
\label{Pq:interactive}
}\hfill{\scriptsize (function)}}\\
\noindent\textcolor{FuncColor}{$\triangleright$\enspace\texttt{Pq({\mdseries\slshape : options})\index{Pq@\texttt{Pq}!interactive, for default process}
\label{Pq:interactive, for default process}
}\hfill{\scriptsize (function)}}\\


 return, for the fp or pc group (let us call it \mbox{\texttt{\mdseries\slshape F}}), of the \mbox{\texttt{\mdseries\slshape i}}th or default interactive \textsf{ANUPQ} process, the $p$\texttt{\symbol{45}}quotient of \mbox{\texttt{\mdseries\slshape F}} specified by \mbox{\texttt{\mdseries\slshape options}}, as a pc group; \mbox{\texttt{\mdseries\slshape F}} must previously have been given (as first argument) to \texttt{PqStart} to start the interactive \textsf{ANUPQ} process (see{\nobreakspace}\texttt{PqStart} (\ref{PqStart})) or restored from file using the function \texttt{PqRestorePcPresentation} (see{\nobreakspace}\texttt{PqRestorePcPresentation} (\ref{PqRestorePcPresentation})). Following the colon \mbox{\texttt{\mdseries\slshape options}} is a selection of the options listed for the
non\texttt{\symbol{45}}interactive \texttt{Pq} function (see{\nobreakspace}\texttt{Pq} (\ref{Pq})), separated by commas like record components (see Section{\nobreakspace} \textbf{Reference: Function Call With Options} in the \textsf{GAP} Reference Manual), except that the options \texttt{SetupFile} or \texttt{PqWorkspace} are ignored by the interactive \texttt{Pq}, and \texttt{RedoPcp} is an option only recognised by the interactive \texttt{Pq} i.e.{\nobreakspace}the following options are recognised by the interactive \texttt{Pq} function: 
\begin{itemize}
\item  \texttt{Prime := \mbox{\texttt{\mdseries\slshape p}}}\index{option Prime} 
\item  \texttt{ClassBound := \mbox{\texttt{\mdseries\slshape n}}}\index{option ClassBound} 
\item  \texttt{Exponent := \mbox{\texttt{\mdseries\slshape n}}}\index{option Exponent} 
\item  \texttt{Relators := \mbox{\texttt{\mdseries\slshape rels}}}\index{option Relators} 
\item  \texttt{Metabelian}\index{option Metabelian} 
\item  \texttt{Identities := \mbox{\texttt{\mdseries\slshape funcs}}}\index{option Identities} 
\item  \texttt{GroupName := \mbox{\texttt{\mdseries\slshape name}}}\index{option GroupName} 
\item  \texttt{OutputLevel := \mbox{\texttt{\mdseries\slshape n}}}\index{option OutputLevel} 
\item  \texttt{RedoPcp}\index{option RedoPcp} 
\end{itemize}
 Detailed descriptions of the above options may be found in
Chapter{\nobreakspace}\hyperref[ANUPQ Options]{`ANUPQ Options'}. 

 As a minimum the \texttt{Pq} function \emph{must} have a value for the \texttt{Prime} option, though \texttt{Prime} need not be passed again in the case it has previously been provided, e.g. to \texttt{PqStart} (see{\nobreakspace}\texttt{PqStart} (\ref{PqStart})) when starting the interactive process. 

 The behaviour of the interactive \texttt{Pq} function depends on the current state of the pc presentation stored by the \texttt{pq} program: 
\begin{enumerate}
\item  If no pc presentation has yet been computed (the case immediately after the \texttt{PqStart} call initiating the process) then the quotient group of the input group of the
process of largest lower exponent\texttt{\symbol{45}}\mbox{\texttt{\mdseries\slshape p}} class bounded by the value of the \texttt{ClassBound} option (see{\nobreakspace}\ref{option ClassBound}) is returned. 
\item  If the current pc presentation of the process was determined by a previous
call to \texttt{Pq} or \texttt{PqEpimorphism}, and the current call has a larger value \texttt{ClassBound} then the class is extended as much as is possible and the quotient group of
the input group of the process of the new lower exponent\texttt{\symbol{45}}\mbox{\texttt{\mdseries\slshape p}} class is returned. 
\item  If the current pc presentation of the process was determined by a previous
call to \texttt{PqPCover} then a consistent pc presentation of a quotient for the current class is
determined before proceeding as in 2. 
\item  If the \texttt{RedoPcp} option is supplied the current pc presentation is scrapped, all options must
be re\texttt{\symbol{45}}supplied (in particular, \texttt{Prime} \emph{must} be supplied) and then the \texttt{Pq} function proceeds as in 1. 
\end{enumerate}
 See Section{\nobreakspace}\hyperref[Attributes and a Property for fp and pc p-groups]{`Attributes and a Property for fp and pc p\texttt{\symbol{45}}groups'} for the attributes and property \texttt{NuclearRank}, \texttt{MultiplicatorRank} and \texttt{IsCapable} which may be applied to the group returned by \texttt{Pq}. 

 The following is one of the examples for the
non\texttt{\symbol{45}}interactive \texttt{Pq} redone with the interactive version. Also, we set the option \texttt{OutputLevel} to 1 (see{\nobreakspace}\ref{option OutputLevel}), in order to see the orders of the quotients of all the classes determined,
and we set the \texttt{InfoANUPQ} level to 2 (see{\nobreakspace}\texttt{InfoANUPQ} (\ref{InfoANUPQ})), so that we catch the timing information.  
\begin{Verbatim}[commandchars=!@|,fontsize=\small,frame=single,label=Example]
  !gapprompt@gap>| !gapinput@F := FreeGroup("a", "b");; a := F.1;; b := F.2;;|
  !gapprompt@gap>| !gapinput@G := F / [a^4, b^4];|
  <fp group on the generators [ a, b ]>
  !gapprompt@gap>| !gapinput@PqStart(G);|
  1
  !gapprompt@gap>| !gapinput@SetInfoLevel(InfoANUPQ, 2); #To see timing information|
  !gapprompt@gap>| !gapinput@Pq(: Prime := 2, ClassBound := 3, OutputLevel := 1 );|
  #I  Lower exponent-2 central series for [grp]
  #I  Group: [grp] to lower exponent-2 central class 1 has order 2^2
  #I  Group: [grp] to lower exponent-2 central class 2 has order 2^5
  #I  Group: [grp] to lower exponent-2 central class 3 has order 2^8
  #I  Computation of presentation took 0.00 seconds
  <pc group of size 256 with 8 generators>
\end{Verbatim}
 }

 

\subsection{\textcolor{Chapter }{PqEpimorphism (interactive)}}
\logpage{[ 5, 3, 2 ]}\nobreak
\hyperdef{L}{X839E6F578227A8EA}{}
{\noindent\textcolor{FuncColor}{$\triangleright$\enspace\texttt{PqEpimorphism({\mdseries\slshape i: options})\index{PqEpimorphism@\texttt{PqEpimorphism}!interactive}
\label{PqEpimorphism:interactive}
}\hfill{\scriptsize (function)}}\\
\noindent\textcolor{FuncColor}{$\triangleright$\enspace\texttt{PqEpimorphism({\mdseries\slshape : options})\index{PqEpimorphism@\texttt{PqEpimorphism}!interactive, for default process}
\label{PqEpimorphism:interactive, for default process}
}\hfill{\scriptsize (function)}}\\


 return, for the fp or pc group (let us call it \mbox{\texttt{\mdseries\slshape F}}), of the \mbox{\texttt{\mdseries\slshape i}}th or default interactive \textsf{ANUPQ} process, an epimorphism from \mbox{\texttt{\mdseries\slshape F}} onto the $p$\texttt{\symbol{45}}quotient of \mbox{\texttt{\mdseries\slshape F}} specified by \mbox{\texttt{\mdseries\slshape options}}; \mbox{\texttt{\mdseries\slshape F}} must previously have been given (as first argument) to \texttt{PqStart} to start the interactive \textsf{ANUPQ} process (see{\nobreakspace}\texttt{PqStart} (\ref{PqStart})). Since the underlying interactions with the \texttt{pq} program effected by the interactive \texttt{PqEpimorphism} are identical to those effected by the interactive \texttt{Pq}, everything said regarding the requirements and behaviour of the interactive \texttt{Pq} function (see{\nobreakspace}\texttt{Pq} (\ref{Pq:interactive})) is also the case for the interactive \texttt{PqEpimorphism}. 

 \emph{Note:} See Section{\nobreakspace}\hyperref[Attributes and a Property for fp and pc p-groups]{`Attributes and a Property for fp and pc p\texttt{\symbol{45}}groups'} for the attributes and property \texttt{NuclearRank}, \texttt{MultiplicatorRank} and \texttt{IsCapable} which may be applied to the image group of the epimorphism returned by \texttt{PqEpimorphism}. }

 

\subsection{\textcolor{Chapter }{PqPCover (interactive)}}
\logpage{[ 5, 3, 3 ]}\nobreak
\hyperdef{L}{X7AA64A2F8509A39A}{}
{\noindent\textcolor{FuncColor}{$\triangleright$\enspace\texttt{PqPCover({\mdseries\slshape i: options})\index{PqPCover@\texttt{PqPCover}!interactive}
\label{PqPCover:interactive}
}\hfill{\scriptsize (function)}}\\
\noindent\textcolor{FuncColor}{$\triangleright$\enspace\texttt{PqPCover({\mdseries\slshape : options})\index{PqPCover@\texttt{PqPCover}!interactive, for default process}
\label{PqPCover:interactive, for default process}
}\hfill{\scriptsize (function)}}\\


 return, for the fp or pc group of the \mbox{\texttt{\mdseries\slshape i}}th or default interactive \textsf{ANUPQ} process, the $p$\texttt{\symbol{45}}covering group of the $p$\texttt{\symbol{45}}quotient \texttt{Pq(\mbox{\texttt{\mdseries\slshape i}} : \mbox{\texttt{\mdseries\slshape options}})} or \texttt{Pq(: \mbox{\texttt{\mdseries\slshape options}})}, modulo the following: 
\begin{enumerate}
\item  If no pc presentation has yet been computed (the case immediately after the \texttt{PqStart} call initiating the process) and the group \mbox{\texttt{\mdseries\slshape F}} of the process is already a $p$\texttt{\symbol{45}}group, in the sense that \texttt{HasIsPGroup(\mbox{\texttt{\mdseries\slshape F}}) and IsPGroup(\mbox{\texttt{\mdseries\slshape F}})} is \texttt{true}, then 
\begin{description}
\item[{\texttt{Prime}}]  defaults to \texttt{PrimePGroup(\mbox{\texttt{\mdseries\slshape F}})}, if not supplied and \texttt{HasPrimePGroup(\mbox{\texttt{\mdseries\slshape F}}) = true}; and 
\item[{\texttt{ClassBound}}]  defaults to \texttt{PClassPGroup(\mbox{\texttt{\mdseries\slshape F}})} if \texttt{HasPClassPGroup(\mbox{\texttt{\mdseries\slshape F}}) = true} if not supplied, or to the usual default of 63, otherwise. 
\end{description}
 
\item  If a pc presentation has been computed and none of \mbox{\texttt{\mdseries\slshape options}} is \texttt{RedoPcp} or if no pc presentation has yet been computed but 1. does not apply then \texttt{PqPCover(\mbox{\texttt{\mdseries\slshape i}} : \mbox{\texttt{\mdseries\slshape options}});} is equivalent to: 
\begin{Verbatim}[commandchars=!@|,fontsize=\small,frame=single,label=]
  Pq(i : options);
  PqPCover(i);
\end{Verbatim}
 
\item  If the \texttt{RedoPcp} option is supplied the current pc presentation is scrapped, and \texttt{PqPCover} proceeds as in 1. or 2. but without the \texttt{RedoPcp} option. 
\end{enumerate}
 }

 \index{automorphisms!of $p$\texttt{\symbol{45}}groups} 

\subsection{\textcolor{Chapter }{PqStandardPresentation (interactive)}}
\logpage{[ 5, 3, 4 ]}\nobreak
\hyperdef{L}{X805F32618005C087}{}
{\noindent\textcolor{FuncColor}{$\triangleright$\enspace\texttt{PqStandardPresentation({\mdseries\slshape [i]: options})\index{PqStandardPresentation@\texttt{PqStandardPresentation}!interactive}
\label{PqStandardPresentation:interactive}
}\hfill{\scriptsize (function)}}\\
\noindent\textcolor{FuncColor}{$\triangleright$\enspace\texttt{StandardPresentation({\mdseries\slshape [i]: options})\index{StandardPresentation@\texttt{StandardPresentation}!interactive}
\label{StandardPresentation:interactive}
}\hfill{\scriptsize (function)}}\\


 return, for the \mbox{\texttt{\mdseries\slshape i}}th or default interactive \textsf{ANUPQ} process, the \mbox{\texttt{\mdseries\slshape p}}\texttt{\symbol{45}}quotient of the group \mbox{\texttt{\mdseries\slshape F}} of the process, specified by \mbox{\texttt{\mdseries\slshape options}}, as an \emph{fp group} which has a standard presentation. Here \mbox{\texttt{\mdseries\slshape options}} is a selection of the options from the following list (see
Chapter{\nobreakspace}\hyperref[ANUPQ Options]{`ANUPQ Options'} for detailed descriptions); this list is the same as for the
non\texttt{\symbol{45}}interactive version of \texttt{PqStandardPresentation} except for the omission of options \texttt{SetupFile} and \texttt{PqWorkspace} (see{\nobreakspace}\texttt{PqStandardPresentation} (\ref{PqStandardPresentation})). 
\begin{itemize}
\item  \texttt{Prime := \mbox{\texttt{\mdseries\slshape p}}}\index{option Prime} 
\item  \texttt{pQuotient := \mbox{\texttt{\mdseries\slshape Q}}}\index{option pQuotient} 
\item  \texttt{ClassBound := \mbox{\texttt{\mdseries\slshape n}}}\index{option ClassBound} 
\item  \texttt{Exponent := \mbox{\texttt{\mdseries\slshape n}}}\index{option Exponent} 
\item  \texttt{Metabelian}\index{option Metabelian} 
\item  \texttt{GroupName := \mbox{\texttt{\mdseries\slshape name}}}\index{option GroupName} 
\item  \texttt{OutputLevel := \mbox{\texttt{\mdseries\slshape n}}}\index{option OutputLevel} 
\item  \texttt{StandardPresentationFile := \mbox{\texttt{\mdseries\slshape filename}}}\index{option StandardPresentationFile} 
\end{itemize}
 Unless \mbox{\texttt{\mdseries\slshape F}} is a pc \mbox{\texttt{\mdseries\slshape p}}\texttt{\symbol{45}}group, or the option \texttt{Prime} has been passed to a previous interactive function for the process to compute
a \mbox{\texttt{\mdseries\slshape p}}\texttt{\symbol{45}}quotient for \mbox{\texttt{\mdseries\slshape F}}, the user \emph{must} supply either the option \texttt{Prime} or the option \texttt{pQuotient} (if both \texttt{Prime} and \texttt{pQuotient} are supplied, the prime \mbox{\texttt{\mdseries\slshape p}} is determined by applying \texttt{PrimePGroup} (see{\nobreakspace}\texttt{PrimePGroup} (\textbf{Reference: PrimePGroup}) in the Reference Manual) to the value of \texttt{pQuotient}). 

 Taking one of the examples for the non\texttt{\symbol{45}}interactive version
of \texttt{StandardPresentation} (see{\nobreakspace}\texttt{StandardPresentation} (\ref{StandardPresentation})) that required two separate calls to the \texttt{pq} program, we now show how it can be done by setting up a dialogue with just the
one \texttt{pq} process, using the interactive version of \texttt{StandardPresentation}: 
\begin{Verbatim}[commandchars=!@|,fontsize=\small,frame=single,label=Example]
  !gapprompt@gap>| !gapinput@F4 := FreeGroup( "a", "b", "c", "d" );;|
  !gapprompt@gap>| !gapinput@a := F4.1;; b := F4.2;; c := F4.3;; d := F4.4;;|
  !gapprompt@gap>| !gapinput@G4 := F4 / [ b^4, b^2 / Comm(Comm (b, a), a), d^16,|
  !gapprompt@>| !gapinput@                a^16 / (c * d), b^8 / (d * c^4) ];|
  <fp group on the generators [ a, b, c, d ]>
  !gapprompt@gap>| !gapinput@SetInfoLevel(InfoANUPQ, 1); #Only essential Info please|
  !gapprompt@gap>| !gapinput@procId := PqStart(G4);; #Start a new interactive process for a new group|
  !gapprompt@gap>| !gapinput@K := Pq( procId : Prime := 2, ClassBound := 1 );|
  <pc group of size 4 with 2 generators>
  !gapprompt@gap>| !gapinput@StandardPresentation( procId : pQuotient := K, ClassBound := 14 );|
  <fp group with 53 generators>
\end{Verbatim}
 \emph{Notes} 

 In contrast to the function \texttt{Pq} (see{\nobreakspace}\texttt{Pq} (\ref{Pq})) which returns a pc group, \texttt{PqStandardPresentation} or \texttt{StandardPresentation} returns an fp group. This is because the output is mainly used for isomorphism
testing for which an fp group is enough. However, the presentation is a
polycyclic presentation and if you need to do any further computation with
this group (e.g.{\nobreakspace}to find the order) you can use the function \texttt{PcGroupFpGroup} (see{\nobreakspace}\texttt{PcGroupFpGroup} (\textbf{Reference: PcGroupFpGroup}) in the \textsf{GAP} Reference Manual) to form a pc group. 

 If the user does not supply a \mbox{\texttt{\mdseries\slshape p}}\texttt{\symbol{45}}quotient \mbox{\texttt{\mdseries\slshape Q}} via the \texttt{pQuotient} option, and the prime \mbox{\texttt{\mdseries\slshape p}} is either supplied, stored, or \mbox{\texttt{\mdseries\slshape F}} is a pc \mbox{\texttt{\mdseries\slshape p}}\texttt{\symbol{45}}group, then a \mbox{\texttt{\mdseries\slshape p}}\texttt{\symbol{45}}quotient \mbox{\texttt{\mdseries\slshape Q}} is computed. (The value of the prime \mbox{\texttt{\mdseries\slshape p}} is stored if passed initially to \texttt{PqStart} or to a subsequent interactive process.) Note that a stored value for \texttt{pQuotient} (from a prior call to \texttt{Pq}) does \emph{not} have precedence over a value for the prime \mbox{\texttt{\mdseries\slshape p}}. If the user does supply a \mbox{\texttt{\mdseries\slshape p}}\texttt{\symbol{45}}quotient \mbox{\texttt{\mdseries\slshape Q}} via the \texttt{pQuotient} option, the package \textsf{AutPGrp} is called to compute the automorphism group of \mbox{\texttt{\mdseries\slshape Q}}; an error will occur that asks the user to install the package \textsf{AutPGrp} if the automorphism group cannot be computed. 

 If any of the interactive functions \texttt{PqStandardPresentation}, \texttt{StandardPresentation}, \texttt{EpimorphismPqStandardPresentation} or \texttt{EpimorphismStandardPresentation} has been called previously for an interactive process, a subsequent call to
any of these functions for the same process returns the previously computed
value. Note that all these functions compute both an epimorphism and an fp
group and store the results in the \texttt{SPepi} and \texttt{SP} fields of the data record associated with the process. See the example for the
interactive \texttt{EpimorphismStandardPresentation} (\texttt{EpimorphismStandardPresentation} (\ref{EpimorphismStandardPresentation:interactive})). 

 The attributes and property \texttt{NuclearRank}, \texttt{MultiplicatorRank} and \texttt{IsCapable} are set for the group returned by \texttt{PqStandardPresentation} or \texttt{StandardPresentation} (see Section{\nobreakspace}\hyperref[Attributes and a Property for fp and pc p-groups]{`Attributes and a Property for fp and pc p\texttt{\symbol{45}}groups'}). }

 

\subsection{\textcolor{Chapter }{EpimorphismPqStandardPresentation (interactive)}}
\logpage{[ 5, 3, 5 ]}\nobreak
\hyperdef{L}{X791392977B2D692D}{}
{\noindent\textcolor{FuncColor}{$\triangleright$\enspace\texttt{EpimorphismPqStandardPresentation({\mdseries\slshape [i]: options})\index{EpimorphismPqStandardPresentation@\texttt{EpimorphismPqStandardPresentation}!interactive}
\label{EpimorphismPqStandardPresentation:interactive}
}\hfill{\scriptsize (function)}}\\
\noindent\textcolor{FuncColor}{$\triangleright$\enspace\texttt{EpimorphismStandardPresentation({\mdseries\slshape [i]: options})\index{EpimorphismStandardPresentation@\texttt{EpimorphismStandardPresentation}!interactive}
\label{EpimorphismStandardPresentation:interactive}
}\hfill{\scriptsize (method)}}\\


 Each of the above functions accepts the same arguments and options as the
interactive form of \texttt{StandardPresentation} (see{\nobreakspace}\texttt{StandardPresentation} (\ref{StandardPresentation:interactive})) and returns an epimorphism from the fp or pc group \mbox{\texttt{\mdseries\slshape F}} of the \mbox{\texttt{\mdseries\slshape i}}th or default interactive \textsf{ANUPQ} process onto the finitely presented group given by a standard presentation,
i.e.{\nobreakspace}if \mbox{\texttt{\mdseries\slshape S}} is the standard presentation computed for the $p$\texttt{\symbol{45}}quotient of \mbox{\texttt{\mdseries\slshape F}} by \texttt{StandardPresentation} then \texttt{EpimorphismStandardPresentation} returns the epimorphism from \mbox{\texttt{\mdseries\slshape F}} to the group with presentation \mbox{\texttt{\mdseries\slshape S}}. The group \mbox{\texttt{\mdseries\slshape F}} must have been given (as first argument) to \texttt{PqStart} to start the interactive \textsf{ANUPQ} process (see{\nobreakspace}\texttt{PqStart} (\ref{PqStart})). 

 Taking our earlier non\texttt{\symbol{45}}interactive example
(see{\nobreakspace}\texttt{EpimorphismPqStandardPresentation} (\ref{EpimorphismPqStandardPresentation})) and modifying it a little, we illustrate, as for the interactive \texttt{StandardPresentation} (see{\nobreakspace}\texttt{StandardPresentation} (\ref{StandardPresentation:interactive})), how something that required two separate calls to the \texttt{pq} program can now be achieved with a dialogue with just one \texttt{pq} process. Also, observe that calls to one of the standard presentation
functions (as mentioned in the notes of{\nobreakspace}\texttt{StandardPresentation} (\ref{StandardPresentation:interactive})) computes and stores both an fp group with a standard presentation and an
epimorphism; subsequent calls to a standard presentation function for the same
process simply return the appropriate stored value.  
\begin{Verbatim}[commandchars=!@|,fontsize=\small,frame=single,label=Example]
  !gapprompt@gap>| !gapinput@F := FreeGroup(6, "F");;|
  !gapprompt@gap>| !gapinput@x := F.1;; y := F.2;; z := F.3;; w := F.4;; a := F.5;; b := F.6;;|
  !gapprompt@gap>| !gapinput@R := [x^3 / w, y^3 / w * a^2 * b^2, w^3 / b,|
  !gapprompt@>| !gapinput@         Comm (y, x) / z, Comm (z, x), Comm (z, y) / a, z^3 ];|
  [ F1^3*F4^-1, F2^3*F4^-1*F5^2*F6^2, F4^3*F6^-1, F2^-1*F1^-1*F2*F1*F3^-1, 
    F3^-1*F1^-1*F3*F1, F3^-1*F2^-1*F3*F2*F5^-1, F3^3 ]
  !gapprompt@gap>| !gapinput@Q := F / R;|
  <fp group on the generators [ F1, F2, F3, F4, F5, F6 ]>
  !gapprompt@gap>| !gapinput@procId := PqStart( Q );;|
  !gapprompt@gap>| !gapinput@G := Pq( procId : Prime := 3, ClassBound := 3 );|
  <pc group of size 729 with 6 generators>
  !gapprompt@gap>| !gapinput@lev := InfoLevel(InfoANUPQ);; # Save current InfoANUPQ level|
  !gapprompt@gap>| !gapinput@SetInfoLevel(InfoANUPQ, 2); # To see computation times|
  !gapprompt@gap>| !gapinput@# It is not necessary to pass the `Prime' option to|
  !gapprompt@gap>| !gapinput@# `EpimorphismStandardPresentation' since it was previously|
  !gapprompt@gap>| !gapinput@# passed to `Pq':|
  !gapprompt@gap>| !gapinput@phi := EpimorphismStandardPresentation( 3 : ClassBound := 3 );|
  #I  Class 1 3-quotient and its 3-covering group computed in 0.00 seconds
  #I  Order of GL subgroup is 48
  #I  No. of soluble autos is 0
  #I    dim U = 1  dim N = 3  dim M = 3
  #I    nice stabilizer with perm rep
  #I  Computing standard presentation for class 2 took 0.00 seconds
  #I  Computing standard presentation for class 3 took 0.01 seconds
  [ F1, F2, F3, F4, F5, F6 ] -> [ f1*f2^2*f3*f4^2*f5^2, f1*f2*f3*f5, f3^2, 
    f4*f6^2, f5, f6 ]
  !gapprompt@gap>| !gapinput@# Image of phi should be isomorphic to G ...|
  !gapprompt@gap>| !gapinput@# let's check the order is correct:|
  !gapprompt@gap>| !gapinput@Size( Image(phi) );|
  729
  !gapprompt@gap>| !gapinput@# `StandardPresentation' and `EpimorphismStandardPresentation'|
  !gapprompt@gap>| !gapinput@# behave like attributes, so no computation is done when|
  !gapprompt@gap>| !gapinput@# either is called again for the same process ...|
  !gapprompt@gap>| !gapinput@StandardPresentation( 3 : ClassBound := 3 );|
  <fp group of size 729 on the generators [ f1, f2, f3, f4, f5, f6 ]>
  !gapprompt@gap>| !gapinput@# No timing data was Info-ed since no computation was done|
  !gapprompt@gap>| !gapinput@SetInfoLevel(InfoANUPQ, lev); # Restore previous InfoANUPQ level|
\end{Verbatim}
 A very similar (essential details are the same) example to the above may be
executed live, by typing: \texttt{PqExample( "EpimorphismStandardPresentation\texttt{\symbol{45}}i" );}. 

 \emph{Note:} The notes for \texttt{PqStandardPresentation} or \texttt{StandardPresentation} (see{\nobreakspace}\texttt{PqStandardPresentation} (\ref{PqStandardPresentation:interactive})) apply also to \texttt{EpimorphismPqStandardPresentation} or \texttt{EpimorphismStandardPresentation} except that their return value is an \emph{epimorphism onto} an fp group, i.e.{\nobreakspace}one should interpret the phrase ``returns an fp group'' as ``returns an epimorphism onto an fp group'' etc. }

 

\subsection{\textcolor{Chapter }{PqDescendants (interactive)}}
\logpage{[ 5, 3, 6 ]}\nobreak
\hyperdef{L}{X795817217C53AB89}{}
{\noindent\textcolor{FuncColor}{$\triangleright$\enspace\texttt{PqDescendants({\mdseries\slshape i: options})\index{PqDescendants@\texttt{PqDescendants}!interactive}
\label{PqDescendants:interactive}
}\hfill{\scriptsize (function)}}\\
\noindent\textcolor{FuncColor}{$\triangleright$\enspace\texttt{PqDescendants({\mdseries\slshape : options})\index{PqDescendants@\texttt{PqDescendants}!interactive, for default process}
\label{PqDescendants:interactive, for default process}
}\hfill{\scriptsize (function)}}\\


 return for the pc group \mbox{\texttt{\mdseries\slshape G}} of the \mbox{\texttt{\mdseries\slshape i}}th or default interactive \textsf{ANUPQ} process, which must be of prime power order with a confluent pc presentation
(see{\nobreakspace}\texttt{IsConfluent} (\textbf{Reference: IsConfluent for pc groups}) in the \textsf{GAP} Reference Manual), a list of proper descendants (pc groups) of \mbox{\texttt{\mdseries\slshape G}}. The group \mbox{\texttt{\mdseries\slshape G}} is usually given as first argument to \texttt{PqStart} when starting the interactive \textsf{ANUPQ} process (see{\nobreakspace}\texttt{PqStart} (\ref{PqStart})). Alternatively, one may initiate the process with an fp group, use \texttt{Pq} interactively (see{\nobreakspace}\texttt{Pq} (\ref{Pq:interactive})) to create a pc group and use \texttt{PqSetPQuotientToGroup} (see{\nobreakspace}\texttt{PqSetPQuotientToGroup} (\ref{PqSetPQuotientToGroup})), which involves \emph{no} computation, to set the pc group returned by \texttt{Pq} as the group of the process. Note that repeating a call to \texttt{PqDescendants} for the same interactive \textsf{ANUPQ} process simply returns the list of descendants originally calculated; a
warning is emitted at \texttt{InfoANUPQ} level 1 reminding you of this should you do this. 

 After the colon, \mbox{\texttt{\mdseries\slshape options}} a selection of the options listed for the non\texttt{\symbol{45}}interactive \texttt{PqDescendants} function (see{\nobreakspace}\texttt{PqDescendants} (\ref{PqDescendants})), should be given, separated by commas like record components (see
Section{\nobreakspace} \textbf{Reference: Function Call With Options} in the \textsf{GAP} Reference Manual), except that the options \texttt{SetupFile} or \texttt{PqWorkspace} are ignored by the interactive \texttt{PqDescendants}, i.e.{\nobreakspace}the following options are recognised by the interactive \texttt{PqDescendants} function: 
\begin{itemize}
\item  \texttt{ClassBound := \mbox{\texttt{\mdseries\slshape n}}}\index{option ClassBound} 
\item  \texttt{Relators := \mbox{\texttt{\mdseries\slshape rels}}}\index{option Relators} 
\item  \texttt{OrderBound := \mbox{\texttt{\mdseries\slshape n}}}\index{option OrderBound} 
\item  \texttt{StepSize := \mbox{\texttt{\mdseries\slshape n}}}, \texttt{StepSize := \mbox{\texttt{\mdseries\slshape list}}} \index{option StepSize} 
\item  \texttt{RankInitialSegmentSubgroups := \mbox{\texttt{\mdseries\slshape n}}}\index{option RankInitialSegmentSubgroups} 
\item  \texttt{SpaceEfficient}\index{option SpaceEfficient} 
\item  \texttt{CapableDescendants}\index{option CapableDescendants} 
\item  \texttt{AllDescendants := false}\index{option AllDescendants} 
\item  \texttt{Exponent := \mbox{\texttt{\mdseries\slshape n}}}\index{option Exponent} 
\item  \texttt{Metabelian}\index{option Metabelian} 
\item  \texttt{GroupName := \mbox{\texttt{\mdseries\slshape name}}}\index{option GroupName} 
\item  \texttt{SubList := \mbox{\texttt{\mdseries\slshape sub}}}\index{option SubList} 
\item  \texttt{BasicAlgorithm}\index{option BasicAlgorithm} 
\item  \texttt{CustomiseOutput := \mbox{\texttt{\mdseries\slshape rec}}}\index{option CustomiseOutput} 
\end{itemize}
 \emph{Notes:} The function \texttt{PqDescendants} uses the automorphism group of \mbox{\texttt{\mdseries\slshape G}} which it computes via the package \textsf{AutPGrp} if the automorphism group of \mbox{\texttt{\mdseries\slshape G}} is not already present. If \textsf{AutPGrp} is not installed an error may be raised. If the automorphism group of \mbox{\texttt{\mdseries\slshape G}} is insoluble the \texttt{pq} program will call \textsf{GAP} together with the \textsf{AutPGrp} package for certain orbit\texttt{\symbol{45}}stabilizer calculations. 

 The attributes and property \texttt{NuclearRank}, \texttt{MultiplicatorRank} and \texttt{IsCapable} are set for each group of the list returned by \texttt{PqDescendants} (see Section{\nobreakspace}\hyperref[Attributes and a Property for fp and pc p-groups]{`Attributes and a Property for fp and pc p\texttt{\symbol{45}}groups'}). 

 Let us now repeat the examples previously given for the
non\texttt{\symbol{45}}interactive \texttt{PqDescendants}, but this time with the interactive version of \texttt{PqDescendants}: 
\begin{Verbatim}[commandchars=!@|,fontsize=\small,frame=single,label=Example]
  !gapprompt@gap>| !gapinput@F := FreeGroup( "a", "b" );; a := F.1;; b := F.2;;|
  !gapprompt@gap>| !gapinput@G := PcGroupFpGroup( F / [ a^2, b^2, Comm(b, a) ] );|
  <pc group of size 4 with 2 generators>
  !gapprompt@gap>| !gapinput@procId := PqStart(G);;|
  !gapprompt@gap>| !gapinput@des := PqDescendants( procId : OrderBound := 6, ClassBound := 5 );;|
  !gapprompt@gap>| !gapinput@Length(des);|
  83
  !gapprompt@gap>| !gapinput@List(des, Size);|
  [ 8, 8, 8, 16, 16, 16, 32, 16, 16, 16, 16, 16, 32, 32, 64, 64, 32, 32, 32, 
    32, 32, 32, 32, 64, 64, 64, 64, 64, 64, 64, 64, 64, 64, 64, 32, 32, 32, 32, 
    64, 64, 64, 64, 64, 64, 64, 64, 64, 64, 64, 32, 32, 32, 32, 32, 64, 64, 64, 
    64, 64, 64, 64, 64, 64, 64, 64, 64, 64, 64, 64, 64, 64, 64, 64, 64, 64, 64, 
    64, 64, 64, 64, 64, 64, 64 ]
  !gapprompt@gap>| !gapinput@List(des, d -> Length( PCentralSeries( d, 2 ) ) - 1 );|
  [ 2, 2, 2, 2, 2, 2, 2, 3, 3, 3, 3, 3, 3, 3, 3, 3, 3, 3, 3, 3, 3, 3, 3, 3, 3, 
    3, 3, 3, 3, 3, 3, 3, 3, 3, 3, 3, 3, 3, 3, 3, 3, 3, 3, 3, 3, 3, 3, 3, 3, 4, 
    4, 4, 4, 4, 4, 4, 4, 4, 4, 4, 4, 4, 4, 4, 4, 4, 4, 4, 4, 4, 4, 4, 4, 4, 4, 
    4, 4, 4, 5, 5, 5, 5, 5 ]
\end{Verbatim}
 In the second example we compute all capable descendants of order 27 of the
elementary abelian group of order 9. 
\begin{Verbatim}[commandchars=!@|,fontsize=\small,frame=single,label=Example]
  !gapprompt@gap>| !gapinput@F := FreeGroup( 2, "g" );;|
  !gapprompt@gap>| !gapinput@G := PcGroupFpGroup( F / [ F.1^3, F.2^3, Comm(F.1, F.2) ] );|
  <pc group of size 9 with 2 generators>
  !gapprompt@gap>| !gapinput@procId := PqStart(G);;|
  !gapprompt@gap>| !gapinput@des := PqDescendants( procId : OrderBound := 3, ClassBound := 2,|
  !gapprompt@>| !gapinput@                             CapableDescendants );|
  [ <pc group of size 27 with 3 generators>, 
    <pc group of size 27 with 3 generators> ]
  !gapprompt@gap>| !gapinput@List(des, d -> Length( PCentralSeries( d, 3 ) ) - 1 );|
  [ 2, 2 ]
  !gapprompt@gap>| !gapinput@# For comparison let us now compute all proper descendants|
  !gapprompt@gap>| !gapinput@# (using the non-interactive Pq function)|
  !gapprompt@gap>| !gapinput@PqDescendants( G : OrderBound := 3, ClassBound := 2);|
  [ <pc group of size 27 with 3 generators>, 
    <pc group of size 27 with 3 generators>, 
    <pc group of size 27 with 3 generators> ]
\end{Verbatim}
 In the third example, we compute all proper capable descendants of the
elementary abelian group of order $5^2$ which have exponent\texttt{\symbol{45}}$5$ class at most $3$, exponent $5$, and are metabelian. 
\begin{Verbatim}[commandchars=!@|,fontsize=\small,frame=single,label=Example]
  !gapprompt@gap>| !gapinput@F := FreeGroup( 2, "g" );;|
  !gapprompt@gap>| !gapinput@G := PcGroupFpGroup( F / [ F.1^5, F.2^5, Comm(F.2, F.1) ] );|
  <pc group of size 25 with 2 generators>
  !gapprompt@gap>| !gapinput@procId := PqStart(G);;|
  !gapprompt@gap>| !gapinput@des := PqDescendants( procId : Metabelian, ClassBound := 3,|
  !gapprompt@>| !gapinput@                             Exponent := 5, CapableDescendants );|
  [ <pc group of size 125 with 3 generators>, 
    <pc group of size 625 with 4 generators>, 
    <pc group of size 3125 with 5 generators> ]
  !gapprompt@gap>| !gapinput@List(des, d -> Length( PCentralSeries( d, 5 ) ) - 1 );|
  [ 2, 3, 3 ]
  !gapprompt@gap>| !gapinput@List(des, d -> Length( DerivedSeries( d ) ) );|
  [ 3, 3, 3 ]
  !gapprompt@gap>| !gapinput@List(des, d -> Maximum( List( d, Order ) ) );|
  [ 5, 5, 5 ]
\end{Verbatim}
 }

 

\subsection{\textcolor{Chapter }{PqSetPQuotientToGroup}}
\logpage{[ 5, 3, 7 ]}\nobreak
\hyperdef{L}{X80CCF6D379E5A3B4}{}
{\noindent\textcolor{FuncColor}{$\triangleright$\enspace\texttt{PqSetPQuotientToGroup({\mdseries\slshape i})\index{PqSetPQuotientToGroup@\texttt{PqSetPQuotientToGroup}}
\label{PqSetPQuotientToGroup}
}\hfill{\scriptsize (function)}}\\
\noindent\textcolor{FuncColor}{$\triangleright$\enspace\texttt{PqSetPQuotientToGroup({\mdseries\slshape })\index{PqSetPQuotientToGroup@\texttt{PqSetPQuotientToGroup}!for default process}
\label{PqSetPQuotientToGroup:for default process}
}\hfill{\scriptsize (function)}}\\


 for the \mbox{\texttt{\mdseries\slshape i}}th or default interactive \textsf{ANUPQ} process, set the $p$\texttt{\symbol{45}}quotient previously computed by the interactive \texttt{Pq} function (see{\nobreakspace}\texttt{Pq} (\ref{Pq:interactive})) to be the group of the process. This function is supplied to enable the
computation of descendants of a $p$\texttt{\symbol{45}}quotient that is already known to the \texttt{pq} program, via the interactive \texttt{PqDescendants} function (see{\nobreakspace}\texttt{PqDescendants} (\ref{PqDescendants:interactive})), thus avoiding the need to re\texttt{\symbol{45}}submit it and have the \texttt{pq} program recompute it. 

 \emph{Note:} See the function \texttt{PqPGSetDescendantToPcp} (\texttt{PqPGSetDescendantToPcp} (\ref{PqPGSetDescendantToPcp})) for a mechanism to make (the $p$\texttt{\symbol{45}}cover of) a particular descendants the current group of
the process. 

 The following example of the usage of \texttt{PqSetPQuotientToGroup}, which is essentially equivalent to what is obtained by running \texttt{PqExample("PqDescendants\texttt{\symbol{45}}1\texttt{\symbol{45}}i");}, redoes the first example of \texttt{PqDescendants} (\ref{PqDescendants:interactive}) (which computes the descendants of the Klein four group). 
\begin{Verbatim}[commandchars=!@|,fontsize=\small,frame=single,label=Example]
  !gapprompt@gap>| !gapinput@F := FreeGroup( "a", "b" );|
  <free group on the generators [ a, b ]>
  !gapprompt@gap>| !gapinput@procId := PqStart( F : Prime := 2 );;|
  !gapprompt@gap>| !gapinput@Pq( procId : ClassBound := 1 );|
  <pc group of size 4 with 2 generators>
  !gapprompt@gap>| !gapinput@PqSetPQuotientToGroup( procId );|
  !gapprompt@gap>| !gapinput@des := PqDescendants( procId : OrderBound := 6, ClassBound := 5 );;|
  !gapprompt@gap>| !gapinput@Length(des);|
  83
  !gapprompt@gap>| !gapinput@List(des, Size);|
  [ 8, 8, 8, 16, 16, 16, 32, 16, 16, 16, 16, 16, 32, 32, 64, 64, 32, 32, 32, 
    32, 32, 32, 32, 64, 64, 64, 64, 64, 64, 64, 64, 64, 64, 64, 32, 32, 32, 32, 
    64, 64, 64, 64, 64, 64, 64, 64, 64, 64, 64, 32, 32, 32, 32, 32, 64, 64, 64, 
    64, 64, 64, 64, 64, 64, 64, 64, 64, 64, 64, 64, 64, 64, 64, 64, 64, 64, 64, 
    64, 64, 64, 64, 64, 64, 64 ]
  !gapprompt@gap>| !gapinput@List(des, d -> Length( PCentralSeries( d, 2 ) ) - 1 );|
  [ 2, 2, 2, 2, 2, 2, 2, 3, 3, 3, 3, 3, 3, 3, 3, 3, 3, 3, 3, 3, 3, 3, 3, 3, 3, 
    3, 3, 3, 3, 3, 3, 3, 3, 3, 3, 3, 3, 3, 3, 3, 3, 3, 3, 3, 3, 3, 3, 3, 3, 4, 
    4, 4, 4, 4, 4, 4, 4, 4, 4, 4, 4, 4, 4, 4, 4, 4, 4, 4, 4, 4, 4, 4, 4, 4, 4, 
    4, 4, 4, 5, 5, 5, 5, 5 ]
\end{Verbatim}
 }

 }

 
\section{\textcolor{Chapter }{Low\texttt{\symbol{45}}level Interactive ANUPQ functions based on menu items
of the pq program}}\label{Low-level Interactive ANUPQ functions based on menu items of the pq program}
\logpage{[ 5, 4, 0 ]}
\hyperdef{L}{X857F050C832A1FE4}{}
{
  The \texttt{pq} program has 5 menus, the details of which the reader will not normally need to
know, but if she wishes to know the details they may be found in the
standalone manual: \texttt{guide.dvi}. Both \texttt{guide.dvi} and the \texttt{pq} program refer to the items of these 5 menus as ``options'', which do \emph{not} correspond in any way to the options used by any of the \textsf{GAP} functions that interface with the \texttt{pq} program. 

 \emph{Warning:} The commands provided in this section are intended to provide something like
the interactive functionality one has when running the standalone, from within \textsf{GAP}. The \texttt{pq} standalone (in particular, its ``advanced'' menus) assumes some expertise of the user; doing the ``wrong'' thing can cause the program to crash. While a number of safeguards have been
provided in the \textsf{GAP} interface to the \texttt{pq} program, these are \emph{not} foolproof, and the user should exercise care and ensure
pre\texttt{\symbol{45}}requisites of the various commands are met. }

 
\section{\textcolor{Chapter }{General commands}}\logpage{[ 5, 5, 0 ]}
\hyperdef{L}{X868B10557D470CF6}{}
{
 The following commands either use a menu item from whatever menu is ``current'' for the \texttt{pq} program, or have general application and are not associated with just one menu
item of the \texttt{pq} program. 

\subsection{\textcolor{Chapter }{PqNrPcGenerators}}
\logpage{[ 5, 5, 1 ]}\nobreak
\hyperdef{L}{X7DE2F6C686C672DD}{}
{\noindent\textcolor{FuncColor}{$\triangleright$\enspace\texttt{PqNrPcGenerators({\mdseries\slshape i})\index{PqNrPcGenerators@\texttt{PqNrPcGenerators}}
\label{PqNrPcGenerators}
}\hfill{\scriptsize (function)}}\\
\noindent\textcolor{FuncColor}{$\triangleright$\enspace\texttt{PqNrPcGenerators({\mdseries\slshape })\index{PqNrPcGenerators@\texttt{PqNrPcGenerators}!for default process}
\label{PqNrPcGenerators:for default process}
}\hfill{\scriptsize (function)}}\\


 for the \mbox{\texttt{\mdseries\slshape i}}th or default interactive \textsf{ANUPQ} process, return the number of pc generators of the lower exponent $p$\texttt{\symbol{45}}class quotient of the group currently determined by the
process. This also applies if the pc presentation is not consistent. }

 

\subsection{\textcolor{Chapter }{PqFactoredOrder}}
\logpage{[ 5, 5, 2 ]}\nobreak
\hyperdef{L}{X87FF98867E8FFB3C}{}
{\noindent\textcolor{FuncColor}{$\triangleright$\enspace\texttt{PqFactoredOrder({\mdseries\slshape i})\index{PqFactoredOrder@\texttt{PqFactoredOrder}}
\label{PqFactoredOrder}
}\hfill{\scriptsize (function)}}\\
\noindent\textcolor{FuncColor}{$\triangleright$\enspace\texttt{PqFactoredOrder({\mdseries\slshape })\index{PqFactoredOrder@\texttt{PqFactoredOrder}!for default process}
\label{PqFactoredOrder:for default process}
}\hfill{\scriptsize (function)}}\\


 for the \mbox{\texttt{\mdseries\slshape i}}th or default interactive \textsf{ANUPQ} process, return an integer pair \texttt{[\mbox{\texttt{\mdseries\slshape p}}, \mbox{\texttt{\mdseries\slshape n}}]} where \mbox{\texttt{\mdseries\slshape p}} is a prime and \mbox{\texttt{\mdseries\slshape n}} is the number of pc generators (see{\nobreakspace}\texttt{PqNrPcGenerators} (\ref{PqNrPcGenerators})) in the pc presentation of the quotient group currently determined by the
process. If this presentation is consistent, then $p^n$ is the order of the quotient group. Otherwise (if tails have been added but
the necessary consistency checks, relation collections, exponent law checks
and redundant generator eliminations have not yet been done), $p^n$ is an upper bound for the order of the group. }

 

\subsection{\textcolor{Chapter }{PqOrder}}
\logpage{[ 5, 5, 3 ]}\nobreak
\hyperdef{L}{X7B4FC9E380001A71}{}
{\noindent\textcolor{FuncColor}{$\triangleright$\enspace\texttt{PqOrder({\mdseries\slshape i})\index{PqOrder@\texttt{PqOrder}}
\label{PqOrder}
}\hfill{\scriptsize (function)}}\\
\noindent\textcolor{FuncColor}{$\triangleright$\enspace\texttt{PqOrder({\mdseries\slshape })\index{PqOrder@\texttt{PqOrder}!for default process}
\label{PqOrder:for default process}
}\hfill{\scriptsize (function)}}\\


 for the \mbox{\texttt{\mdseries\slshape i}}th or default interactive \textsf{ANUPQ} process, return $p^n$ where \texttt{[\mbox{\texttt{\mdseries\slshape p}}, \mbox{\texttt{\mdseries\slshape n}}]} is the pair as returned by \texttt{PqFactoredOrder} (see{\nobreakspace}\texttt{PqFactoredOrder} (\ref{PqFactoredOrder})). }

 

\subsection{\textcolor{Chapter }{PqPClass}}
\logpage{[ 5, 5, 4 ]}\nobreak
\hyperdef{L}{X87C7F7EB7C10DC4B}{}
{\noindent\textcolor{FuncColor}{$\triangleright$\enspace\texttt{PqPClass({\mdseries\slshape i})\index{PqPClass@\texttt{PqPClass}}
\label{PqPClass}
}\hfill{\scriptsize (function)}}\\
\noindent\textcolor{FuncColor}{$\triangleright$\enspace\texttt{PqPClass({\mdseries\slshape })\index{PqPClass@\texttt{PqPClass}!for default process}
\label{PqPClass:for default process}
}\hfill{\scriptsize (function)}}\\


 for the \mbox{\texttt{\mdseries\slshape i}}th or default interactive \textsf{ANUPQ} process, return the lower exponent $p$\texttt{\symbol{45}}class of the quotient group currently determined by the
process. }

 

\subsection{\textcolor{Chapter }{PqWeight}}
\logpage{[ 5, 5, 5 ]}\nobreak
\hyperdef{L}{X7EA3B6917908C20A}{}
{\noindent\textcolor{FuncColor}{$\triangleright$\enspace\texttt{PqWeight({\mdseries\slshape i, j})\index{PqWeight@\texttt{PqWeight}}
\label{PqWeight}
}\hfill{\scriptsize (function)}}\\
\noindent\textcolor{FuncColor}{$\triangleright$\enspace\texttt{PqWeight({\mdseries\slshape j})\index{PqWeight@\texttt{PqWeight}!for default process}
\label{PqWeight:for default process}
}\hfill{\scriptsize (function)}}\\


 for the \mbox{\texttt{\mdseries\slshape i}}th or default interactive \textsf{ANUPQ} process, return the weight of the \mbox{\texttt{\mdseries\slshape j}}th pc generator of the lower exponent $p$\texttt{\symbol{45}}class quotient of the group currently determined by the
process, or \texttt{fail} if there is no such numbered pc generator. }

 

\subsection{\textcolor{Chapter }{PqCurrentGroup}}
\logpage{[ 5, 5, 6 ]}\nobreak
\hyperdef{L}{X79862B83817E20E1}{}
{\noindent\textcolor{FuncColor}{$\triangleright$\enspace\texttt{PqCurrentGroup({\mdseries\slshape i})\index{PqCurrentGroup@\texttt{PqCurrentGroup}}
\label{PqCurrentGroup}
}\hfill{\scriptsize (function)}}\\
\noindent\textcolor{FuncColor}{$\triangleright$\enspace\texttt{PqCurrentGroup({\mdseries\slshape })\index{PqCurrentGroup@\texttt{PqCurrentGroup}!for default process}
\label{PqCurrentGroup:for default process}
}\hfill{\scriptsize (function)}}\\


 for the \mbox{\texttt{\mdseries\slshape i}}th or default interactive \textsf{ANUPQ} process, return the group whose pc presentation is determined by the process
as a \textsf{GAP} pc group (either a lower exponent $p$\texttt{\symbol{45}}class quotient of the start group or the $p$\texttt{\symbol{45}}cover of such a quotient). 

 \emph{Notes:} See Section{\nobreakspace}\hyperref[Attributes and a Property for fp and pc p-groups]{`Attributes and a Property for fp and pc p\texttt{\symbol{45}}groups'} for the attributes and property \texttt{NuclearRank}, \texttt{MultiplicatorRank} and \texttt{IsCapable} which may be applied to the group returned by \texttt{PqCurrentGroup}. }

 

\subsection{\textcolor{Chapter }{PqDisplayPcPresentation}}
\logpage{[ 5, 5, 7 ]}\nobreak
\hyperdef{L}{X805A50687D82B9EC}{}
{\noindent\textcolor{FuncColor}{$\triangleright$\enspace\texttt{PqDisplayPcPresentation({\mdseries\slshape i: [OutputLevel := lev]})\index{PqDisplayPcPresentation@\texttt{PqDisplayPcPresentation}}
\label{PqDisplayPcPresentation}
}\hfill{\scriptsize (function)}}\\
\noindent\textcolor{FuncColor}{$\triangleright$\enspace\texttt{PqDisplayPcPresentation({\mdseries\slshape : [OutputLevel := lev]})\index{PqDisplayPcPresentation@\texttt{PqDisplayPcPresentation}!for default process}
\label{PqDisplayPcPresentation:for default process}
}\hfill{\scriptsize (function)}}\\


 for the \mbox{\texttt{\mdseries\slshape i}}th or default interactive \textsf{ANUPQ} process, direct the \texttt{pq} program to display the pc presentation of the lower exponent $p$\texttt{\symbol{45}}class quotient of the group currently determined by the
process. 

 Except if the last command communicating with the \texttt{pq} program was a $p$\texttt{\symbol{45}}group generation command (for which there is only a
verbose output level), to set the amount of information this command displays
you may wish to call \texttt{PqSetOutputLevel} first (see{\nobreakspace}\texttt{PqSetOutputLevel} (\ref{PqSetOutputLevel})), or equivalently pass the option \texttt{OutputLevel} (see{\nobreakspace}\ref{option OutputLevel}). 

 \emph{Note:} For those familiar with the \texttt{pq} program, \texttt{PqDisplayPcPresentation} performs menu item 4 of the current menu of the \texttt{pq} program. }

 

\subsection{\textcolor{Chapter }{PqSetOutputLevel}}
\logpage{[ 5, 5, 8 ]}\nobreak
\hyperdef{L}{X80410979854280E1}{}
{\noindent\textcolor{FuncColor}{$\triangleright$\enspace\texttt{PqSetOutputLevel({\mdseries\slshape i, lev})\index{PqSetOutputLevel@\texttt{PqSetOutputLevel}}
\label{PqSetOutputLevel}
}\hfill{\scriptsize (function)}}\\
\noindent\textcolor{FuncColor}{$\triangleright$\enspace\texttt{PqSetOutputLevel({\mdseries\slshape lev})\index{PqSetOutputLevel@\texttt{PqSetOutputLevel}!for default process}
\label{PqSetOutputLevel:for default process}
}\hfill{\scriptsize (function)}}\\


 for the \mbox{\texttt{\mdseries\slshape i}}th or default interactive \textsf{ANUPQ} process, direct the \texttt{pq} program to set the output level of the \texttt{pq} program to \mbox{\texttt{\mdseries\slshape lev}}. 

 \emph{Note:} For those familiar with the \texttt{pq} program, \texttt{PqSetOutputLevel} performs menu item 5 of the main (or advanced) $p$\texttt{\symbol{45}}Quotient menu, or the Standard Presentation menu. }

 

\subsection{\textcolor{Chapter }{PqEvaluateIdentities}}
\logpage{[ 5, 5, 9 ]}\nobreak
\hyperdef{L}{X7B8C72B37AE667F7}{}
{\noindent\textcolor{FuncColor}{$\triangleright$\enspace\texttt{PqEvaluateIdentities({\mdseries\slshape i: [Identities := funcs]})\index{PqEvaluateIdentities@\texttt{PqEvaluateIdentities}}
\label{PqEvaluateIdentities}
}\hfill{\scriptsize (function)}}\\
\noindent\textcolor{FuncColor}{$\triangleright$\enspace\texttt{PqEvaluateIdentities({\mdseries\slshape : [Identities := funcs]})\index{PqEvaluateIdentities@\texttt{PqEvaluateIdentities}!for default process}
\label{PqEvaluateIdentities:for default process}
}\hfill{\scriptsize (function)}}\\


 for the \mbox{\texttt{\mdseries\slshape i}}th or default interactive \textsf{ANUPQ} process, invoke the evaluation of identities defined by the \texttt{Identities} option, and eliminate any redundant pc generators formed. Since a previous
value of \texttt{Identities} is saved in the data record of the process, it is unnecessary to pass the \texttt{Identities} if set previously. 

 \emph{Note:} This function is mainly implemented at the \textsf{GAP} level. It does not correspond to a menu item of the \texttt{pq} program. }

 }

 
\section{\textcolor{Chapter }{Commands from the Main $p$\texttt{\symbol{45}}Quotient menu}}\logpage{[ 5, 6, 0 ]}
\hyperdef{L}{X7BDD5B278719C630}{}
{
 

\subsection{\textcolor{Chapter }{PqPcPresentation}}
\logpage{[ 5, 6, 1 ]}\nobreak
\hyperdef{L}{X7BF135DD84A781EB}{}
{\noindent\textcolor{FuncColor}{$\triangleright$\enspace\texttt{PqPcPresentation({\mdseries\slshape i: options})\index{PqPcPresentation@\texttt{PqPcPresentation}}
\label{PqPcPresentation}
}\hfill{\scriptsize (function)}}\\
\noindent\textcolor{FuncColor}{$\triangleright$\enspace\texttt{PqPcPresentation({\mdseries\slshape : options})\index{PqPcPresentation@\texttt{PqPcPresentation}!for default process}
\label{PqPcPresentation:for default process}
}\hfill{\scriptsize (function)}}\\


 for the \mbox{\texttt{\mdseries\slshape i}}th or default interactive \textsf{ANUPQ} process, direct the \texttt{pq} program to compute the pc presentation of the quotient (determined by \mbox{\texttt{\mdseries\slshape options}}) of the group of the process, which for process \mbox{\texttt{\mdseries\slshape i}} is stored as \texttt{ANUPQData.io[\mbox{\texttt{\mdseries\slshape i}}].group}. 

 The possible \mbox{\texttt{\mdseries\slshape options}} are the same as for the interactive \texttt{Pq} (see{\nobreakspace}\texttt{Pq} (\ref{Pq:interactive})) function, except for \texttt{RedoPcp} (which, in any case, would be superfluous), namely: \texttt{Prime}, \texttt{ClassBound}, \texttt{Exponent}, \texttt{Relators}, \texttt{GroupName}, \texttt{Metabelian}, \texttt{Identities} and \texttt{OutputLevel} (see Chapter{\nobreakspace}\hyperref[ANUPQ Options]{`ANUPQ Options'} for a detailed description for these options). The option \texttt{Prime} is required unless already provided to \texttt{PqStart}. 

 \emph{Notes} 

 The pc presentation is held by the \texttt{pq} program. In contrast to \texttt{Pq} (see{\nobreakspace}\texttt{Pq} (\ref{Pq:interactive})), no \textsf{GAP} pc group is returned; see{\nobreakspace}\texttt{PqCurrentGroup} (\texttt{PqCurrentGroup} (\ref{PqCurrentGroup})) if you need the corresponding \textsf{GAP} pc group. 

 \texttt{PqPcPresentation(\mbox{\texttt{\mdseries\slshape i}}: \mbox{\texttt{\mdseries\slshape options}});} is roughly equivalent to the following sequence of
low\texttt{\symbol{45}}level commands: 

 
\begin{Verbatim}[commandchars=!@|,fontsize=\small,frame=single,label=]
  PqPcPresentation(i: opts); #class 1 call
  for c in [2 .. class] do
      PqNextClass(i);
  od;
\end{Verbatim}
 where \mbox{\texttt{\mdseries\slshape opts}} is \mbox{\texttt{\mdseries\slshape options}} except with the \texttt{ClassBound} option set to 1, and \mbox{\texttt{\mdseries\slshape class}} is either the maximum class of a \mbox{\texttt{\mdseries\slshape p}}\texttt{\symbol{45}}quotient of the group of the process or the
user\texttt{\symbol{45}}supplied value of the option \texttt{ClassBound} (whichever is smaller). If the \texttt{Identities} option has been set, both the first \texttt{PqPcPresentation} class 1 call and the \texttt{PqNextClass} calls invoke \texttt{PqEvaluateIdentities(\mbox{\texttt{\mdseries\slshape i}});} as their final step. 

 For those familiar with the \texttt{pq} program, \texttt{PqPcPresentation} performs menu item 1 of the main $p$\texttt{\symbol{45}}Quotient menu. }

 

\subsection{\textcolor{Chapter }{PqSavePcPresentation}}
\logpage{[ 5, 6, 2 ]}\nobreak
\hyperdef{L}{X7D947CEA82C44898}{}
{\noindent\textcolor{FuncColor}{$\triangleright$\enspace\texttt{PqSavePcPresentation({\mdseries\slshape i, filename})\index{PqSavePcPresentation@\texttt{PqSavePcPresentation}}
\label{PqSavePcPresentation}
}\hfill{\scriptsize (function)}}\\
\noindent\textcolor{FuncColor}{$\triangleright$\enspace\texttt{PqSavePcPresentation({\mdseries\slshape filename})\index{PqSavePcPresentation@\texttt{PqSavePcPresentation}!for default process}
\label{PqSavePcPresentation:for default process}
}\hfill{\scriptsize (function)}}\\


 for the \mbox{\texttt{\mdseries\slshape i}}th or default interactive \textsf{ANUPQ} process, direct the \texttt{pq} program to save the pc presentation previously computed for the quotient of
the group of that process to the file with name \mbox{\texttt{\mdseries\slshape filename}}. If the first character of the string \mbox{\texttt{\mdseries\slshape filename}} is not \texttt{/}, \mbox{\texttt{\mdseries\slshape filename}} is assumed to be the path of a writable file relative to the directory in
which \textsf{GAP} was started. A saved file may be restored by \texttt{PqRestorePcPresentation} (see{\nobreakspace}\texttt{PqRestorePcPresentation} (\ref{PqRestorePcPresentation})). 

 \emph{Note:} For those familiar with the \texttt{pq} program, \texttt{PqSavePcPresentation} performs menu item 2 of the main $p$\texttt{\symbol{45}}Quotient menu. }

 

\subsection{\textcolor{Chapter }{PqRestorePcPresentation}}
\logpage{[ 5, 6, 3 ]}\nobreak
\hyperdef{L}{X7FD001C0798EF219}{}
{\noindent\textcolor{FuncColor}{$\triangleright$\enspace\texttt{PqRestorePcPresentation({\mdseries\slshape i, filename})\index{PqRestorePcPresentation@\texttt{PqRestorePcPresentation}}
\label{PqRestorePcPresentation}
}\hfill{\scriptsize (function)}}\\
\noindent\textcolor{FuncColor}{$\triangleright$\enspace\texttt{PqRestorePcPresentation({\mdseries\slshape filename})\index{PqRestorePcPresentation@\texttt{PqRestorePcPresentation}!for default process}
\label{PqRestorePcPresentation:for default process}
}\hfill{\scriptsize (function)}}\\


 for the \mbox{\texttt{\mdseries\slshape i}}th or default interactive \textsf{ANUPQ} process, direct the \texttt{pq} program to restore the pc presentation previously saved to \mbox{\texttt{\mdseries\slshape filename}}, by \texttt{PqSavePcPresentation} (see{\nobreakspace}\texttt{PqSavePcPresentation} (\ref{PqSavePcPresentation})). If the first character of the string \mbox{\texttt{\mdseries\slshape filename}} is not \texttt{/}, \mbox{\texttt{\mdseries\slshape filename}} is assumed to be the path of a readable file relative to the directory in
which \textsf{GAP} was started. 

 \emph{Note:} For those familiar with the \texttt{pq} program, \texttt{PqRestorePcPresentation} performs menu item 3 of the main $p$\texttt{\symbol{45}}Quotient menu. }

 

\subsection{\textcolor{Chapter }{PqNextClass}}
\logpage{[ 5, 6, 4 ]}\nobreak
\hyperdef{L}{X832F69597C095A27}{}
{\noindent\textcolor{FuncColor}{$\triangleright$\enspace\texttt{PqNextClass({\mdseries\slshape i: [QueueFactor]})\index{PqNextClass@\texttt{PqNextClass}}
\label{PqNextClass}
}\hfill{\scriptsize (function)}}\\
\noindent\textcolor{FuncColor}{$\triangleright$\enspace\texttt{PqNextClass({\mdseries\slshape : [QueueFactor]})\index{PqNextClass@\texttt{PqNextClass}!for default process}
\label{PqNextClass:for default process}
}\hfill{\scriptsize (function)}}\\


 for the \mbox{\texttt{\mdseries\slshape i}}th or default interactive \textsf{ANUPQ} process, direct the \texttt{pq} program to calculate the next class of \texttt{ANUPQData.io[\mbox{\texttt{\mdseries\slshape i}}].group}. 

 \index{option QueueFactor} \texttt{PqNextClass} accepts the option \texttt{QueueFactor} (see also{\nobreakspace}\ref{option QueueFactor}) which should be a positive integer if automorphisms have been previously
supplied. If the \texttt{pq} program requires a queue factor and none is supplied via the option \texttt{QueueFactor} a default of 15 is taken. 

 \emph{Notes} 

 The single command: \texttt{PqNextClass(\mbox{\texttt{\mdseries\slshape i}});} is equivalent to executing 
\begin{Verbatim}[commandchars=!@|,fontsize=\small,frame=single,label=]
  PqComputePCover(i);
  PqCollectDefiningRelations(i);
  PqDoExponentChecks(i);
  PqEliminateRedundantGenerators(i);
\end{Verbatim}
 If the \texttt{Identities} option is set the \texttt{PqEliminateRedundantGenerators(\mbox{\texttt{\mdseries\slshape i}});} step is essentially replaced by \texttt{PqEvaluateIdentities(\mbox{\texttt{\mdseries\slshape i}});} (which invokes its own elimination of redundant generators). 

 For those familiar with the \texttt{pq} program, \texttt{PqNextClass} performs menu item 6 of the main $p$\texttt{\symbol{45}}Quotient menu. }

 

\subsection{\textcolor{Chapter }{PqComputePCover}}
\logpage{[ 5, 6, 5 ]}\nobreak
\hyperdef{L}{X7CF9DF7A84D387CB}{}
{\noindent\textcolor{FuncColor}{$\triangleright$\enspace\texttt{PqComputePCover({\mdseries\slshape i})\index{PqComputePCover@\texttt{PqComputePCover}}
\label{PqComputePCover}
}\hfill{\scriptsize (function)}}\\
\noindent\textcolor{FuncColor}{$\triangleright$\enspace\texttt{PqComputePCover({\mdseries\slshape })\index{PqComputePCover@\texttt{PqComputePCover}!for default process}
\label{PqComputePCover:for default process}
}\hfill{\scriptsize (function)}}\\


 for the \mbox{\texttt{\mdseries\slshape i}}th or default interactive \textsf{ANUPQ} processi, directi, the \texttt{pq} program to compute the $p$\texttt{\symbol{45}}covering group of \texttt{ANUPQData.io[\mbox{\texttt{\mdseries\slshape i}}].group}. In contrast to the function \texttt{PqPCover} (see{\nobreakspace}\texttt{PqPCover} (\ref{PqPCover})), this function does not return a \textsf{GAP} pc group. 

 \emph{Notes} 

 The single command: \texttt{PqComputePCover(\mbox{\texttt{\mdseries\slshape i}});} is equivalent to executing 
\begin{Verbatim}[commandchars=!@|,fontsize=\small,frame=single,label=]
  PqSetupTablesForNextClass(i);
  PqTails(i, 0);
  PqDoConsistencyChecks(i, 0, 0);
  PqEliminateRedundantGenerators(i);
\end{Verbatim}
 For those familiar with the \texttt{pq} program, \texttt{PqComputePCover} performs menu item 7 of the main $p$\texttt{\symbol{45}}Quotient menu. }

 }

 
\section{\textcolor{Chapter }{Commands from the Advanced $p$\texttt{\symbol{45}}Quotient menu}}\logpage{[ 5, 7, 0 ]}
\hyperdef{L}{X7D27BE937B1DE16E}{}
{
 

\subsection{\textcolor{Chapter }{PqCollect}}
\logpage{[ 5, 7, 1 ]}\nobreak
\hyperdef{L}{X8633356078CA4115}{}
{\noindent\textcolor{FuncColor}{$\triangleright$\enspace\texttt{PqCollect({\mdseries\slshape i, word})\index{PqCollect@\texttt{PqCollect}}
\label{PqCollect}
}\hfill{\scriptsize (function)}}\\
\noindent\textcolor{FuncColor}{$\triangleright$\enspace\texttt{PqCollect({\mdseries\slshape word})\index{PqCollect@\texttt{PqCollect}!for default process}
\label{PqCollect:for default process}
}\hfill{\scriptsize (function)}}\\


 for the \mbox{\texttt{\mdseries\slshape i}}th or default interactive \textsf{ANUPQ} process, instruct the \texttt{pq} program to do a collection on \mbox{\texttt{\mdseries\slshape word}}, a word in the current pc generators (the form of \mbox{\texttt{\mdseries\slshape word}} required is described below). \texttt{PqCollect} returns the resulting word of the collection as a list of generator number,
exponent pairs (the same form as the second allowed input form of \mbox{\texttt{\mdseries\slshape word}}; see below). 

 The argument \mbox{\texttt{\mdseries\slshape word}} may be input in either of the following ways: 
\begin{enumerate}
\item  \mbox{\texttt{\mdseries\slshape word}} may be a string, where the \mbox{\texttt{\mdseries\slshape i}}th pc generator is represented by \texttt{x\mbox{\texttt{\mdseries\slshape i}}}, e.g.{\nobreakspace}\texttt{"x3*x2\texttt{\symbol{94}}2*x1"}. This way is quite versatile as parentheses and
left\texttt{\symbol{45}}normed commutators
\texttt{\symbol{45}}\texttt{\symbol{45}} using square brackets, in the same
way as \texttt{PqGAPRelators} (see{\nobreakspace}\texttt{PqGAPRelators} (\ref{PqGAPRelators})) \texttt{\symbol{45}}\texttt{\symbol{45}} are permitted; \mbox{\texttt{\mdseries\slshape word}} is checked for correct syntax via \texttt{PqParseWord} (see{\nobreakspace}\texttt{PqParseWord} (\ref{PqParseWord})). 
\item  Otherwise, \mbox{\texttt{\mdseries\slshape word}} must be a list of generator number, exponent pairs of integers,
i.e.{\nobreakspace} each pair represents a ``syllable'' so that \texttt{[ [3, 1], [2, 2], [1, 1] ]} represents the same word as that of the example given for the first allowed
form of \mbox{\texttt{\mdseries\slshape word}}. 
\end{enumerate}
 \emph{Note:} For those familiar with the \texttt{pq} program, \texttt{PqCollect} performs menu item 1 of the Advanced $p$\texttt{\symbol{45}}Quotient menu. }

 

\subsection{\textcolor{Chapter }{PqSolveEquation}}
\logpage{[ 5, 7, 2 ]}\nobreak
\hyperdef{L}{X86CCB77281FDC0FC}{}
{\noindent\textcolor{FuncColor}{$\triangleright$\enspace\texttt{PqSolveEquation({\mdseries\slshape i, a, b})\index{PqSolveEquation@\texttt{PqSolveEquation}}
\label{PqSolveEquation}
}\hfill{\scriptsize (function)}}\\
\noindent\textcolor{FuncColor}{$\triangleright$\enspace\texttt{PqSolveEquation({\mdseries\slshape a, b})\index{PqSolveEquation@\texttt{PqSolveEquation}!for default process}
\label{PqSolveEquation:for default process}
}\hfill{\scriptsize (function)}}\\


 for the \mbox{\texttt{\mdseries\slshape i}}th or default interactive \textsf{ANUPQ} process, direct the \texttt{pq} program to solve $\mbox{\texttt{\mdseries\slshape a}} * \mbox{\texttt{\mdseries\slshape x}} = \mbox{\texttt{\mdseries\slshape b}}$ for \mbox{\texttt{\mdseries\slshape x}}, where \mbox{\texttt{\mdseries\slshape a}} and \mbox{\texttt{\mdseries\slshape b}} are words in the pc generators. For the representation of these words see the
description of the function \texttt{PqCollect} (\texttt{PqCollect} (\ref{PqCollect})). 

 \emph{Note:} For those familiar with the \texttt{pq} program, \texttt{PqSolveEquation} performs menu item 2 of the Advanced $p$\texttt{\symbol{45}}Quotient menu. }

 

\subsection{\textcolor{Chapter }{PqCommutator}}
\logpage{[ 5, 7, 3 ]}\nobreak
\hyperdef{L}{X789B120B7E2F3017}{}
{\noindent\textcolor{FuncColor}{$\triangleright$\enspace\texttt{PqCommutator({\mdseries\slshape i, words, pow})\index{PqCommutator@\texttt{PqCommutator}}
\label{PqCommutator}
}\hfill{\scriptsize (function)}}\\
\noindent\textcolor{FuncColor}{$\triangleright$\enspace\texttt{PqCommutator({\mdseries\slshape words, pow})\index{PqCommutator@\texttt{PqCommutator}!for default process}
\label{PqCommutator:for default process}
}\hfill{\scriptsize (function)}}\\


 for the \mbox{\texttt{\mdseries\slshape i}}th or default interactive \textsf{ANUPQ} process, instruct the \texttt{pq} program to compute the left normed commutator of the list \mbox{\texttt{\mdseries\slshape words}} of words in the current pc generators raised to the integer power \mbox{\texttt{\mdseries\slshape pow}}, and return the resulting word as a list of generator number, exponent pairs.
The form required for each word of \mbox{\texttt{\mdseries\slshape words}} is the same as that required for the \mbox{\texttt{\mdseries\slshape word}} argument of \texttt{PqCollect} (see{\nobreakspace}\texttt{PqCollect} (\ref{PqCollect})). The form of the output word is also the same as for \texttt{PqCollect}. 

 \emph{Note:} For those familiar with the \texttt{pq} program, \texttt{PqCommutator} performs menu item 3 of the Advanced $p$\texttt{\symbol{45}}Quotient menu. }

 

\subsection{\textcolor{Chapter }{PqSetupTablesForNextClass}}
\logpage{[ 5, 7, 4 ]}\nobreak
\hyperdef{L}{X7A61A15E78F52743}{}
{\noindent\textcolor{FuncColor}{$\triangleright$\enspace\texttt{PqSetupTablesForNextClass({\mdseries\slshape i})\index{PqSetupTablesForNextClass@\texttt{PqSetupTablesForNextClass}}
\label{PqSetupTablesForNextClass}
}\hfill{\scriptsize (function)}}\\
\noindent\textcolor{FuncColor}{$\triangleright$\enspace\texttt{PqSetupTablesForNextClass({\mdseries\slshape })\index{PqSetupTablesForNextClass@\texttt{PqSetupTablesForNextClass}!for default process}
\label{PqSetupTablesForNextClass:for default process}
}\hfill{\scriptsize (function)}}\\


 for the \mbox{\texttt{\mdseries\slshape i}}th or default interactive \textsf{ANUPQ} process, direct the \texttt{pq} program to set up tables for the next class. As as
side\texttt{\symbol{45}}effect, after \texttt{PqSetupTablesForNextClass(\mbox{\texttt{\mdseries\slshape i}})} the value returned by \texttt{PqPClass(\mbox{\texttt{\mdseries\slshape i}})} will be one more than it was previously. 

 \emph{Note:} For those familiar with the \texttt{pq} program, \texttt{PqSetupTablesForNextClass} performs menu item 6 of the Advanced $p$\texttt{\symbol{45}}Quotient menu. }

 

\subsection{\textcolor{Chapter }{PqTails}}
\logpage{[ 5, 7, 5 ]}\nobreak
\hyperdef{L}{X84F7DD8582B368D3}{}
{\noindent\textcolor{FuncColor}{$\triangleright$\enspace\texttt{PqTails({\mdseries\slshape i, weight})\index{PqTails@\texttt{PqTails}}
\label{PqTails}
}\hfill{\scriptsize (function)}}\\
\noindent\textcolor{FuncColor}{$\triangleright$\enspace\texttt{PqTails({\mdseries\slshape weight})\index{PqTails@\texttt{PqTails}!for default process}
\label{PqTails:for default process}
}\hfill{\scriptsize (function)}}\\


 for the \mbox{\texttt{\mdseries\slshape i}}th or default interactive \textsf{ANUPQ} process, direct the \texttt{pq} program to compute and add tails of weight \mbox{\texttt{\mdseries\slshape weight}} if \mbox{\texttt{\mdseries\slshape weight}} is in the integer range \texttt{[2 .. PqPClass(\mbox{\texttt{\mdseries\slshape i}})]} (assuming \mbox{\texttt{\mdseries\slshape i}} is the number of the process, even in the default case) or for all weights if \texttt{\mbox{\texttt{\mdseries\slshape weight}} = 0}. 

 If \mbox{\texttt{\mdseries\slshape weight}} is non\texttt{\symbol{45}}zero, then tails that introduce new generators for
only weight \mbox{\texttt{\mdseries\slshape weight}} are computed and added, and in this case and if \texttt{\mbox{\texttt{\mdseries\slshape weight}} {\textless} PqPClass(\mbox{\texttt{\mdseries\slshape i}})}, it is assumed that the tails that introduce new generators for each weight
from \texttt{PqPClass(\mbox{\texttt{\mdseries\slshape i}})} down to weight \texttt{\mbox{\texttt{\mdseries\slshape weight}} + 1} have already been added. You may wish to call \texttt{PqSetMetabelian} (see{\nobreakspace}\texttt{PqSetMetabelian} (\ref{PqSetMetabelian})) prior to calling \texttt{PqTails}. 

 \emph{Notes} 

 For its use in the context of finding the next class see \texttt{PqNextClass} (\ref{PqNextClass}); in particular, a call to \texttt{PqSetupTablesForNextClass} (see{\nobreakspace}\texttt{PqSetupTablesForNextClass} (\ref{PqSetupTablesForNextClass})) needs to have been made prior to calling \texttt{PqTails}. 

 The single command: \texttt{PqTails(\mbox{\texttt{\mdseries\slshape i}}, \mbox{\texttt{\mdseries\slshape weight}});} is equivalent to 
\begin{Verbatim}[commandchars=!@|,fontsize=\small,frame=single,label=]
  PqComputeTails(i, weight);
  PqAddTails(i, weight);
\end{Verbatim}
 For those familiar with the \texttt{pq} program, \texttt{PqTails} uses menu item 7 of the Advanced $p$\texttt{\symbol{45}}Quotient menu. }

 

\subsection{\textcolor{Chapter }{PqComputeTails}}
\logpage{[ 5, 7, 6 ]}\nobreak
\hyperdef{L}{X8389F6437ED634B8}{}
{\noindent\textcolor{FuncColor}{$\triangleright$\enspace\texttt{PqComputeTails({\mdseries\slshape i, weight})\index{PqComputeTails@\texttt{PqComputeTails}}
\label{PqComputeTails}
}\hfill{\scriptsize (function)}}\\
\noindent\textcolor{FuncColor}{$\triangleright$\enspace\texttt{PqComputeTails({\mdseries\slshape weight})\index{PqComputeTails@\texttt{PqComputeTails}!for default process}
\label{PqComputeTails:for default process}
}\hfill{\scriptsize (function)}}\\


 for the \mbox{\texttt{\mdseries\slshape i}}th or default interactive \textsf{ANUPQ} process, direct the \texttt{pq} program to compute tails of weight \mbox{\texttt{\mdseries\slshape weight}} if \mbox{\texttt{\mdseries\slshape weight}} is in the integer range \texttt{[2 .. PqPClass(\mbox{\texttt{\mdseries\slshape i}})]} (assuming \mbox{\texttt{\mdseries\slshape i}} is the number of the process, even in the default case) or for all weights if \texttt{\mbox{\texttt{\mdseries\slshape weight}} = 0}. See \texttt{PqTails} (\texttt{PqTails} (\ref{PqTails})) for more details. 

 \emph{Note:} For those familiar with the \texttt{pq} program, \texttt{PqComputeTails} uses menu item 7 of the Advanced $p$\texttt{\symbol{45}}Quotient menu. }

 

\subsection{\textcolor{Chapter }{PqAddTails}}
\logpage{[ 5, 7, 7 ]}\nobreak
\hyperdef{L}{X83CD6D888372DFB8}{}
{\noindent\textcolor{FuncColor}{$\triangleright$\enspace\texttt{PqAddTails({\mdseries\slshape i, weight})\index{PqAddTails@\texttt{PqAddTails}}
\label{PqAddTails}
}\hfill{\scriptsize (function)}}\\
\noindent\textcolor{FuncColor}{$\triangleright$\enspace\texttt{PqAddTails({\mdseries\slshape weight})\index{PqAddTails@\texttt{PqAddTails}!for default process}
\label{PqAddTails:for default process}
}\hfill{\scriptsize (function)}}\\


 for the \mbox{\texttt{\mdseries\slshape i}}th or default interactive \textsf{ANUPQ} process, direct the \texttt{pq} program to add the tails of weight \mbox{\texttt{\mdseries\slshape weight}}, previously computed by \texttt{PqComputeTails} (see{\nobreakspace}\texttt{PqComputeTails} (\ref{PqComputeTails})), if \mbox{\texttt{\mdseries\slshape weight}} is in the integer range \texttt{[2 .. PqPClass(\mbox{\texttt{\mdseries\slshape i}})]} (assuming \mbox{\texttt{\mdseries\slshape i}} is the number of the process, even in the default case) or for all weights if \texttt{\mbox{\texttt{\mdseries\slshape weight}} = 0}. See \texttt{PqTails} (\texttt{PqTails} (\ref{PqTails})) for more details. 

 \emph{Note:} For those familiar with the \texttt{pq} program, \texttt{PqAddTails} uses menu item 7 of the Advanced $p$\texttt{\symbol{45}}Quotient menu. }

 

\subsection{\textcolor{Chapter }{PqDoConsistencyChecks}}
\logpage{[ 5, 7, 8 ]}\nobreak
\hyperdef{L}{X82AA8FAE85826BB9}{}
{\noindent\textcolor{FuncColor}{$\triangleright$\enspace\texttt{PqDoConsistencyChecks({\mdseries\slshape i, weight, type})\index{PqDoConsistencyChecks@\texttt{PqDoConsistencyChecks}}
\label{PqDoConsistencyChecks}
}\hfill{\scriptsize (function)}}\\
\noindent\textcolor{FuncColor}{$\triangleright$\enspace\texttt{PqDoConsistencyChecks({\mdseries\slshape weight, type})\index{PqDoConsistencyChecks@\texttt{PqDoConsistencyChecks}!for default process}
\label{PqDoConsistencyChecks:for default process}
}\hfill{\scriptsize (function)}}\\


 for the \mbox{\texttt{\mdseries\slshape i}}th or default interactive \textsf{ANUPQ} process, do consistency checks for weight \mbox{\texttt{\mdseries\slshape weight}} if \mbox{\texttt{\mdseries\slshape weight}} is in the integer range \texttt{[3 .. PqPClass(\mbox{\texttt{\mdseries\slshape i}})]} (assuming \mbox{\texttt{\mdseries\slshape i}} is the number of the process) or for all weights if \texttt{\mbox{\texttt{\mdseries\slshape weight}} = 0}, and for type \mbox{\texttt{\mdseries\slshape type}} if \mbox{\texttt{\mdseries\slshape type}} is in the range \texttt{[1, 2, 3]} (see below) or for all types if \texttt{\mbox{\texttt{\mdseries\slshape type}} = 0}. (For its use in the context of finding the next class see \texttt{PqNextClass} (\ref{PqNextClass}).) 

 The \emph{type} of a consistency check is defined as follows. \texttt{PqDoConsistencyChecks(\mbox{\texttt{\mdseries\slshape i}}, \mbox{\texttt{\mdseries\slshape weight}}, \mbox{\texttt{\mdseries\slshape type}})} for \mbox{\texttt{\mdseries\slshape weight}} in \texttt{[3 .. PqPClass(\mbox{\texttt{\mdseries\slshape i}})]} and the given value of \mbox{\texttt{\mdseries\slshape type}} invokes the equivalent of the following \texttt{PqDoConsistencyCheck} calls (see{\nobreakspace}\texttt{PqDoConsistencyCheck} (\ref{PqDoConsistencyCheck})): 
\begin{description}
\item[{\texttt{\mbox{\texttt{\mdseries\slshape type}} = 1}:}]  \texttt{PqDoConsistencyCheck(\mbox{\texttt{\mdseries\slshape i}}, \mbox{\texttt{\mdseries\slshape a}}, \mbox{\texttt{\mdseries\slshape a}}, \mbox{\texttt{\mdseries\slshape a}})} checks \texttt{2 * PqWeight(\mbox{\texttt{\mdseries\slshape i}}, \mbox{\texttt{\mdseries\slshape a}}) + 1 = \mbox{\texttt{\mdseries\slshape weight}}}, for pc generators of index \mbox{\texttt{\mdseries\slshape a}}. 
\item[{\texttt{\mbox{\texttt{\mdseries\slshape type}} = 2}:}]  \texttt{PqDoConsistencyCheck(\mbox{\texttt{\mdseries\slshape i}}, \mbox{\texttt{\mdseries\slshape b}}, \mbox{\texttt{\mdseries\slshape b}}, \mbox{\texttt{\mdseries\slshape a}})} checks for pc generators of indices \mbox{\texttt{\mdseries\slshape b}}, \mbox{\texttt{\mdseries\slshape a}} satisfyingx both \texttt{\mbox{\texttt{\mdseries\slshape b}} {\textgreater} \mbox{\texttt{\mdseries\slshape a}}} and \texttt{PqWeight(\mbox{\texttt{\mdseries\slshape i}}, \mbox{\texttt{\mdseries\slshape b}}) + PqWeight(\mbox{\texttt{\mdseries\slshape i}}, \mbox{\texttt{\mdseries\slshape a}}) + 1 = \mbox{\texttt{\mdseries\slshape weight}}}. 
\item[{\texttt{\mbox{\texttt{\mdseries\slshape type}} = 3}:}]  \texttt{PqDoConsistencyCheck(\mbox{\texttt{\mdseries\slshape i}}, \mbox{\texttt{\mdseries\slshape c}}, \mbox{\texttt{\mdseries\slshape b}}, \mbox{\texttt{\mdseries\slshape a}})} checks for pc generators of indices \mbox{\texttt{\mdseries\slshape c}}, \mbox{\texttt{\mdseries\slshape b}}, \mbox{\texttt{\mdseries\slshape a}} satisfying \texttt{\mbox{\texttt{\mdseries\slshape c}} {\textgreater} \mbox{\texttt{\mdseries\slshape b}} {\textgreater} \mbox{\texttt{\mdseries\slshape a}}} and the sum of the weights of these generators equals \mbox{\texttt{\mdseries\slshape weight}}. 
\end{description}
 \emph{Notes} 

 \texttt{PqWeight(\mbox{\texttt{\mdseries\slshape i}}, \mbox{\texttt{\mdseries\slshape j}})} returns the weight of the \mbox{\texttt{\mdseries\slshape j}}th pc generator, for process \mbox{\texttt{\mdseries\slshape i}} (see{\nobreakspace}\texttt{PqWeight} (\ref{PqWeight})). 

 It is assumed that tails for the given weight (or weights) have already been
added (see{\nobreakspace}\texttt{PqTails} (\ref{PqTails})). 

 For those familiar with the \texttt{pq} program, \texttt{PqDoConsistencyChecks} performs menu item 8 of the Advanced $p$\texttt{\symbol{45}}Quotient menu. }

 

\subsection{\textcolor{Chapter }{PqCollectDefiningRelations}}
\logpage{[ 5, 7, 9 ]}\nobreak
\hyperdef{L}{X7A01EE0382689928}{}
{\noindent\textcolor{FuncColor}{$\triangleright$\enspace\texttt{PqCollectDefiningRelations({\mdseries\slshape i})\index{PqCollectDefiningRelations@\texttt{PqCollectDefiningRelations}}
\label{PqCollectDefiningRelations}
}\hfill{\scriptsize (function)}}\\
\noindent\textcolor{FuncColor}{$\triangleright$\enspace\texttt{PqCollectDefiningRelations({\mdseries\slshape })\index{PqCollectDefiningRelations@\texttt{PqCollectDefiningRelations}!for default process}
\label{PqCollectDefiningRelations:for default process}
}\hfill{\scriptsize (function)}}\\


 for the \mbox{\texttt{\mdseries\slshape i}}th or default interactive \textsf{ANUPQ} process, direct the \texttt{pq} program to collect the images of the defining relations of the original fp
group of the process, with respect to the current pc presentation, in the
context of finding the next class (see{\nobreakspace}\texttt{PqNextClass} (\ref{PqNextClass})). If the tails operation is not complete then the relations may be evaluated
incorrectly. 

 \emph{Note:} For those familiar with the \texttt{pq} program, \texttt{PqCollectDefiningRelations} performs menu item 9 of the Advanced $p$\texttt{\symbol{45}}Quotient menu. }

 

\subsection{\textcolor{Chapter }{PqCollectWordInDefiningGenerators}}
\logpage{[ 5, 7, 10 ]}\nobreak
\hyperdef{L}{X7C2B4C1E7BEB19D5}{}
{\noindent\textcolor{FuncColor}{$\triangleright$\enspace\texttt{PqCollectWordInDefiningGenerators({\mdseries\slshape i, word})\index{PqCollectWordInDefiningGenerators@\texttt{PqCollectWordInDefiningGenerators}}
\label{PqCollectWordInDefiningGenerators}
}\hfill{\scriptsize (function)}}\\
\noindent\textcolor{FuncColor}{$\triangleright$\enspace\texttt{PqCollectWordInDefiningGenerators({\mdseries\slshape word})\index{PqCollectWordInDefiningGenerators@\texttt{PqCollectWordInDefiningGenerators}!for default process}
\label{PqCollectWordInDefiningGenerators:for default process}
}\hfill{\scriptsize (function)}}\\


 for the \mbox{\texttt{\mdseries\slshape i}}th or default interactive \textsf{ANUPQ} process, take a user\texttt{\symbol{45}}defined word \mbox{\texttt{\mdseries\slshape word}} in the defining generators of the original presentation of the fp or pc group
of the process. Each generator is mapped into the current pc presentation, and
the resulting word is collected with respect to the current pc presentation.
The result of the collection is returned as a list of generator number,
exponent pairs. 

 The \mbox{\texttt{\mdseries\slshape word}} argument may be input in either of the two ways described for \texttt{PqCollect} (see{\nobreakspace}\texttt{PqCollect} (\ref{PqCollect})). 

 \emph{Note:} For those familiar with the \texttt{pq} program, \texttt{PqCollectDefiningGenerators} performs menu item 23 of the Advanced $p$\texttt{\symbol{45}}Quotient menu. }

 

\subsection{\textcolor{Chapter }{PqCommutatorDefiningGenerators}}
\logpage{[ 5, 7, 11 ]}\nobreak
\hyperdef{L}{X7DB05AA587D7A052}{}
{\noindent\textcolor{FuncColor}{$\triangleright$\enspace\texttt{PqCommutatorDefiningGenerators({\mdseries\slshape i, words, pow})\index{PqCommutatorDefiningGenerators@\texttt{PqCommutatorDefiningGenerators}}
\label{PqCommutatorDefiningGenerators}
}\hfill{\scriptsize (function)}}\\
\noindent\textcolor{FuncColor}{$\triangleright$\enspace\texttt{PqCommutatorDefiningGenerators({\mdseries\slshape words, pow})\index{PqCommutatorDefiningGenerators@\texttt{PqCommutatorDefiningGenerators}!for default process}
\label{PqCommutatorDefiningGenerators:for default process}
}\hfill{\scriptsize (function)}}\\


 for the \mbox{\texttt{\mdseries\slshape i}}th or default interactive \textsf{ANUPQ} process, take a list \mbox{\texttt{\mdseries\slshape words}} of user\texttt{\symbol{45}}defined words in the defining generators of the
original presentation of the fp or pc group of the process, and an integer
power \mbox{\texttt{\mdseries\slshape pow}}. Each generator is mapped into the current pc presentation. The list \mbox{\texttt{\mdseries\slshape words}} is interpreted as a left\texttt{\symbol{45}}normed commutator which is then
raised to \mbox{\texttt{\mdseries\slshape pow}} and collected with respect to the current pc presentation. The result of the
collection is returned as a list of generator number, exponent pairs. 

 \emph{Note} For those familiar with the \texttt{pq} program, \texttt{PqCommutatorDefiningGenerators} performs menu item 24 of the Advanced $p$\texttt{\symbol{45}}Quotient menu. }

 

\subsection{\textcolor{Chapter }{PqDoExponentChecks}}
\logpage{[ 5, 7, 12 ]}\nobreak
\hyperdef{L}{X80E563A97EDE083E}{}
{\noindent\textcolor{FuncColor}{$\triangleright$\enspace\texttt{PqDoExponentChecks({\mdseries\slshape i: [Bounds := list]})\index{PqDoExponentChecks@\texttt{PqDoExponentChecks}}
\label{PqDoExponentChecks}
}\hfill{\scriptsize (function)}}\\
\noindent\textcolor{FuncColor}{$\triangleright$\enspace\texttt{PqDoExponentChecks({\mdseries\slshape : [Bounds := list]})\index{PqDoExponentChecks@\texttt{PqDoExponentChecks}!for default process}
\label{PqDoExponentChecks:for default process}
}\hfill{\scriptsize (function)}}\\


 for the \mbox{\texttt{\mdseries\slshape i}}th or default interactive \textsf{ANUPQ} process, direct the \texttt{pq} program to do exponent checks for weights (inclusively) between the bounds of \texttt{Bounds} or for all weights if \texttt{Bounds} is not given. The value \mbox{\texttt{\mdseries\slshape list}} of \texttt{Bounds} (assuming the interactive process is numbered \mbox{\texttt{\mdseries\slshape i}}) should be a list of two integers \mbox{\texttt{\mdseries\slshape low}}, \mbox{\texttt{\mdseries\slshape high}} satisfying  $1 \le \mbox{\texttt{\mdseries\slshape low}} \le \mbox{\texttt{\mdseries\slshape high}} \le $ \texttt{PqPClass(\mbox{\texttt{\mdseries\slshape i}})} (see{\nobreakspace}\texttt{PqPClass} (\ref{PqPClass})). If no exponent law has been specified, no exponent checks are performed. 

 \emph{Note:} For those familiar with the \texttt{pq} program, \texttt{PqDoExponentChecks} performs menu item 10 of the Advanced $p$\texttt{\symbol{45}}Quotient menu. }

 

\subsection{\textcolor{Chapter }{PqEliminateRedundantGenerators}}
\logpage{[ 5, 7, 13 ]}\nobreak
\hyperdef{L}{X7BBAB2787B305516}{}
{\noindent\textcolor{FuncColor}{$\triangleright$\enspace\texttt{PqEliminateRedundantGenerators({\mdseries\slshape i})\index{PqEliminateRedundantGenerators@\texttt{PqEliminateRedundantGenerators}}
\label{PqEliminateRedundantGenerators}
}\hfill{\scriptsize (function)}}\\
\noindent\textcolor{FuncColor}{$\triangleright$\enspace\texttt{PqEliminateRedundantGenerators({\mdseries\slshape })\index{PqEliminateRedundantGenerators@\texttt{PqEliminateRedundantGenerators}!for default process}
\label{PqEliminateRedundantGenerators:for default process}
}\hfill{\scriptsize (function)}}\\


 for the \mbox{\texttt{\mdseries\slshape i}}th or default interactive \textsf{ANUPQ} process, direct the \texttt{pq} program to eliminate redundant generators of the current $p$\texttt{\symbol{45}}quotient. 

 \emph{Note:} For those familiar with the \texttt{pq} program, \texttt{PqEliminateRedundantGenerators} performs menu item 11 of the Advanced $p$\texttt{\symbol{45}}Quotient menu. }

 

\subsection{\textcolor{Chapter }{PqRevertToPreviousClass}}
\logpage{[ 5, 7, 14 ]}\nobreak
\hyperdef{L}{X783E53B5853F45E6}{}
{\noindent\textcolor{FuncColor}{$\triangleright$\enspace\texttt{PqRevertToPreviousClass({\mdseries\slshape i})\index{PqRevertToPreviousClass@\texttt{PqRevertToPreviousClass}}
\label{PqRevertToPreviousClass}
}\hfill{\scriptsize (function)}}\\
\noindent\textcolor{FuncColor}{$\triangleright$\enspace\texttt{PqRevertToPreviousClass({\mdseries\slshape })\index{PqRevertToPreviousClass@\texttt{PqRevertToPreviousClass}!for default process}
\label{PqRevertToPreviousClass:for default process}
}\hfill{\scriptsize (function)}}\\


 for the \mbox{\texttt{\mdseries\slshape i}}th or default interactive \textsf{ANUPQ} process, direct the \texttt{pq} program to abandon the current class and revert to the previous class. 

 \emph{Note:} For those familiar with the \texttt{pq} program, \texttt{PqRevertToPreviousClass} performs menu item 12 of the Advanced $p$\texttt{\symbol{45}}Quotient menu. }

 

\subsection{\textcolor{Chapter }{PqSetMaximalOccurrences}}
\logpage{[ 5, 7, 15 ]}\nobreak
\hyperdef{L}{X8265A6DB81CD24DB}{}
{\noindent\textcolor{FuncColor}{$\triangleright$\enspace\texttt{PqSetMaximalOccurrences({\mdseries\slshape i, noccur})\index{PqSetMaximalOccurrences@\texttt{PqSetMaximalOccurrences}}
\label{PqSetMaximalOccurrences}
}\hfill{\scriptsize (function)}}\\
\noindent\textcolor{FuncColor}{$\triangleright$\enspace\texttt{PqSetMaximalOccurrences({\mdseries\slshape noccur})\index{PqSetMaximalOccurrences@\texttt{PqSetMaximalOccurrences}!for default process}
\label{PqSetMaximalOccurrences:for default process}
}\hfill{\scriptsize (function)}}\\


 for the \mbox{\texttt{\mdseries\slshape i}}th or default interactive \textsf{ANUPQ} process, direct the \texttt{pq} program to set maximal occurrences of the weight 1 generators in the
definitions of pcp generators of the group of the process. This can be used to
avoid the definition of generators of which one knows for theoretical reasons
that they would be eliminated later on. 

 The argument \mbox{\texttt{\mdseries\slshape noccur}} must be a list of non\texttt{\symbol{45}}negative integers of length the
number of weight 1 generators (i.e.{\nobreakspace}the rank of the class 1 $p$\texttt{\symbol{45}}quotient of the group of the process). An entry of \texttt{0} for a particular generator indicates that there is no limit on the number of
occurrences for the generator. 

 \emph{Note:} For those familiar with the \texttt{pq} program, \texttt{PqSetMaximalOccurrences} performs menu item 13 of the Advanced $p$\texttt{\symbol{45}}Quotient menu. }

 

\subsection{\textcolor{Chapter }{PqSetMetabelian}}
\logpage{[ 5, 7, 16 ]}\nobreak
\hyperdef{L}{X87A35ABB7E11595E}{}
{\noindent\textcolor{FuncColor}{$\triangleright$\enspace\texttt{PqSetMetabelian({\mdseries\slshape i})\index{PqSetMetabelian@\texttt{PqSetMetabelian}}
\label{PqSetMetabelian}
}\hfill{\scriptsize (function)}}\\
\noindent\textcolor{FuncColor}{$\triangleright$\enspace\texttt{PqSetMetabelian({\mdseries\slshape })\index{PqSetMetabelian@\texttt{PqSetMetabelian}!for default process}
\label{PqSetMetabelian:for default process}
}\hfill{\scriptsize (function)}}\\


 for the \mbox{\texttt{\mdseries\slshape i}}th or default interactive \textsf{ANUPQ} process, direct the \texttt{pq} program to enforce metabelian\texttt{\symbol{45}}ness. 

 \emph{Note:} For those familiar with the \texttt{pq} program, \texttt{PqSetMetabelian} performs menu item 14 of the Advanced $p$\texttt{\symbol{45}}Quotient menu. }

 

\subsection{\textcolor{Chapter }{PqDoConsistencyCheck}}
\logpage{[ 5, 7, 17 ]}\nobreak
\hyperdef{L}{X7C7D17F878AA1BAA}{}
{\noindent\textcolor{FuncColor}{$\triangleright$\enspace\texttt{PqDoConsistencyCheck({\mdseries\slshape i, c, b, a})\index{PqDoConsistencyCheck@\texttt{PqDoConsistencyCheck}}
\label{PqDoConsistencyCheck}
}\hfill{\scriptsize (function)}}\\
\noindent\textcolor{FuncColor}{$\triangleright$\enspace\texttt{PqDoConsistencyCheck({\mdseries\slshape c, b, a})\index{PqDoConsistencyCheck@\texttt{PqDoConsistencyCheck}!for default process}
\label{PqDoConsistencyCheck:for default process}
}\hfill{\scriptsize (function)}}\\
\noindent\textcolor{FuncColor}{$\triangleright$\enspace\texttt{PqJacobi({\mdseries\slshape i, c, b, a})\index{PqJacobi@\texttt{PqJacobi}}
\label{PqJacobi}
}\hfill{\scriptsize (function)}}\\
\noindent\textcolor{FuncColor}{$\triangleright$\enspace\texttt{PqJacobi({\mdseries\slshape c, b, a})\index{PqJacobi@\texttt{PqJacobi}!for default process}
\label{PqJacobi:for default process}
}\hfill{\scriptsize (function)}}\\


 for the \mbox{\texttt{\mdseries\slshape i}}th or default interactive \textsf{ANUPQ} process, direct the \texttt{pq} program to do the consistency check for the pc generators with indices \mbox{\texttt{\mdseries\slshape c}}, \mbox{\texttt{\mdseries\slshape b}}, \mbox{\texttt{\mdseries\slshape a}} which should be non\texttt{\symbol{45}}increasing positive integers,
i.e.{\nobreakspace}$\mbox{\texttt{\mdseries\slshape c}} \ge \mbox{\texttt{\mdseries\slshape b}} \ge \mbox{\texttt{\mdseries\slshape a}}$. 

 There are 3 types of consistency checks: 
\[ \begin{array}{rclrl} (a^n)a &=& a(a^n) && {\rm (Type\ 1)} \\ (b^n)a &=&
b^{(n-1)}(ba), b(a^n) = (ba)a^{(n-1)} && {\rm (Type\ 2)} \\ c(ba) &=& (cb)a &&
{\rm (Type\ 3)} \\ \end{array} \]
 The reason some people talk about Jacobi relations instead of consistency
checks becomes clear when one looks at the consistency check of type 3: 
\[ \begin{array}{rcl} c(ba) &=& a c[c,a] b[b,a] = acb [c,a][c,a,b][b,a] = \dots
\\ (cb)a &=& b c[c,b] a = a b[b,a] c[c,a] [c,b][c,b,a] \\ &=& abc [b,a]
[b,a,c] [c,a] [c,b] [c,b,a] = \dots \\ \end{array} \]
 Each collection would normally carry on further. But one can see already that
no other commutators of weight 3 will occur. After all terms of weight one and
weight two have been moved to the left we end up with: 
\[ \begin{array}{rcl} & &abc [b,a] [c,a] [c,b] [c,a,b] \dots \\ &=&abc [b,a]
[c,a] [c,b] [c,b,a] [b,a,c] \dots \\ \end{array} \]
 Modulo terms of weight 4 this is equivalent to 
\[ [c,a,b] [b,c,a] [a,b,c] = 1 \]
 which is the Jacobi identity. 

 See also \texttt{PqDoConsistencyChecks} (\texttt{PqDoConsistencyChecks} (\ref{PqDoConsistencyChecks})). 

 \emph{Note:} For those familiar with the \texttt{pq} program, \texttt{PqDoConsistencyCheck} and \texttt{PqJacobi} perform menu item 15 of the Advanced $p$\texttt{\symbol{45}}Quotient menu. }

 

\subsection{\textcolor{Chapter }{PqCompact}}
\logpage{[ 5, 7, 18 ]}\nobreak
\hyperdef{L}{X843953F48319FDE1}{}
{\noindent\textcolor{FuncColor}{$\triangleright$\enspace\texttt{PqCompact({\mdseries\slshape i})\index{PqCompact@\texttt{PqCompact}}
\label{PqCompact}
}\hfill{\scriptsize (function)}}\\
\noindent\textcolor{FuncColor}{$\triangleright$\enspace\texttt{PqCompact({\mdseries\slshape })\index{PqCompact@\texttt{PqCompact}!for default process}
\label{PqCompact:for default process}
}\hfill{\scriptsize (function)}}\\


 for the \mbox{\texttt{\mdseries\slshape i}}th or default interactive \textsf{ANUPQ} process, direct the \texttt{pq} program to do a compaction of its work space. This function is safe to perform
only at certain points in time. 

 \emph{Note:} For those familiar with the \texttt{pq} program, \texttt{PqCompact} performs menu item 16 of the Advanced $p$\texttt{\symbol{45}}Quotient menu. }

 

\subsection{\textcolor{Chapter }{PqEchelonise}}
\logpage{[ 5, 7, 19 ]}\nobreak
\hyperdef{L}{X80344EEC78A09ACF}{}
{\noindent\textcolor{FuncColor}{$\triangleright$\enspace\texttt{PqEchelonise({\mdseries\slshape i})\index{PqEchelonise@\texttt{PqEchelonise}}
\label{PqEchelonise}
}\hfill{\scriptsize (function)}}\\
\noindent\textcolor{FuncColor}{$\triangleright$\enspace\texttt{PqEchelonise({\mdseries\slshape })\index{PqEchelonise@\texttt{PqEchelonise}!for default process}
\label{PqEchelonise:for default process}
}\hfill{\scriptsize (function)}}\\


 for the \mbox{\texttt{\mdseries\slshape i}}th or default interactive \textsf{ANUPQ} process, direct the \texttt{pq} program to echelonise the word most recently collected by \texttt{PqCollect} or \texttt{PqCommutator} against the relations of the current pc presentation, and return the number of
the generator made redundant or \texttt{fail} if no generator was made redundant. A call to \texttt{PqCollect} (see{\nobreakspace}\texttt{PqCollect} (\ref{PqCollect})) or \texttt{PqCommutator} (see{\nobreakspace}\texttt{PqCommutator} (\ref{PqCommutator})) needs to be performed prior to using this command. 

 \emph{Note:} For those familiar with the \texttt{pq} program, \texttt{PqEchelonise} performs menu item 17 of the Advanced $p$\texttt{\symbol{45}}Quotient menu. }

 

\subsection{\textcolor{Chapter }{PqSupplyAutomorphisms}}
\logpage{[ 5, 7, 20 ]}\nobreak
\hyperdef{L}{X852947327FC39DDF}{}
{\noindent\textcolor{FuncColor}{$\triangleright$\enspace\texttt{PqSupplyAutomorphisms({\mdseries\slshape i, mlist})\index{PqSupplyAutomorphisms@\texttt{PqSupplyAutomorphisms}}
\label{PqSupplyAutomorphisms}
}\hfill{\scriptsize (function)}}\\
\noindent\textcolor{FuncColor}{$\triangleright$\enspace\texttt{PqSupplyAutomorphisms({\mdseries\slshape mlist})\index{PqSupplyAutomorphisms@\texttt{PqSupplyAutomorphisms}!for default process}
\label{PqSupplyAutomorphisms:for default process}
}\hfill{\scriptsize (function)}}\\


 for the \mbox{\texttt{\mdseries\slshape i}}th or default interactive \textsf{ANUPQ} process, supply the automorphism data provided by the list \mbox{\texttt{\mdseries\slshape mlist}} of matrices with non\texttt{\symbol{45}}negative integer coefficients. Each
matrix in \mbox{\texttt{\mdseries\slshape mlist}} describes one automorphism in the following way. 
\begin{itemize}
\item  The rows of each matrix correspond to the pc generators of weight one. 
\item  Each row is the exponent vector of the image of the corresponding weight one
generator under the respective automorphism. 
\end{itemize}
 \emph{Note:} For those familiar with the \texttt{pq} program, \texttt{PqSupplyAutomorphisms} uses menu item 18 of the Advanced $p$\texttt{\symbol{45}}Quotient menu. }

 

\subsection{\textcolor{Chapter }{PqExtendAutomorphisms}}
\logpage{[ 5, 7, 21 ]}\nobreak
\hyperdef{L}{X7B919B537C042DAD}{}
{\noindent\textcolor{FuncColor}{$\triangleright$\enspace\texttt{PqExtendAutomorphisms({\mdseries\slshape i})\index{PqExtendAutomorphisms@\texttt{PqExtendAutomorphisms}}
\label{PqExtendAutomorphisms}
}\hfill{\scriptsize (function)}}\\
\noindent\textcolor{FuncColor}{$\triangleright$\enspace\texttt{PqExtendAutomorphisms({\mdseries\slshape })\index{PqExtendAutomorphisms@\texttt{PqExtendAutomorphisms}!for default process}
\label{PqExtendAutomorphisms:for default process}
}\hfill{\scriptsize (function)}}\\


 for the \mbox{\texttt{\mdseries\slshape i}}th or default interactive \textsf{ANUPQ} process, direct the \texttt{pq} program to extend automorphisms of the $p$\texttt{\symbol{45}}quotient of the previous class to the $p$\texttt{\symbol{45}}quotient of the present class. 

 \emph{Note:} For those familiar with the \texttt{pq} program, \texttt{PqExtendAutomorphisms} uses menu item 18 of the Advanced $p$\texttt{\symbol{45}}Quotient menu. }

 

\subsection{\textcolor{Chapter }{PqApplyAutomorphisms}}
\logpage{[ 5, 7, 22 ]}\nobreak
\hyperdef{L}{X7CFD54657DE4BD39}{}
{\noindent\textcolor{FuncColor}{$\triangleright$\enspace\texttt{PqApplyAutomorphisms({\mdseries\slshape i, qfac})\index{PqApplyAutomorphisms@\texttt{PqApplyAutomorphisms}}
\label{PqApplyAutomorphisms}
}\hfill{\scriptsize (function)}}\\
\noindent\textcolor{FuncColor}{$\triangleright$\enspace\texttt{PqApplyAutomorphisms({\mdseries\slshape qfac})\index{PqApplyAutomorphisms@\texttt{PqApplyAutomorphisms}!for default process}
\label{PqApplyAutomorphisms:for default process}
}\hfill{\scriptsize (function)}}\\


 for the \mbox{\texttt{\mdseries\slshape i}}th or default interactive \textsf{ANUPQ} process, direct the \texttt{pq} program to apply automorphisms; \mbox{\texttt{\mdseries\slshape qfac}} is the queue factor e.g. \texttt{15}. 

 \emph{Note:} For those familiar with the \texttt{pq} program, \texttt{PqCloseRelations} performs menu item 19 of the Advanced $p$\texttt{\symbol{45}}Quotient menu. }

 

\subsection{\textcolor{Chapter }{PqDisplayStructure}}
\logpage{[ 5, 7, 23 ]}\nobreak
\hyperdef{L}{X81A3D5BF87A34934}{}
{\noindent\textcolor{FuncColor}{$\triangleright$\enspace\texttt{PqDisplayStructure({\mdseries\slshape i: [Bounds := list]})\index{PqDisplayStructure@\texttt{PqDisplayStructure}}
\label{PqDisplayStructure}
}\hfill{\scriptsize (function)}}\\
\noindent\textcolor{FuncColor}{$\triangleright$\enspace\texttt{PqDisplayStructure({\mdseries\slshape : [Bounds := list]})\index{PqDisplayStructure@\texttt{PqDisplayStructure}!for default process}
\label{PqDisplayStructure:for default process}
}\hfill{\scriptsize (function)}}\\


 for the \mbox{\texttt{\mdseries\slshape i}}th or default interactive \textsf{ANUPQ} process, direct the \texttt{pq} program to display the structure for the pcp generators numbered (inclusively)
between the bounds of \texttt{Bounds} or for all generators if \texttt{Bounds} is not given. The value \mbox{\texttt{\mdseries\slshape list}} of \texttt{Bounds} (assuming the interactive process is numbered \mbox{\texttt{\mdseries\slshape i}}) should be a list of two integers  \mbox{\texttt{\mdseries\slshape low}}, \mbox{\texttt{\mdseries\slshape high}} satisfying $1 \le \mbox{\texttt{\mdseries\slshape low}} \le \mbox{\texttt{\mdseries\slshape high}} \le $ \texttt{PqNrPcGenerators(\mbox{\texttt{\mdseries\slshape i}})} (see{\nobreakspace}\texttt{PqNrPcGenerators} (\ref{PqNrPcGenerators})). \texttt{PqDisplayStructure} also accepts the option \texttt{OutputLevel} (see \ref{option OutputLevel}). 

 \emph{Explanation of output} 

 New generators are defined as commutators of previous generators and
generators of class 1 or as $p$\texttt{\symbol{45}}th powers of generators that have themselves been defined
as $p$\texttt{\symbol{45}}th powers. A generator is never defined as $p$\texttt{\symbol{45}}th power of a commutator. 

 Therefore, there are two cases: all the numbers on the righthand side are
either the same or they differ. Below, \texttt{g\mbox{\texttt{\mdseries\slshape i}}} refers to the \mbox{\texttt{\mdseries\slshape i}}th defining generator. 
\begin{itemize}
\item  If the righthand side numbers are all the same, then the generator is a $p$\texttt{\symbol{45}}th power (of a $p$\texttt{\symbol{45}}th power of a $p$\texttt{\symbol{45}}th power, etc.). The number of repeated digits say how
often a $p$\texttt{\symbol{45}}th power has to be taken. 

 In the following example, the generator number 31 is the eleventh power of
generator 17 which in turn is an eleventh power and so on: 

 \texttt{\symbol{92}}begintt \#I 31 is defined on 17\texttt{\symbol{94}}11 = 1
1 1 1 1 \texttt{\symbol{92}}endtt So generator 31 is obtained by taking the
eleventh power of generator 1 five times. 
\item  If the numbers are not all the same, the generator is defined by a commutator.
If the first two generator numbers differ, the generator is defined as a
left\texttt{\symbol{45}}normed commutator of the weight one generators, e.g. 

 \texttt{\symbol{92}}begintt \#I 19 is defined on [11, 1] = 2 1 1 1 1
\texttt{\symbol{92}}endtt Here, generator 19 is defined as the commutator of
generator 11 and generator 1 which is the same as the
left\texttt{\symbol{45}}normed commutator \texttt{[x2, x1, x1, x1, x1]}. One can check this by tracing back the definition of generator 11 until one
gets to a generator of class 1. 
\item  If the first two generator numbers are identical, then the left most component
of the left\texttt{\symbol{45}}normed commutator is a $p$\texttt{\symbol{45}}th power, e.g. 

 \texttt{\symbol{92}}begintt \#I 25 is defined on [14, 1] = 1 1 2 1 1
\texttt{\symbol{92}}endtt 

 In this example, generator 25 is defined as commutator of generator 14 and
generator 1. The left\texttt{\symbol{45}}normed commutator is 
\[ [(x1^{11})^{11}, x2, x1, x1] \]
 Again, this can be verified by tracing back the definitions. 
\end{itemize}
 \emph{Note:} For those familiar with the \texttt{pq} program, \texttt{PqDisplayStructure} performs menu item 20 of the Advanced $p$\texttt{\symbol{45}}Quotient menu. }

 

\subsection{\textcolor{Chapter }{PqDisplayAutomorphisms}}
\logpage{[ 5, 7, 24 ]}\nobreak
\hyperdef{L}{X853834B478C4EEE2}{}
{\noindent\textcolor{FuncColor}{$\triangleright$\enspace\texttt{PqDisplayAutomorphisms({\mdseries\slshape i: [Bounds := list]})\index{PqDisplayAutomorphisms@\texttt{PqDisplayAutomorphisms}}
\label{PqDisplayAutomorphisms}
}\hfill{\scriptsize (function)}}\\
\noindent\textcolor{FuncColor}{$\triangleright$\enspace\texttt{PqDisplayAutomorphisms({\mdseries\slshape : [Bounds := list]})\index{PqDisplayAutomorphisms@\texttt{PqDisplayAutomorphisms}!for default process}
\label{PqDisplayAutomorphisms:for default process}
}\hfill{\scriptsize (function)}}\\


 for the \mbox{\texttt{\mdseries\slshape i}}th or default interactive \textsf{ANUPQ} process, direct the \texttt{pq} program to display the automorphism actions on the pcp generators numbered
(inclusively) between the bounds of \texttt{Bounds} or for all generators if \texttt{Bounds} is not given. The value \mbox{\texttt{\mdseries\slshape list}} of \texttt{Bounds} (assuming the interactive process is numbered \mbox{\texttt{\mdseries\slshape i}}) should be a list of two integers  \mbox{\texttt{\mdseries\slshape low}}, \mbox{\texttt{\mdseries\slshape high}} satisfying $1 \le \mbox{\texttt{\mdseries\slshape low}} \le \mbox{\texttt{\mdseries\slshape high}} \le $ \texttt{PqNrPcGenerators(\mbox{\texttt{\mdseries\slshape i}})} (see{\nobreakspace}\texttt{PqNrPcGenerators} (\ref{PqNrPcGenerators})). \texttt{PqDisplayStructure} also accepts the option \texttt{OutputLevel} (see \ref{option OutputLevel}). 

 \emph{Note:} For those familiar with the \texttt{pq} program, \texttt{PqDisplayAutomorphisms} performs menu item 21 of the Advanced $p$\texttt{\symbol{45}}Quotient menu. }

 

\subsection{\textcolor{Chapter }{PqWritePcPresentation}}
\logpage{[ 5, 7, 25 ]}\nobreak
\hyperdef{L}{X80687A138626EB2D}{}
{\noindent\textcolor{FuncColor}{$\triangleright$\enspace\texttt{PqWritePcPresentation({\mdseries\slshape i, filename})\index{PqWritePcPresentation@\texttt{PqWritePcPresentation}}
\label{PqWritePcPresentation}
}\hfill{\scriptsize (function)}}\\
\noindent\textcolor{FuncColor}{$\triangleright$\enspace\texttt{PqWritePcPresentation({\mdseries\slshape filename})\index{PqWritePcPresentation@\texttt{PqWritePcPresentation}!for default process}
\label{PqWritePcPresentation:for default process}
}\hfill{\scriptsize (function)}}\\


 for the \mbox{\texttt{\mdseries\slshape i}}th or default interactive \textsf{ANUPQ} process, direct the \texttt{pq} program to write a pc presentation of a previously\texttt{\symbol{45}}computed
quotient of the group of that process, to the file with name \mbox{\texttt{\mdseries\slshape filename}}. Here the group of a process is the one given as first argument when \texttt{PqStart} was called to initiate that process (for process \mbox{\texttt{\mdseries\slshape i}} the group is stored as \texttt{ANUPQData.io[\mbox{\texttt{\mdseries\slshape i}}].group}). If the first character of the string \mbox{\texttt{\mdseries\slshape filename}} is not \texttt{/}, \mbox{\texttt{\mdseries\slshape filename}} is assumed to be the path of a writable file relative to the directory in
which \textsf{GAP} was started. If a pc presentation has not been previously computed by the \texttt{pq} program, then \texttt{pq} is called to compute it first, effectively invoking \texttt{PqPcPresentation} (see{\nobreakspace}\texttt{PqPcPresentation} (\ref{PqPcPresentation})). 

 \emph{Note:} For those familiar with the \texttt{pq} program, \texttt{PqPcWritePresentation} performs menu item 25 of the Advanced $p$\texttt{\symbol{45}}Quotient menu. 

               }

 }

 
\section{\textcolor{Chapter }{Commands from the Standard Presentation menu}}\logpage{[ 5, 8, 0 ]}
\hyperdef{L}{X7B94EDA385FDD904}{}
{
 

\subsection{\textcolor{Chapter }{PqSPComputePcpAndPCover}}
\logpage{[ 5, 8, 1 ]}\nobreak
\hyperdef{L}{X78D1963D85C71B7B}{}
{\noindent\textcolor{FuncColor}{$\triangleright$\enspace\texttt{PqSPComputePcpAndPCover({\mdseries\slshape i: options})\index{PqSPComputePcpAndPCover@\texttt{PqSPComputePcpAndPCover}}
\label{PqSPComputePcpAndPCover}
}\hfill{\scriptsize (function)}}\\
\noindent\textcolor{FuncColor}{$\triangleright$\enspace\texttt{PqSPComputePcpAndPCover({\mdseries\slshape : options})\index{PqSPComputePcpAndPCover@\texttt{PqSPComputePcpAndPCover}!for default process}
\label{PqSPComputePcpAndPCover:for default process}
}\hfill{\scriptsize (function)}}\\


 for the \mbox{\texttt{\mdseries\slshape i}}th or default interactive \textsf{ANUPQ} process, directs the \texttt{pq} program to compute for the group of that process a pc presentation up to the $p$\texttt{\symbol{45}}quotient of maximum class or the value of the option \texttt{ClassBound} and the $p$\texttt{\symbol{45}}cover of that quotient, and sets up tabular information
required for computation of a standard presentation. Here the group of a
process is the one given as first argument when \texttt{PqStart} was called to initiate that process (for process \mbox{\texttt{\mdseries\slshape i}} the group is stored as \texttt{ANUPQData.io[\mbox{\texttt{\mdseries\slshape i}}].group}). 

 The possible \mbox{\texttt{\mdseries\slshape options}} are \texttt{Prime}, \texttt{ClassBound}, \texttt{Relators}, \texttt{Exponent}, \texttt{Metabelian} and \texttt{OutputLevel} (see Chapter{\nobreakspace}\hyperref[ANUPQ Options]{`ANUPQ Options'} for detailed descriptions of these options). The option \texttt{Prime} is normally determined via \texttt{PrimePGroup}, and so is not required unless the group doesn't know it's a $p$\texttt{\symbol{45}}group and \texttt{HasPrimePGroup} returns \texttt{false}. 

 \emph{Note:} For those familiar with the \texttt{pq} program, \texttt{PqSPComputePcpAndPCover} performs option 1 of the Standard Presentation menu. }

 

\subsection{\textcolor{Chapter }{PqSPStandardPresentation}}
\logpage{[ 5, 8, 2 ]}\nobreak
\hyperdef{L}{X876F30187E89E467}{}
{\noindent\textcolor{FuncColor}{$\triangleright$\enspace\texttt{PqSPStandardPresentation({\mdseries\slshape i[, mlist]: [options]})\index{PqSPStandardPresentation@\texttt{PqSPStandardPresentation}}
\label{PqSPStandardPresentation}
}\hfill{\scriptsize (function)}}\\
\noindent\textcolor{FuncColor}{$\triangleright$\enspace\texttt{PqSPStandardPresentation({\mdseries\slshape [mlist]: [options]})\index{PqSPStandardPresentation@\texttt{PqSPStandardPresentation}!for default process}
\label{PqSPStandardPresentation:for default process}
}\hfill{\scriptsize (function)}}\\


 for the \mbox{\texttt{\mdseries\slshape i}}th or default interactive \textsf{ANUPQ} process, inputs data given by \mbox{\texttt{\mdseries\slshape options}} to compute a standard presentation for the group of that process. If argument \mbox{\texttt{\mdseries\slshape mlist}} is given it is assumed to be the automorphism group data required. Otherwise
it is assumed that a call to either \texttt{Pq} (see{\nobreakspace}\texttt{Pq} (\ref{Pq:interactive})) or \texttt{PqEpimorphism} (see{\nobreakspace}\texttt{PqEpimorphism} (\ref{PqEpimorphism:interactive})) has generated a $p$\texttt{\symbol{45}}quotient and that \textsf{GAP} can compute its automorphism group from which the necessary automorphism group
data can be derived. The group of the process is the one given as first
argument when \texttt{PqStart} was called to initiate the process (for process \mbox{\texttt{\mdseries\slshape i}} the group is stored as \texttt{ANUPQData.io[\mbox{\texttt{\mdseries\slshape i}}].group} and the $p$\texttt{\symbol{45}}quotient if existent is stored as \texttt{ANUPQData.io[\mbox{\texttt{\mdseries\slshape i}}].pQuotient}). If \mbox{\texttt{\mdseries\slshape mlist}} is not given and a $p$\texttt{\symbol{45}}quotient of the group has not been previously computed a
class 1 $p$\texttt{\symbol{45}}quotient is computed. 

 \texttt{PqSPStandardPresentation} accepts three options, all optional: 
\begin{itemize}
\item  \texttt{ClassBound := \mbox{\texttt{\mdseries\slshape n}}}\index{option ClassBound} 
\item  \texttt{PcgsAutomorphisms}\index{option PcgsAutomorphisms} 
\item  \texttt{StandardPresentationFile := \mbox{\texttt{\mdseries\slshape filename}}}\index{option StandardPresentationFile} 
\end{itemize}
 If \texttt{ClassBound} is omitted it defaults to 63. 

 Detailed descriptions of the above options may be found in
Chapter{\nobreakspace}\hyperref[ANUPQ Options]{`ANUPQ Options'}. 

 \emph{Note:} For those familiar with the \texttt{pq} program, \texttt{PqSPPcPresentation} performs menu item 2 of the Standard Presentation menu. }

 

\subsection{\textcolor{Chapter }{PqSPSavePresentation}}
\logpage{[ 5, 8, 3 ]}\nobreak
\hyperdef{L}{X7E1B2B088322F48A}{}
{\noindent\textcolor{FuncColor}{$\triangleright$\enspace\texttt{PqSPSavePresentation({\mdseries\slshape i, filename})\index{PqSPSavePresentation@\texttt{PqSPSavePresentation}}
\label{PqSPSavePresentation}
}\hfill{\scriptsize (function)}}\\
\noindent\textcolor{FuncColor}{$\triangleright$\enspace\texttt{PqSPSavePresentation({\mdseries\slshape filename})\index{PqSPSavePresentation@\texttt{PqSPSavePresentation}!for default process}
\label{PqSPSavePresentation:for default process}
}\hfill{\scriptsize (function)}}\\


 for the \mbox{\texttt{\mdseries\slshape i}}th or default interactive \textsf{ANUPQ} process, directs the \texttt{pq} program to save the standard presentation previously computed for the group of
that process to the file with name \mbox{\texttt{\mdseries\slshape filename}}, where the group of a process is the one given as first argument when \texttt{PqStart} was called to initiate that process. If the first character of the string \mbox{\texttt{\mdseries\slshape filename}} is not \texttt{/}, \mbox{\texttt{\mdseries\slshape filename}} is assumed to be the path of a writable file relative to the directory in
which \textsf{GAP} was started. 

 \emph{Note:} For those familiar with the \texttt{pq} program, \texttt{PqSPSavePresentation} performs menu item 3 of the Standard Presentation menu. }

 

\subsection{\textcolor{Chapter }{PqSPCompareTwoFilePresentations}}
\logpage{[ 5, 8, 4 ]}\nobreak
\hyperdef{L}{X81D2F9FF7C241D1E}{}
{\noindent\textcolor{FuncColor}{$\triangleright$\enspace\texttt{PqSPCompareTwoFilePresentations({\mdseries\slshape i, f1, f2})\index{PqSPCompareTwoFilePresentations@\texttt{PqSPCompareTwoFilePresentations}}
\label{PqSPCompareTwoFilePresentations}
}\hfill{\scriptsize (function)}}\\
\noindent\textcolor{FuncColor}{$\triangleright$\enspace\texttt{PqSPCompareTwoFilePresentations({\mdseries\slshape f1, f2})\index{PqSPCompareTwoFilePresentations@\texttt{PqSPCompareTwoFilePresentations}!for default process}
\label{PqSPCompareTwoFilePresentations:for default process}
}\hfill{\scriptsize (function)}}\\


 for the \mbox{\texttt{\mdseries\slshape i}}th or default interactive \textsf{ANUPQ} process, direct the \texttt{pq} program to compare the presentations in the files with names \mbox{\texttt{\mdseries\slshape f1}} and \mbox{\texttt{\mdseries\slshape f2}} and returns \texttt{true} if they are identical and \texttt{false} otherwise. For each of the strings \mbox{\texttt{\mdseries\slshape f1}} and \mbox{\texttt{\mdseries\slshape f2}}, if the first character is not a \texttt{/} then it is assumed to be the path of a readable file relative to the directory
in which \textsf{GAP} was started. 

 \emph{Notes} 

 The presentations in files \mbox{\texttt{\mdseries\slshape f1}} and \mbox{\texttt{\mdseries\slshape f2}} must have been generated by the \texttt{pq} program but they do \emph{not} need to be \emph{standard} presentations. If If the presentations in files \mbox{\texttt{\mdseries\slshape f1}} and \mbox{\texttt{\mdseries\slshape f2}} \emph{have} been generated by \texttt{PqSPStandardPresentation} (see{\nobreakspace}\texttt{PqSPStandardPresentation} (\ref{PqSPStandardPresentation})) then a \texttt{false} response from \texttt{PqSPCompareTwoFilePresentations} says the groups defined by those presentations are \emph{not} isomorphic. 

 For those familiar with the \texttt{pq} program, \texttt{PqSPCompareTwoFilePresentations} performs menu item 6 of the Standard Presentation menu. }

 

\subsection{\textcolor{Chapter }{PqSPIsomorphism}}
\logpage{[ 5, 8, 5 ]}\nobreak
\hyperdef{L}{X7D1277E584590ECE}{}
{\noindent\textcolor{FuncColor}{$\triangleright$\enspace\texttt{PqSPIsomorphism({\mdseries\slshape i})\index{PqSPIsomorphism@\texttt{PqSPIsomorphism}}
\label{PqSPIsomorphism}
}\hfill{\scriptsize (function)}}\\
\noindent\textcolor{FuncColor}{$\triangleright$\enspace\texttt{PqSPIsomorphism({\mdseries\slshape })\index{PqSPIsomorphism@\texttt{PqSPIsomorphism}!for default process}
\label{PqSPIsomorphism:for default process}
}\hfill{\scriptsize (function)}}\\


 for the \mbox{\texttt{\mdseries\slshape i}}th or default interactive \textsf{ANUPQ} process, direct the \texttt{pq} program to compute the isomorphism mapping from the $p$\texttt{\symbol{45}}group of the process to its standard presentation. This
function provides a description only; for a \textsf{GAP} object, use \texttt{EpimorphismStandardPresentation} (see{\nobreakspace}\texttt{EpimorphismStandardPresentation} (\ref{EpimorphismStandardPresentation:interactive})). 

 \emph{Note:} For those familiar with the \texttt{pq} program, \texttt{PqSPIsomorphism} performs menu item 8 of the Standard Presentation menu. }

 }

 
\section{\textcolor{Chapter }{Commands from the Main $p$\texttt{\symbol{45}}Group Generation menu}}\logpage{[ 5, 9, 0 ]}
\hyperdef{L}{X8490031B7AA7F237}{}
{
 Note that the $p$\texttt{\symbol{45}}group generation commands can only be applied once the \texttt{pq} program has produced a pc presentation of some quotient group of the ``group of the process''. 

\subsection{\textcolor{Chapter }{PqPGSupplyAutomorphisms}}
\logpage{[ 5, 9, 1 ]}\nobreak
\hyperdef{L}{X794A8D667A9D725A}{}
{\noindent\textcolor{FuncColor}{$\triangleright$\enspace\texttt{PqPGSupplyAutomorphisms({\mdseries\slshape i[, mlist]: options})\index{PqPGSupplyAutomorphisms@\texttt{PqPGSupplyAutomorphisms}}
\label{PqPGSupplyAutomorphisms}
}\hfill{\scriptsize (function)}}\\
\noindent\textcolor{FuncColor}{$\triangleright$\enspace\texttt{PqPGSupplyAutomorphisms({\mdseries\slshape [mlist]: options})\index{PqPGSupplyAutomorphisms@\texttt{PqPGSupplyAutomorphisms}!for default process}
\label{PqPGSupplyAutomorphisms:for default process}
}\hfill{\scriptsize (function)}}\\


 for the \mbox{\texttt{\mdseries\slshape i}}th or default interactive \textsf{ANUPQ} process, supply the \texttt{pq} program with the automorphism group data needed for the current quotient of
the group of that process (for process \mbox{\texttt{\mdseries\slshape i}} the group is stored as \texttt{ANUPQData.io[\mbox{\texttt{\mdseries\slshape i}}].group}). For a description of the format of \mbox{\texttt{\mdseries\slshape mlist}} see{\nobreakspace}\texttt{PqSupplyAutomorphisms} (\ref{PqSupplyAutomorphisms}). The options possible are \texttt{NumberOfSolubleAutomorphisms} and \texttt{RelativeOrders}. (Detailed descriptions of these options may be found in
Chapter{\nobreakspace}\hyperref[ANUPQ Options]{`ANUPQ Options'}.) 

 If \mbox{\texttt{\mdseries\slshape mlist}} is omitted, the automorphism data is determined from the group of the process
which must have been a $p$\texttt{\symbol{45}}group in pc presentation. 

 \emph{Note:} For those familiar with the \texttt{pq} program, \texttt{PqPGSupplyAutomorphisms} performs menu item 1 of the main $p$\texttt{\symbol{45}}Group Generation menu. }

 

\subsection{\textcolor{Chapter }{PqPGExtendAutomorphisms}}
\logpage{[ 5, 9, 2 ]}\nobreak
\hyperdef{L}{X87F251077C65B6DD}{}
{\noindent\textcolor{FuncColor}{$\triangleright$\enspace\texttt{PqPGExtendAutomorphisms({\mdseries\slshape i})\index{PqPGExtendAutomorphisms@\texttt{PqPGExtendAutomorphisms}}
\label{PqPGExtendAutomorphisms}
}\hfill{\scriptsize (function)}}\\
\noindent\textcolor{FuncColor}{$\triangleright$\enspace\texttt{PqPGExtendAutomorphisms({\mdseries\slshape })\index{PqPGExtendAutomorphisms@\texttt{PqPGExtendAutomorphisms}!for default process}
\label{PqPGExtendAutomorphisms:for default process}
}\hfill{\scriptsize (function)}}\\


 for the \mbox{\texttt{\mdseries\slshape i}}th or default interactive \textsf{ANUPQ} process, direct the \texttt{pq} program to compute the extensions of the automorphisms of the $p$\texttt{\symbol{45}}quotient of the previous class to the $p$\texttt{\symbol{45}}quotient of the current class. You may wish to set the \texttt{InfoLevel} of \texttt{InfoANUPQ} to 2 (or more) in order to see the output from the \texttt{pq} program (see{\nobreakspace}\texttt{InfoANUPQ} (\ref{InfoANUPQ})). 

 \emph{Note:} For those familiar with the \texttt{pq} program, \texttt{PqPGExtendAutomorphisms} performs menu item 2 of the main or advanced $p$\texttt{\symbol{45}}Group Generation menu. }

 

\subsection{\textcolor{Chapter }{PqPGConstructDescendants}}
\logpage{[ 5, 9, 3 ]}\nobreak
\hyperdef{L}{X82FD51A27D269E43}{}
{\noindent\textcolor{FuncColor}{$\triangleright$\enspace\texttt{PqPGConstructDescendants({\mdseries\slshape i: options})\index{PqPGConstructDescendants@\texttt{PqPGConstructDescendants}}
\label{PqPGConstructDescendants}
}\hfill{\scriptsize (function)}}\\
\noindent\textcolor{FuncColor}{$\triangleright$\enspace\texttt{PqPGConstructDescendants({\mdseries\slshape : options})\index{PqPGConstructDescendants@\texttt{PqPGConstructDescendants}!for default process}
\label{PqPGConstructDescendants:for default process}
}\hfill{\scriptsize (function)}}\\


 for the \mbox{\texttt{\mdseries\slshape i}}th or default interactive \textsf{ANUPQ} process, direct the \texttt{pq} program to construct descendants prescribed by \mbox{\texttt{\mdseries\slshape options}}, and return the number of descendants constructed (compare
function{\nobreakspace}\texttt{PqDescendants} (\ref{PqDescendants}) which returns the list of descendants). The options possible are \texttt{ClassBound}, \texttt{OrderBound}, \texttt{StepSize}, \texttt{PcgsAutomorphisms}, \texttt{RankInitialSegmentSubgroups}, \texttt{SpaceEfficient}, \texttt{CapableDescendants}, \texttt{AllDescendants}, \texttt{Exponent}, \texttt{Metabelian}, \texttt{BasicAlgorithm}, \texttt{CustomiseOutput}. (Detailed descriptions of these options may be found in
Chapter{\nobreakspace}\hyperref[ANUPQ Options]{`ANUPQ Options'}.) 

 \texttt{PqPGConstructDescendants} requires that the \texttt{pq} program has previously computed a pc presentation and a $p$\texttt{\symbol{45}}cover for a $p$\texttt{\symbol{45}}quotient of some class of the group of the process. 

 \emph{Note:} For those familiar with the \texttt{pq} program, \texttt{PqPGConstructDescendants} performs menu item 5 of the main $p$\texttt{\symbol{45}}Group Generation menu. }

 

\subsection{\textcolor{Chapter }{PqPGSetDescendantToPcp (with class)}}
\logpage{[ 5, 9, 4 ]}\nobreak
\hyperdef{L}{X7E9D511385D48A98}{}
{\noindent\textcolor{FuncColor}{$\triangleright$\enspace\texttt{PqPGSetDescendantToPcp({\mdseries\slshape i, cls, n})\index{PqPGSetDescendantToPcp@\texttt{PqPGSetDescendantToPcp}!with class}
\label{PqPGSetDescendantToPcp:with class}
}\hfill{\scriptsize (function)}}\\
\noindent\textcolor{FuncColor}{$\triangleright$\enspace\texttt{PqPGSetDescendantToPcp({\mdseries\slshape cls, n})\index{PqPGSetDescendantToPcp@\texttt{PqPGSetDescendantToPcp}!with class, for default process}
\label{PqPGSetDescendantToPcp:with class, for default process}
}\hfill{\scriptsize (function)}}\\
\noindent\textcolor{FuncColor}{$\triangleright$\enspace\texttt{PqPGSetDescendantToPcp({\mdseries\slshape i: [Filename := name]})\index{PqPGSetDescendantToPcp@\texttt{PqPGSetDescendantToPcp}}
\label{PqPGSetDescendantToPcp}
}\hfill{\scriptsize (function)}}\\
\noindent\textcolor{FuncColor}{$\triangleright$\enspace\texttt{PqPGSetDescendantToPcp({\mdseries\slshape : [Filename := name]})\index{PqPGSetDescendantToPcp@\texttt{PqPGSetDescendantToPcp}!for default process}
\label{PqPGSetDescendantToPcp:for default process}
}\hfill{\scriptsize (function)}}\\
\noindent\textcolor{FuncColor}{$\triangleright$\enspace\texttt{PqPGRestoreDescendantFromFile({\mdseries\slshape i, cls, n})\index{PqPGRestoreDescendantFromFile@\texttt{PqPGRestoreDescendantFromFile}!with class}
\label{PqPGRestoreDescendantFromFile:with class}
}\hfill{\scriptsize (function)}}\\
\noindent\textcolor{FuncColor}{$\triangleright$\enspace\texttt{PqPGRestoreDescendantFromFile({\mdseries\slshape cls, n})\index{PqPGRestoreDescendantFromFile@\texttt{PqPGRestoreDescendantFromFile}!with class, for default process}
\label{PqPGRestoreDescendantFromFile:with class, for default process}
}\hfill{\scriptsize (function)}}\\
\noindent\textcolor{FuncColor}{$\triangleright$\enspace\texttt{PqPGRestoreDescendantFromFile({\mdseries\slshape i: [Filename := name]})\index{PqPGRestoreDescendantFromFile@\texttt{PqPGRestoreDescendantFromFile}}
\label{PqPGRestoreDescendantFromFile}
}\hfill{\scriptsize (function)}}\\
\noindent\textcolor{FuncColor}{$\triangleright$\enspace\texttt{PqPGRestoreDescendantFromFile({\mdseries\slshape : [Filename := name]})\index{PqPGRestoreDescendantFromFile@\texttt{PqPGRestoreDescendantFromFile}!for default process}
\label{PqPGRestoreDescendantFromFile:for default process}
}\hfill{\scriptsize (function)}}\\


 for the \mbox{\texttt{\mdseries\slshape i}}th or default interactive \textsf{ANUPQ} process, direct the \texttt{pq} program to restore group \mbox{\texttt{\mdseries\slshape n}} of class \mbox{\texttt{\mdseries\slshape cls}} from a temporary file, where \mbox{\texttt{\mdseries\slshape cls}} and \mbox{\texttt{\mdseries\slshape n}} are positive integers, or the group stored in \mbox{\texttt{\mdseries\slshape name}}. \texttt{PqPGSetDescendantToPcp} and \texttt{PqPGRestoreDescendantFromFile} are synonyms; they make sense only after a prior call to construct descendants
by say \texttt{PqPGConstructDescendants} (see{\nobreakspace}\texttt{PqPGConstructDescendants} (\ref{PqPGConstructDescendants})) or the interactive \texttt{PqDescendants} (see{\nobreakspace}\texttt{PqDescendants} (\ref{PqDescendants:interactive})). In the \texttt{Filename} option forms, the option defaults to the last filename in which a presentation
was stored by the \texttt{pq} program. 

 \emph{Notes} 

 Since the \texttt{PqPGSetDescendantToPcp} and \texttt{PqPGRestoreDescendantFromFile} are intended to be used in calculation of further descendants the \texttt{pq} program computes the $p$\texttt{\symbol{45}}cover of the restored descendant. Hence, \texttt{PqCurrentGroup} used immediately after one of these commands returns the $p$\texttt{\symbol{45}}cover of the restored descendant rather than the
descendant itself. 

 For those familiar with the \texttt{pq} program, \texttt{PqPGSetDescendantToPcp} and \texttt{PqPGRestoreDescendantFromFile} perform menu item 3 of the main or advanced $p$\texttt{\symbol{45}}Group Generation menu. }

 }

 
\section{\textcolor{Chapter }{Commands from the Advanced $p$\texttt{\symbol{45}}Group Generation menu}}\logpage{[ 5, 10, 0 ]}
\hyperdef{L}{X7E8FBC2D83C43481}{}
{
 The functions below perform the component algorithms of \texttt{PqPGConstructDescendants} (see{\nobreakspace}\texttt{PqPGConstructDescendants} (\ref{PqPGConstructDescendants})). You can get some idea of their usage by trying \texttt{PqExample("Nott\texttt{\symbol{45}}APG\texttt{\symbol{45}}Rel\texttt{\symbol{45}}i");}. You can get some idea of the breakdown of \texttt{PqPGConstructDescendants} into these functions by comparing the previous output with \texttt{PqExample("Nott\texttt{\symbol{45}}PG\texttt{\symbol{45}}Rel\texttt{\symbol{45}}i");}. 

 These functions are intended for use only by ``experts''; please contact the authors of the package if you genuinely have a need for
them and need any amplified descriptions. 

\subsection{\textcolor{Chapter }{PqAPGDegree}}
\logpage{[ 5, 10, 1 ]}\nobreak
\hyperdef{L}{X7E539E4B78A6CE8D}{}
{\noindent\textcolor{FuncColor}{$\triangleright$\enspace\texttt{PqAPGDegree({\mdseries\slshape i, step, rank: [Exponent := n]})\index{PqAPGDegree@\texttt{PqAPGDegree}}
\label{PqAPGDegree}
}\hfill{\scriptsize (function)}}\\
\noindent\textcolor{FuncColor}{$\triangleright$\enspace\texttt{PqAPGDegree({\mdseries\slshape step, rank: [Exponent := n]})\index{PqAPGDegree@\texttt{PqAPGDegree}!for default process}
\label{PqAPGDegree:for default process}
}\hfill{\scriptsize (function)}}\\


      for the \mbox{\texttt{\mdseries\slshape i}}th or default interactive \textsf{ANUPQ} process, direct the \texttt{pq} program to invoke menu item 6 of the Advanced $p$\texttt{\symbol{45}}Group Generation menu. Here the
step\texttt{\symbol{45}}size \mbox{\texttt{\mdseries\slshape step}} and the rank \mbox{\texttt{\mdseries\slshape rank}} are positive integers and are the arguments required by the \texttt{pq} program. See{\nobreakspace}\ref{option Exponent} for the one recognised option \texttt{Exponent}.   }

 

\subsection{\textcolor{Chapter }{PqAPGPermutations}}
\logpage{[ 5, 10, 2 ]}\nobreak
\hyperdef{L}{X7F1182A77AB7650C}{}
{\noindent\textcolor{FuncColor}{$\triangleright$\enspace\texttt{PqAPGPermutations({\mdseries\slshape i: options})\index{PqAPGPermutations@\texttt{PqAPGPermutations}}
\label{PqAPGPermutations}
}\hfill{\scriptsize (function)}}\\
\noindent\textcolor{FuncColor}{$\triangleright$\enspace\texttt{PqAPGPermutations({\mdseries\slshape : options})\index{PqAPGPermutations@\texttt{PqAPGPermutations}!for default process}
\label{PqAPGPermutations:for default process}
}\hfill{\scriptsize (function)}}\\


   for the \mbox{\texttt{\mdseries\slshape i}}th or default interactive \textsf{ANUPQ} process, direct the \texttt{pq} program to perform menu item 7 of the Advanced $p$\texttt{\symbol{45}}Group Generation menu. Here the options \mbox{\texttt{\mdseries\slshape options}} recognised are \texttt{PcgsAutomorphisms}, \texttt{SpaceEfficient}, \texttt{PrintAutomorphisms} and \texttt{PrintPermutations} (see Chapter{\nobreakspace}\hyperref[ANUPQ Options]{`ANUPQ Options'} for details).   }

 

\subsection{\textcolor{Chapter }{PqAPGOrbits}}
\logpage{[ 5, 10, 3 ]}\nobreak
\hyperdef{L}{X7DD9E7507AA888CC}{}
{\noindent\textcolor{FuncColor}{$\triangleright$\enspace\texttt{PqAPGOrbits({\mdseries\slshape i: options})\index{PqAPGOrbits@\texttt{PqAPGOrbits}}
\label{PqAPGOrbits}
}\hfill{\scriptsize (function)}}\\
\noindent\textcolor{FuncColor}{$\triangleright$\enspace\texttt{PqAPGOrbits({\mdseries\slshape : options})\index{PqAPGOrbits@\texttt{PqAPGOrbits}!for default process}
\label{PqAPGOrbits:for default process}
}\hfill{\scriptsize (function)}}\\


     for the \mbox{\texttt{\mdseries\slshape i}}th or default interactive \textsf{ANUPQ} process, direct the \texttt{pq} to perform menu item 8 of the Advanced $p$\texttt{\symbol{45}}Group Generation menu. 

 Here the options \mbox{\texttt{\mdseries\slshape options}} recognised are \texttt{PcgsAutomorphisms}, \texttt{SpaceEfficient} and \texttt{CustomiseOutput} (see Chapter{\nobreakspace}\hyperref[ANUPQ Options]{`ANUPQ Options'} for details). For the \texttt{CustomiseOutput} option only the setting of the \texttt{orbit} is recognised (all other fields if set are ignored).   }

 

\subsection{\textcolor{Chapter }{PqAPGOrbitRepresentatives}}
\logpage{[ 5, 10, 4 ]}\nobreak
\hyperdef{L}{X795A756C78BE5923}{}
{\noindent\textcolor{FuncColor}{$\triangleright$\enspace\texttt{PqAPGOrbitRepresentatives({\mdseries\slshape i: options})\index{PqAPGOrbitRepresentatives@\texttt{PqAPGOrbitRepresentatives}}
\label{PqAPGOrbitRepresentatives}
}\hfill{\scriptsize (function)}}\\
\noindent\textcolor{FuncColor}{$\triangleright$\enspace\texttt{PqAPGOrbitRepresentatives({\mdseries\slshape : options})\index{PqAPGOrbitRepresentatives@\texttt{PqAPGOrbitRepresentatives}!for default process}
\label{PqAPGOrbitRepresentatives:for default process}
}\hfill{\scriptsize (function)}}\\


    for the \mbox{\texttt{\mdseries\slshape i}}th or default interactive \textsf{ANUPQ} process, direct the \texttt{pq} to perform item 9 of the Advanced $p$\texttt{\symbol{45}}Group Generation menu. 

 The options \mbox{\texttt{\mdseries\slshape options}} may be any selection of the following: \texttt{PcgsAutomorphisms}, \texttt{SpaceEfficient}, \texttt{Exponent}, \texttt{Metabelian}, \texttt{CapableDescendants} (or \texttt{AllDescendants}), \texttt{CustomiseOutput} (where only the \texttt{group} and \texttt{autgroup} fields are recognised) and \texttt{Filename} (see Chapter{\nobreakspace}\hyperref[ANUPQ Options]{`ANUPQ Options'} for details). If \texttt{Filename} is omitted the reduced $p$\texttt{\symbol{45}}cover is written to the file \texttt{"redPCover"} in the temporary directory whose name is stored in \texttt{ANUPQData.tmpdir}.    }

 

\subsection{\textcolor{Chapter }{PqAPGSingleStage}}
\logpage{[ 5, 10, 5 ]}\nobreak
\hyperdef{L}{X7E5D6D5782FE190E}{}
{\noindent\textcolor{FuncColor}{$\triangleright$\enspace\texttt{PqAPGSingleStage({\mdseries\slshape i: options})\index{PqAPGSingleStage@\texttt{PqAPGSingleStage}}
\label{PqAPGSingleStage}
}\hfill{\scriptsize (function)}}\\
\noindent\textcolor{FuncColor}{$\triangleright$\enspace\texttt{PqAPGSingleStage({\mdseries\slshape : options})\index{PqAPGSingleStage@\texttt{PqAPGSingleStage}!for default process}
\label{PqAPGSingleStage:for default process}
}\hfill{\scriptsize (function)}}\\


    for the \mbox{\texttt{\mdseries\slshape i}}th or default interactive \textsf{ANUPQ} process, direct the \texttt{pq} to perform option 5 of the Advanced $p$\texttt{\symbol{45}}Group Generation menu. 

 The possible options are \texttt{StepSize}, \texttt{PcgsAutomorphisms}, \texttt{RankInitialSegmentSubgroups}, \texttt{SpaceEfficient}, \texttt{CapableDescendants}, \texttt{AllDescendants}, \texttt{Exponent}, \texttt{Metabelian}, \texttt{BasicAlgorithm} and \texttt{CustomiseOutput}. (Detailed descriptions of these options may be found in
Chapter{\nobreakspace}\hyperref[ANUPQ Options]{`ANUPQ Options'}.)    }

 }

 
\section{\textcolor{Chapter }{Primitive Interactive ANUPQ Process Read/Write Functions}}\label{Primitive Interactive ANUPQ Process Read/Write Functions}
\logpage{[ 5, 11, 0 ]}
\hyperdef{L}{X84A92E8087762DEE}{}
{
  For those familiar with using the \texttt{pq} program as a standalone we provide primitive read/write tools to communicate
directly with an interactive \textsf{ANUPQ} process, started via \texttt{PqStart}. For the most part, it is up to the user to translate the output strings from \texttt{pq} program into a form useful in \textsf{GAP}. 

\subsection{\textcolor{Chapter }{PqRead}}
\logpage{[ 5, 11, 1 ]}\nobreak
\hyperdef{L}{X8464AF7E83A523C1}{}
{\noindent\textcolor{FuncColor}{$\triangleright$\enspace\texttt{PqRead({\mdseries\slshape i})\index{PqRead@\texttt{PqRead}}
\label{PqRead}
}\hfill{\scriptsize (function)}}\\
\noindent\textcolor{FuncColor}{$\triangleright$\enspace\texttt{PqRead({\mdseries\slshape })\index{PqRead@\texttt{PqRead}!for default process}
\label{PqRead:for default process}
}\hfill{\scriptsize (function)}}\\


 read a complete line of \textsf{ANUPQ} output, from the \mbox{\texttt{\mdseries\slshape i}}th or default interactive \textsf{ANUPQ} process, if there is output to be read and returns \texttt{fail} otherwise. When successful, the line is returned as a string complete with
trailing newline, colon, or question\texttt{\symbol{45}}mark character. Please
note that it is possible to be ``too quick'' (i.e.{\nobreakspace}the return can be \texttt{fail} purely because the output from \textsf{ANUPQ} is not there yet), but if \texttt{PqRead} finds any output at all, it waits for a complete line. \texttt{PqRead} also writes the line read via \texttt{Info} at \texttt{InfoANUPQ} level 2. It doesn't try to distinguish banner and menu output from other
output of the \texttt{pq} program. }

 

\subsection{\textcolor{Chapter }{PqReadAll}}
\logpage{[ 5, 11, 2 ]}\nobreak
\hyperdef{L}{X7ADD82207D049A87}{}
{\noindent\textcolor{FuncColor}{$\triangleright$\enspace\texttt{PqReadAll({\mdseries\slshape i})\index{PqReadAll@\texttt{PqReadAll}}
\label{PqReadAll}
}\hfill{\scriptsize (function)}}\\
\noindent\textcolor{FuncColor}{$\triangleright$\enspace\texttt{PqReadAll({\mdseries\slshape })\index{PqReadAll@\texttt{PqReadAll}!for default process}
\label{PqReadAll:for default process}
}\hfill{\scriptsize (function)}}\\


 read and return as many \emph{complete} lines of \textsf{ANUPQ} output, from the \mbox{\texttt{\mdseries\slshape i}}th or default interactive \textsf{ANUPQ} process, as there are to be read, \emph{at the time of the call}, as a list of strings with any trailing newlines removed and returns the
empty list otherwise. \texttt{PqReadAll} also writes each line read via \texttt{Info} at \texttt{InfoANUPQ} level 2. It doesn't try to distinguish banner and menu output from other
output of the \texttt{pq} program. Whenever \texttt{PqReadAll} finds only a partial line, it waits for the complete line, thus increasing the
probability that it has captured all the output to be had from \textsf{ANUPQ}. }

 

\subsection{\textcolor{Chapter }{PqReadUntil}}
\logpage{[ 5, 11, 3 ]}\nobreak
\hyperdef{L}{X7EBF3C3681671879}{}
{\noindent\textcolor{FuncColor}{$\triangleright$\enspace\texttt{PqReadUntil({\mdseries\slshape i, IsMyLine})\index{PqReadUntil@\texttt{PqReadUntil}}
\label{PqReadUntil}
}\hfill{\scriptsize (function)}}\\
\noindent\textcolor{FuncColor}{$\triangleright$\enspace\texttt{PqReadUntil({\mdseries\slshape IsMyLine})\index{PqReadUntil@\texttt{PqReadUntil}!for default process}
\label{PqReadUntil:for default process}
}\hfill{\scriptsize (function)}}\\
\noindent\textcolor{FuncColor}{$\triangleright$\enspace\texttt{PqReadUntil({\mdseries\slshape i, IsMyLine, Modify})\index{PqReadUntil@\texttt{PqReadUntil}!with modify map}
\label{PqReadUntil:with modify map}
}\hfill{\scriptsize (function)}}\\
\noindent\textcolor{FuncColor}{$\triangleright$\enspace\texttt{PqReadUntil({\mdseries\slshape IsMyLine, Modify})\index{PqReadUntil@\texttt{PqReadUntil}!with modify map, for default process}
\label{PqReadUntil:with modify map, for default process}
}\hfill{\scriptsize (function)}}\\


 read complete lines of \textsf{ANUPQ} output, from the \mbox{\texttt{\mdseries\slshape i}}th or default interactive \textsf{ANUPQ} process, ``chomps'' them (i.e.{\nobreakspace}removes any trailing newline character), emits them
to \texttt{Info} at \texttt{InfoANUPQ} level 2 (without trying to distinguish banner and menu output from other
output of the \texttt{pq} program), and applies the function \mbox{\texttt{\mdseries\slshape Modify}} (where \mbox{\texttt{\mdseries\slshape Modify}} is just the identity map/function for the first two forms) until a ``chomped'' line \mbox{\texttt{\mdseries\slshape line}} for which \texttt{\mbox{\texttt{\mdseries\slshape IsMyLine}}( \mbox{\texttt{\mdseries\slshape Modify}}(\mbox{\texttt{\mdseries\slshape line}}) )} is true. \texttt{PqReadUntil} returns the list of \mbox{\texttt{\mdseries\slshape Modify}}\texttt{\symbol{45}}ed ``chomped'' lines read. 

 \emph{Notes:} When provided by the user, \mbox{\texttt{\mdseries\slshape Modify}} should be a function that accepts a single string argument. 

 \mbox{\texttt{\mdseries\slshape IsMyLine}} should be a function that is able to accept the output of \mbox{\texttt{\mdseries\slshape Modify}} (or take a single string argument when \mbox{\texttt{\mdseries\slshape Modify}} is not provided) and should return a boolean. 

 If \texttt{\mbox{\texttt{\mdseries\slshape IsMyLine}}( \mbox{\texttt{\mdseries\slshape Modify}}(\mbox{\texttt{\mdseries\slshape line}}) )} is never true, \texttt{PqReadUntil} will wait indefinitely. }

 

\subsection{\textcolor{Chapter }{PqWrite}}
\logpage{[ 5, 11, 4 ]}\nobreak
\hyperdef{L}{X83CBE77F78057E21}{}
{\noindent\textcolor{FuncColor}{$\triangleright$\enspace\texttt{PqWrite({\mdseries\slshape i, string})\index{PqWrite@\texttt{PqWrite}}
\label{PqWrite}
}\hfill{\scriptsize (function)}}\\
\noindent\textcolor{FuncColor}{$\triangleright$\enspace\texttt{PqWrite({\mdseries\slshape string})\index{PqWrite@\texttt{PqWrite}!for default process}
\label{PqWrite:for default process}
}\hfill{\scriptsize (function)}}\\


 write \mbox{\texttt{\mdseries\slshape string}} to the \mbox{\texttt{\mdseries\slshape i}}th or default interactive \textsf{ANUPQ} process; \mbox{\texttt{\mdseries\slshape string}} must be in exactly the form the \textsf{ANUPQ} standalone expects. The command is echoed via \texttt{Info} at \texttt{InfoANUPQ} level 3 (with a ``\texttt{ToPQ{\textgreater} }'' prompt); i.e.{\nobreakspace}do \texttt{SetInfoLevel(InfoANUPQ, 3);} to see what is transmitted to the \texttt{pq} program. \texttt{PqWrite} returns \texttt{true} if successful in writing to the stream of the interactive \textsf{ANUPQ} process, and \texttt{fail} otherwise. 

 \emph{Note:} If \texttt{PqWrite} returns \texttt{fail} it means that the \textsf{ANUPQ} process has died. }

 }

 }

 
\chapter{\textcolor{Chapter }{ANUPQ Options}}\label{ANUPQ Options}
\logpage{[ 6, 0, 0 ]}
\hyperdef{L}{X7B0946E97E7EB359}{}
{
  
\section{\textcolor{Chapter }{Overview}}\logpage{[ 6, 1, 0 ]}
\hyperdef{L}{X8389AD927B74BA4A}{}
{
  In this chapter we describe in detail all the options used by functions of the \textsf{ANUPQ} package. Note that by ``options'' we mean \textsf{GAP} options that are passed to functions after the arguments and separated from
the arguments by a colon as described in Chapter{\nobreakspace} \textbf{Reference: Function Calls} in the Reference Manual. The user is strongly advised to read
Section{\nobreakspace}\hyperref[Hints and Warnings regarding the use of Options]{`Hints and Warnings regarding the use of Options'}. 

\subsection{\textcolor{Chapter }{AllANUPQoptions}}
\logpage{[ 6, 1, 1 ]}\nobreak
\hyperdef{L}{X7846E25D8065D78F}{}
{\noindent\textcolor{FuncColor}{$\triangleright$\enspace\texttt{AllANUPQoptions({\mdseries\slshape })\index{AllANUPQoptions@\texttt{AllANUPQoptions}}
\label{AllANUPQoptions}
}\hfill{\scriptsize (function)}}\\


 lists all the \textsf{GAP} options defined for functions of the \textsf{ANUPQ} package: 
\begin{Verbatim}[commandchars=!@|,fontsize=\small,frame=single,label=Example]
  !gapprompt@gap>| !gapinput@AllANUPQoptions();|
  [ "AllDescendants", "BasicAlgorithm", "Bounds", "CapableDescendants", 
    "ClassBound", "CustomiseOutput", "Exponent", "Filename", "GroupName", 
    "Identities", "Metabelian", "NumberOfSolubleAutomorphisms", "OrderBound", 
    "OutputLevel", "PcgsAutomorphisms", "PqWorkspace", "Prime", 
    "PrintAutomorphisms", "PrintPermutations", "QueueFactor", 
    "RankInitialSegmentSubgroups", "RedoPcp", "RelativeOrders", "Relators", 
    "SetupFile", "SpaceEfficient", "StandardPresentationFile", "StepSize", 
    "SubList", "TreeDepth", "pQuotient" ]
\end{Verbatim}
 }

 The following global variable gives a partial breakdown of where the above
options are used. 

\subsection{\textcolor{Chapter }{ANUPQoptions}}
\logpage{[ 6, 1, 2 ]}\nobreak
\hyperdef{L}{X82D560C77D195089}{}
{\noindent\textcolor{FuncColor}{$\triangleright$\enspace\texttt{ANUPQoptions\index{ANUPQoptions@\texttt{ANUPQoptions}}
\label{ANUPQoptions}
}\hfill{\scriptsize (global variable)}}\\


 is a record of lists of names of admissible \textsf{ANUPQ} options, such that each field is either the name of a ``key'' \textsf{ANUPQ} function or \texttt{other} (for a miscellaneous list of functions) and the corresponding value is the
list of option names that are admissible for the function (or miscellaneous
list of functions). 

 Also, from within a \textsf{GAP} session, you may use \textsf{GAP}'s help browser (see Chapter{\nobreakspace} \textbf{Reference: The Help System} in the \textsf{GAP} Reference Manual); to find out about any particular \textsf{ANUPQ} option, simply type: ``\texttt{?option \mbox{\texttt{\mdseries\slshape option}}}'', where \mbox{\texttt{\mdseries\slshape option}} is one of the options listed above without any quotes, e.g. 
\begin{Verbatim}[commandchars=!@|,fontsize=\small,frame=single,label=Example]
  !gapprompt@gap>| !gapinput@?option Prime|
\end{Verbatim}
 will display the sections in this manual that describe the \texttt{Prime} option. In fact the first 4 are for the functions that have \texttt{Prime} as an option and the last actually describes the option. So follow up by
choosing  
\begin{Verbatim}[commandchars=!@|,fontsize=\small,frame=single,label=Example]
  !gapprompt@gap>| !gapinput@?5|
\end{Verbatim}
 This is also the pattern for other options (the last section of the list
always describes the option; the other sections are the functions with which
the option may be used). 

 In the section following we describe in detail all \textsf{ANUPQ} options. To continue onto the next section on\texttt{\symbol{45}}line using \textsf{GAP}'s help browser, type: 
\begin{Verbatim}[commandchars=!@|,fontsize=\small,frame=single,label=Example]
  !gapprompt@gap>| !gapinput@?>|
\end{Verbatim}
 }

 }

 
\section{\textcolor{Chapter }{Detailed descriptions of ANUPQ Options}}\label{Detailed descriptions of ANUPQ Options}
\logpage{[ 6, 2, 0 ]}
\hyperdef{L}{X87EE6DD67C9996A3}{}
{
  
\begin{description}
\item[{\texttt{Prime := \mbox{\texttt{\mdseries\slshape p}}} \label{option Prime}\index{option Prime}}]  Specifies that the $p$\texttt{\symbol{45}}quotient for the prime \mbox{\texttt{\mdseries\slshape p}} should be computed. 
\item[{\texttt{ClassBound := \mbox{\texttt{\mdseries\slshape n}}} \label{option ClassBound}\index{option ClassBound}}]  Specifies that the $p$\texttt{\symbol{45}}quotient to be computed has lower
exponent\texttt{\symbol{45}}$p$ class at most \mbox{\texttt{\mdseries\slshape n}}. If this option is omitted a default of 63 (which is the maximum possible for
the \texttt{pq} program) is taken, except for \texttt{PqDescendants} (see{\nobreakspace}\texttt{PqDescendants} (\ref{PqDescendants})) and in a special case of \texttt{PqPCover} (see{\nobreakspace}\texttt{PqPCover} (\ref{PqPCover})). Let \mbox{\texttt{\mdseries\slshape F}} be the argument (or start group of the process in the interactive case) for
the function; then for \texttt{PqDescendants} the default is \texttt{PClassPGroup(\mbox{\texttt{\mdseries\slshape F}}) + 1}, and for the special case of \texttt{PqPCover} the default is \texttt{PClassPGroup(\mbox{\texttt{\mdseries\slshape F}})}. 
\item[{\texttt{pQuotient := \mbox{\texttt{\mdseries\slshape Q}}} \label{option pQuotient}\index{option pQuotient}}]  This option is only available for the standard presentation functions. It
specifies that a $p$\texttt{\symbol{45}}quotient of the group argument of the function or group of
the process is the pc \mbox{\texttt{\mdseries\slshape p}}\texttt{\symbol{45}}group \mbox{\texttt{\mdseries\slshape Q}}, where \mbox{\texttt{\mdseries\slshape Q}} is of class \emph{less than} the provided (or default) value of \texttt{ClassBound}. If \texttt{pQuotient} is provided, then the option \texttt{Prime} if also provided, is ignored; the prime \mbox{\texttt{\mdseries\slshape p}} is discovered by computing \texttt{PrimePGroup(\mbox{\texttt{\mdseries\slshape Q}})}. 
\item[{\texttt{Exponent := \mbox{\texttt{\mdseries\slshape n}}} \label{option Exponent}\index{option Exponent}}]  Specifies that the $p$\texttt{\symbol{45}}quotient to be computed has exponent \mbox{\texttt{\mdseries\slshape n}}. For an interactive process, \texttt{Exponent} defaults to a previously supplied value for the process. Otherwise (and
non\texttt{\symbol{45}}interactively), the default is 0, which means that no
exponent law is enforced. 
\item[{\texttt{Relators := \mbox{\texttt{\mdseries\slshape rels}}} \label{option Relators}\index{option Relators}}]  Specifies that the relators sent to the \texttt{pq} program should be \mbox{\texttt{\mdseries\slshape rels}} instead of the relators of the argument group \mbox{\texttt{\mdseries\slshape F}} (or start group in the interactive case) of the calling function; \mbox{\texttt{\mdseries\slshape rels}} should be a list of \emph{strings} in the string representations of the generators of \mbox{\texttt{\mdseries\slshape F}}, and \mbox{\texttt{\mdseries\slshape F}} must be an \emph{fp group} (even if the calling function accepts a pc group). This option provides a way
of giving relators to the \texttt{pq} program, without having them pre\texttt{\symbol{45}}expanded by \textsf{GAP}, which can sometimes effect a performance loss of the order of 100 (see
Section{\nobreakspace}\hyperref[The Relators Option]{`The Relators Option'}). 

 \emph{Notes} 
\begin{enumerate}
\item  The \texttt{pq} program does not use \texttt{/} to indicate multiplication by an inverse and uses square brackets to represent
(left normed) commutators. Also, even though the \texttt{pq} program accepts relations, all elements of \mbox{\texttt{\mdseries\slshape rels}} \emph{must} be in relator form, i.e.{\nobreakspace}a relation of form \texttt{\mbox{\texttt{\mdseries\slshape w1}} = \mbox{\texttt{\mdseries\slshape w2}}} must be written as \texttt{\mbox{\texttt{\mdseries\slshape w1}}*(\mbox{\texttt{\mdseries\slshape w2}})\texttt{\symbol{94}}\texttt{\symbol{45}}1} and then put in a pair of double\texttt{\symbol{45}}quotes to make it a
string. See the example below. 
\item  To ensure there are no syntax errors in \mbox{\texttt{\mdseries\slshape rels}}, each relator is parsed for validity via \texttt{PqParseWord} (see{\nobreakspace}\texttt{PqParseWord} (\ref{PqParseWord})). If they are ok, a message to say so is \texttt{Info}\texttt{\symbol{45}}ed at \texttt{InfoANUPQ} level 2. 
\end{enumerate}
 
\item[{\texttt{Metabelian} \label{option Metabelian}\index{option Metabelian}}]  Specifies that the largest metabelian $p$\texttt{\symbol{45}}quotient subject to any other conditions specified by
other options be constructed. By default this restriction is not enforced. 
\item[{\texttt{GroupName := \mbox{\texttt{\mdseries\slshape name}}} \label{option GroupName}\index{option GroupName}}]  Specifies that the \texttt{pq} program should refer to the group by the name \mbox{\texttt{\mdseries\slshape name}} (a string). If \texttt{GroupName} is not set and the group has been assigned a name via \texttt{SetName} (see{\nobreakspace} \textbf{Reference: Name}) it is set as the name the \texttt{pq} program should use. Otherwise, the ``generic'' name \texttt{"[grp]"} is set as a default. 
\item[{\texttt{Identities := \mbox{\texttt{\mdseries\slshape funcs}}} \label{option Identities}\index{option Identities}}]  Specifies that the pc presentation should satisfy the laws defined by each
function in the list \mbox{\texttt{\mdseries\slshape funcs}}. This option may be called by \texttt{Pq}, \texttt{PqEpimorphism}, or \texttt{PqPCover} (see{\nobreakspace}\texttt{Pq} (\ref{Pq})). Each function in the list \mbox{\texttt{\mdseries\slshape funcs}} must return a word in its arguments (there may be any number of arguments).
Let \mbox{\texttt{\mdseries\slshape identity}} be one such function in \mbox{\texttt{\mdseries\slshape funcs}}. Then as each lower exponent \mbox{\texttt{\mdseries\slshape p}}\texttt{\symbol{45}}class quotient is formed, instances $\mbox{\texttt{\mdseries\slshape identity}}(\mbox{\texttt{\mdseries\slshape w1}}, \dots, \mbox{\texttt{\mdseries\slshape wn}})$ are added as relators to the pc presentation, where $\mbox{\texttt{\mdseries\slshape w1}}, \dots, \mbox{\texttt{\mdseries\slshape wn}}$ are words in the pc generators of the quotient. At each class the class and
number of pc generators is \texttt{Info}\texttt{\symbol{45}}ed at \texttt{InfoANUPQ} level 1, the number of instances is \texttt{Info}\texttt{\symbol{45}}ed at \texttt{InfoANUPQ} level 2, and the instances that are evaluated are \texttt{Info}\texttt{\symbol{45}}ed at \texttt{InfoANUPQ} level 3. As usual timing information is \texttt{Info}\texttt{\symbol{45}}ed at \texttt{InfoANUPQ} level 2; and details of the processing of each instance from the \texttt{pq} program (which is often quite \emph{voluminous}) is \texttt{Info}\texttt{\symbol{45}}ed at \texttt{InfoANUPQ} level 3. Try the examples \texttt{"B2\texttt{\symbol{45}}4\texttt{\symbol{45}}Id"} and \texttt{"11gp\texttt{\symbol{45}}3\texttt{\symbol{45}}Engel\texttt{\symbol{45}}Id"} which demonstrate the usage of the \texttt{Identities} option; these are run using \texttt{PqExample} (see{\nobreakspace}\texttt{PqExample} (\ref{PqExample})). Take note of Note 1.{\nobreakspace}below in relation to the example \texttt{"B2\texttt{\symbol{45}}4\texttt{\symbol{45}}Id"}; the companion example \texttt{"B2\texttt{\symbol{45}}4"} generates the same group using the \texttt{Exponent} option. These examples are discussed at length in Section{\nobreakspace}\hyperref[The Identities Option and PqEvaluateIdentities Function]{`The Identities Option and PqEvaluateIdentities Function'}. 

 \emph{Notes} 
\begin{enumerate}
\item  Setting the \texttt{InfoANUPQ} level to 3 or more when setting the \texttt{Identities} option may slow down the computation considerably, by overloading \textsf{GAP} with io operations. 
\item  The \texttt{Identities} option is implemented at the \textsf{GAP} level. An identity that is just an exponent law should be specified using the \texttt{Exponent} option (see{\nobreakspace}\texttt{option Exponent} (\ref{option Exponent})), which is implemented at the C level and is highly optimised and so is much
more efficient. 
\item  The number of instances of each identity tends to grow combinatorially with
the class. So \emph{care} should be exercised in using the \texttt{Identities} option, by including other restrictions, e.g.{\nobreakspace}by using the \texttt{ClassBound} option (see{\nobreakspace}\texttt{option ClassBound} (\ref{option ClassBound})). 
\end{enumerate}
 
\item[{\texttt{OutputLevel := \mbox{\texttt{\mdseries\slshape n}}} \label{option OutputLevel}\index{option OutputLevel}}]  Specifies the level of ``verbosity'' of the information output by the ANU \texttt{pq} program when computing a pc presentation; \mbox{\texttt{\mdseries\slshape n}} must be an integer in the range 0 to 3. \texttt{OutputLevel := 0} displays at most one line of output and is the default; \texttt{OutputLevel := 1} displays (usually) slightly more output and \texttt{OutputLevel}s of 2 and 3 are two levels of verbose output. To see these messages from the \texttt{pq} program, the \texttt{InfoANUPQ} level must be set to at least 1 (see{\nobreakspace}\texttt{InfoANUPQ} (\ref{InfoANUPQ})). See Section{\nobreakspace}\hyperref[Hints and Warnings regarding the use of Options]{`Hints and Warnings regarding the use of Options'} for an example of how \texttt{OutputLevel} can be used as a troubleshooting tool. 
\item[{\texttt{RedoPcp} \label{option RedoPcp}\index{option RedoPcp}}]  Specifies that the current pc presentation (for an interactive process) stored
by the \texttt{pq} program be scrapped and clears the current values stored for the options \texttt{Prime}, \texttt{ClassBound}, \texttt{Exponent} and \texttt{Metabelian} and also clears the \texttt{pQuotient}, \texttt{pQepi} and \texttt{pCover} fields of the data record of the process. 
\item[{\texttt{SetupFile := \mbox{\texttt{\mdseries\slshape filename}}} \label{option SetupFile}\index{option SetupFile}}]  Non\texttt{\symbol{45}}interactively, this option directs that \texttt{pq} should not be called and that an input file with name \mbox{\texttt{\mdseries\slshape filename}} (a string), containing the commands necessary for the ANU \texttt{pq} standalone, be constructed. The commands written to \mbox{\texttt{\mdseries\slshape filename}} are also \texttt{Info}\texttt{\symbol{45}}ed behind a ``\texttt{ToPQ{\textgreater} }'' prompt at \texttt{InfoANUPQ} level 4 (see{\nobreakspace}\texttt{InfoANUPQ} (\ref{InfoANUPQ})). Except in the case following, the calling function returns \texttt{true}. If the calling function is the non\texttt{\symbol{45}}interactive version of
one of \texttt{Pq}, \texttt{PqPCover} or \texttt{PqEpimorphism} and the group provided as argument is trivial given with an empty set of
generators, then no setup file is written and \texttt{fail} is returned (the \texttt{pq} program cannot do anything useful with such a group). Interactively, \texttt{SetupFile} is ignored. 

 \emph{Note:} Since commands emitted to the \texttt{pq} program may depend on knowing what the ``current state'' is, to form a setup file some ``close enough guesses'' may sometimes be necessary; when this occurs a warning is \texttt{Info}\texttt{\symbol{45}}ed at \texttt{InfoANUPQ} or \texttt{InfoWarning} level 1. To determine whether the ``close enough guesses'' give an accurate setup file, it is necessary to run the command without the \texttt{SetupFile} option, after either setting the \texttt{InfoANUPQ} level to at least 4 (the setup file script can then be compared with the ``\texttt{ToPQ{\textgreater} }'' commands that are \texttt{Info}\texttt{\symbol{45}}ed) or setting a \texttt{pq} command log file by using \texttt{ToPQLog} (see{\nobreakspace}\texttt{ToPQLog} (\ref{ToPQLog})). 
\item[{\texttt{PqWorkspace := \mbox{\texttt{\mdseries\slshape workspace}}} \label{option PqWorkspace}\index{option PqWorkspace}}]  Non\texttt{\symbol{45}}interactively, this option sets the memory used by the \texttt{pq} program. It sets the maximum number of integer\texttt{\symbol{45}}sized
elements to allocate in its main storage array. By default, the \texttt{pq} program sets this figure to 10000000. Interactively, \texttt{PqWorkspace} is ignored; the memory used in this case may be set by giving \texttt{PqStart} a second argument (see{\nobreakspace}\texttt{PqStart} (\ref{PqStart})). 
\item[{\texttt{PcgsAutomorphisms} \label{option PcgsAutomorphisms}\index{option PcgsAutomorphisms}}] 
\item[{\texttt{PcgsAutomorphisms := false}  }]  Let \mbox{\texttt{\mdseries\slshape G}} be the group associated with the calling function (or associated interactive
process). Passing the option \texttt{PcgsAutomorphisms} without a value (or equivalently setting it to \texttt{true}), specifies that a polycyclic generating sequence for the automorphism group
(which must be \emph{soluble}) of \mbox{\texttt{\mdseries\slshape G}}, be computed and passed to the \texttt{pq} program. This increases the efficiency of the computation; it also prevents
the \texttt{pq} from calling \textsf{GAP} for orbit\texttt{\symbol{45}}stabilizer calculations. By default, \texttt{PcgsAutomorphisms} is set to the value returned by \texttt{IsSolvable( AutomorphismGroup( \mbox{\texttt{\mdseries\slshape G}} ) )}, and uses the package \textsf{AutPGrp} to compute \texttt{AutomorphismGroup( \mbox{\texttt{\mdseries\slshape G}} )} if it is installed. This flag is set to \texttt{true} or \texttt{false} in the background according to the above criterion by the function \texttt{PqDescendants} (see{\nobreakspace}\texttt{PqDescendants} (\ref{PqDescendants}) and{\nobreakspace}\texttt{PqDescendants} (\ref{PqDescendants:interactive})). 

 \emph{Note:} If \texttt{PcgsAutomorphisms} is used when the automorphism group of \mbox{\texttt{\mdseries\slshape G}} is insoluble, an error message occurs. 
\item[{\texttt{OrderBound := \mbox{\texttt{\mdseries\slshape n}}} \label{option OrderBound}\index{option OrderBound}}]  Specifies that only descendants of size at most $p^{\mbox{\texttt{\mdseries\slshape n}}}$, where \mbox{\texttt{\mdseries\slshape n}} is a non\texttt{\symbol{45}}negative integer, be generated. Note that you
cannot set both \texttt{OrderBound} and \texttt{StepSize}. 
\item[{\texttt{StepSize := \mbox{\texttt{\mdseries\slshape n}}} \label{option StepSize}\index{option StepSize}}] 
\item[{\texttt{StepSize := \mbox{\texttt{\mdseries\slshape list}}}  }]  For a positive integer \mbox{\texttt{\mdseries\slshape n}}, \texttt{StepSize} specifies that only those immediate descendants which are a factor $p^{\mbox{\texttt{\mdseries\slshape n}}}$ bigger than their parent group be generated. 

 For a list \mbox{\texttt{\mdseries\slshape list}} of positive integers such that the sum of the length of \mbox{\texttt{\mdseries\slshape list}} and the exponent\texttt{\symbol{45}}$p$ class of \mbox{\texttt{\mdseries\slshape G}} is equal to the class bound defined by the option \texttt{ClassBound}, \texttt{StepSize} specifies that the integers of \mbox{\texttt{\mdseries\slshape list}} are the step sizes for each additional class. 
\item[{\texttt{RankInitialSegmentSubgroups := \mbox{\texttt{\mdseries\slshape n}}} \label{option RankInitialSegmentSubgroups}\index{option RankInitialSegmentSubgroups}}]  Sets the rank of the initial segment subgroup chosen to be \mbox{\texttt{\mdseries\slshape n}}. By default, this has value 0. 
\item[{\texttt{SpaceEfficient} \label{option SpaceEfficient}\index{option SpaceEfficient}}]  Specifies that the \texttt{pq} program performs certain calculations of $p$\texttt{\symbol{45}}group generation more slowly but with greater space
efficiency. This flag is frequently necessary for groups of large Frattini
quotient rank. The space saving occurs because only one permutation is stored
at any one time. This option is only available if the \texttt{PcgsAutomorphisms} flag is set to \texttt{true} (see{\nobreakspace}\texttt{option PcgsAutomorphisms} (\ref{option PcgsAutomorphisms})). For an interactive process, \texttt{SpaceEfficient} defaults to a previously supplied value for the process. Otherwise (and
non\texttt{\symbol{45}}interactively), \texttt{SpaceEfficient} is by default \texttt{false}. 
\item[{\texttt{CapableDescendants} \label{option CapableDescendants}\index{option CapableDescendants}}]  By default, \emph{all} (i.e.{\nobreakspace}capable and terminal) descendants are computed. If this
flag is set, only capable descendants are computed. Setting this option is
equivalent to setting \texttt{AllDescendants := false} (see{\nobreakspace}\texttt{option AllDescendants} (\ref{option AllDescendants})), except if both \texttt{CapableDescendants} and \texttt{AllDescendants} are passed, \texttt{AllDescendants} is essentially ignored. 
\item[{\texttt{AllDescendants := false} \label{option AllDescendants}\index{option AllDescendants}}]  By default, \emph{all} descendants are constructed. If this flag is set to \texttt{false}, only capable descendants are computed. Passing \texttt{AllDescendants} without a value (which is equivalent to setting it to \texttt{true}) is superfluous. This option is provided only for backward compatibility with
the \textsf{GAP} 3 version of the \textsf{ANUPQ} package, where by default \texttt{AllDescendants} was set to \texttt{false} (rather than \texttt{true}). It is preferable to use \texttt{CapableDescendants} (see{\nobreakspace}\texttt{option CapableDescendants} (\ref{option CapableDescendants})). 
\item[{\texttt{TreeDepth := \mbox{\texttt{\mdseries\slshape class}}} \label{option TreeDepth}\index{option TreeDepth}}]  Specifies that the descendants tree developed by \texttt{PqDescendantsTreeCoclassOne} (see{\nobreakspace}\texttt{PqDescendantsTreeCoclassOne} (\ref{PqDescendantsTreeCoclassOne})) should be extended to class \mbox{\texttt{\mdseries\slshape class}}, where \mbox{\texttt{\mdseries\slshape class}} is a positive integer. 
\item[{\texttt{SubList := \mbox{\texttt{\mdseries\slshape sub}}} \label{option SubList}\index{option SubList}}]  Suppose that \mbox{\texttt{\mdseries\slshape L}} is the list of descendants generated, then for a list \mbox{\texttt{\mdseries\slshape sub}} of integers this option causes \texttt{PqDescendants} to return \texttt{Sublist( \mbox{\texttt{\mdseries\slshape L}}, \mbox{\texttt{\mdseries\slshape sub}} )}. If an integer \mbox{\texttt{\mdseries\slshape n}} is supplied, \texttt{PqDescendants} returns \texttt{\mbox{\texttt{\mdseries\slshape L}}[\mbox{\texttt{\mdseries\slshape n}}]}. 
\item[{\texttt{NumberOfSolubleAutomorphisms := \mbox{\texttt{\mdseries\slshape n}}} \label{option NumberOfSolubleAutomorphisms}\index{option NumberOfSolubleAutomorphisms}}]  Specifies that the number of soluble automorphisms of the automorphism group
supplied by \texttt{PqPGSupplyAutomorphisms} (see{\nobreakspace}\texttt{PqPGSupplyAutomorphisms} (\ref{PqPGSupplyAutomorphisms})) in a $p$\texttt{\symbol{45}}group generation calculation is \mbox{\texttt{\mdseries\slshape n}}. By default, \mbox{\texttt{\mdseries\slshape n}} is taken to be $0$; \mbox{\texttt{\mdseries\slshape n}} must be a non\texttt{\symbol{45}}negative integer. If $\mbox{\texttt{\mdseries\slshape n}} \ge 0$ then a value for the option \texttt{RelativeOrders} (see{\nobreakspace}\ref{option RelativeOrders}) must also be supplied. 
\item[{\texttt{RelativeOrders := \mbox{\texttt{\mdseries\slshape list}}} \label{option RelativeOrders}\index{option RelativeOrders}}]  Specifies the relative orders of each soluble automorphism of the automorphism
group supplied by \texttt{PqPGSupplyAutomorphisms} (see{\nobreakspace}\texttt{PqPGSupplyAutomorphisms} (\ref{PqPGSupplyAutomorphisms})) in a $p$\texttt{\symbol{45}}group generation calculation. The list \mbox{\texttt{\mdseries\slshape list}} must consist of \mbox{\texttt{\mdseries\slshape n}} positive integers, where \mbox{\texttt{\mdseries\slshape n}} is the value of the option \texttt{NumberOfSolubleAutomorphisms} (see{\nobreakspace}\ref{option NumberOfSolubleAutomorphisms}). By default \mbox{\texttt{\mdseries\slshape list}} is empty. 
\item[{\texttt{BasicAlgorithm} \label{option BasicAlgorithm}\index{option BasicAlgorithm}}]  Specifies that an algorithm that the \texttt{pq} program calls its ``default'' algorithm be used for $p$\texttt{\symbol{45}}group generation. By default this algorithm is \emph{not} used. If this option is supplied the settings of options \texttt{RankInitialSegmentSubgroups}, \texttt{AllDescendants}, \texttt{Exponent} and \texttt{Metabelian} are ignored. 
\item[{\texttt{CustomiseOutput := \mbox{\texttt{\mdseries\slshape rec}}} \label{option CustomiseOutput}\index{option CustomiseOutput}}]  Specifies that fine tuning of the output is desired. The record \mbox{\texttt{\mdseries\slshape rec}} should have any subset (or all) of the the following fields: 
\begin{description}
\item[{\texttt{perm := \mbox{\texttt{\mdseries\slshape list}}}}]  where \mbox{\texttt{\mdseries\slshape list}} is a list of booleans which determine whether the permutation group output for
the automorphism group should contain: the degree, the extended automorphisms,
the automorphism matrices, and the permutations, respectively. 
\item[{\texttt{orbit := \mbox{\texttt{\mdseries\slshape list}}}}]  where \mbox{\texttt{\mdseries\slshape list}} is a list of booleans which determine whether the orbit output of the action
of the automorphism group should contain: a summary, and a complete listing of
orbits, respectively. (It's possible to have \emph{both} a summary and a complete listing.) 
\item[{\texttt{group := \mbox{\texttt{\mdseries\slshape list}}}}]  where \mbox{\texttt{\mdseries\slshape list}} is a list of booleans which determine whether the group output should contain:
the standard matrix of each allowable subgroup, the presentation of reduced $p$\texttt{\symbol{45}}covering groups, the presentation of immediate
descendants, the nuclear rank of descendants, and the $p$\texttt{\symbol{45}}multiplicator rank of descendants, respectively. 
\item[{\texttt{autgroup := \mbox{\texttt{\mdseries\slshape list}}}}]  where \mbox{\texttt{\mdseries\slshape list}} is a list of booleans which determine whether the automorphism group output
should contain: the commutator matrix, the automorphism group description of
descendants, and the automorphism group order of descendants, respectively. 
\item[{\texttt{trace := \mbox{\texttt{\mdseries\slshape val}}}}]  where \mbox{\texttt{\mdseries\slshape val}} is a boolean which if \texttt{true} specifies algorithm trace data is desired. By default, one does not get
algorithm trace data. 
\end{description}
 Not providing a field (or mis\texttt{\symbol{45}}spelling it!), specifies that
the default output is desired. As a convenience, \texttt{1} is also accepted as \texttt{true}, and any value that is neither \texttt{1} nor \texttt{true} is taken as \texttt{false}. Also for each \mbox{\texttt{\mdseries\slshape list}} above, an unbound list entry is taken as \texttt{false}. Thus, for example 
\begin{Verbatim}[commandchars=!@|,fontsize=\small,frame=single,label=Example]
  CustomiseOutput := rec(group := [,,1], autgroup := [,1])
\end{Verbatim}
 specifies for the group output that only the presentation of immediate
descendants is desired, for the automorphism group output only the
automorphism group description of descendants should be printed, that there
should be no algorithm trace data, and that the default output should be
provided for the permutation group and orbit output. 
\item[{\texttt{StandardPresentationFile := \mbox{\texttt{\mdseries\slshape filename}}} \label{option StandardPresentationFile}\index{option StandardPresentationFile}}]  Specifies that the file to which the standard presentation is written has name \mbox{\texttt{\mdseries\slshape filename}}. If the first character of the string \mbox{\texttt{\mdseries\slshape filename}} is not \texttt{/}, \mbox{\texttt{\mdseries\slshape filename}} is assumed to be the path of a writable file relative to the directory in
which \textsf{GAP} was started. If this option is omitted it is written to the file with the name
generated by the command \texttt{Filename( ANUPQData.tmpdir, "SPres" );}, i.e.{\nobreakspace}the file with name \texttt{"SPres"} in the temporary directory in which the \texttt{pq} program executes. 
\item[{\texttt{QueueFactor := \mbox{\texttt{\mdseries\slshape n}}} \label{option QueueFactor}\index{option QueueFactor}}]  Specifies a queue factor of \mbox{\texttt{\mdseries\slshape n}}, where \mbox{\texttt{\mdseries\slshape n}} must be a positive integer. This option may be used with \texttt{PqNextClass} (see{\nobreakspace}\texttt{PqNextClass} (\ref{PqNextClass})). 

 The queue factor is used when the \texttt{pq} program uses automorphisms to close a set of elements of the $p$\texttt{\symbol{45}}multiplicator under their action. 

 The algorithm used is a spinning algorithm: it starts with a set of vectors in
echelonized form (elements of the $p$\texttt{\symbol{45}}multiplicator) and closes the span of these vectors under
the action of the automorphisms. For this each automorphism is applied to each
vector and it is checked if the result is contained in the span computed so
far. If not, the span becomes bigger and the vector is put into a queue and
the automorphisms are applied to this vector at a later stage. The process
terminates when the automorphisms have been applied to all vectors and no new
vectors have been produced. 

 For each new vector it is decided, if its processing should be delayed. If the
vector contains too many non\texttt{\symbol{45}}zero entries, it is put into a
second queue. The elements in this queue are processed only when there are no
elements in the first queue left. 

 The queue factor is a percentage figure. A vector is put into the second queue
if the percentage of its non\texttt{\symbol{45}}zero entries exceeds the queue
factor. 
\item[{\texttt{Bounds := \mbox{\texttt{\mdseries\slshape list}}} \label{option Bounds}\index{option Bounds}}]  Specifies a lower and upper bound on the indices of a list, where \mbox{\texttt{\mdseries\slshape list}} is a pair of positive non\texttt{\symbol{45}}decreasing integers.
See{\nobreakspace}\texttt{PqDisplayStructure} (\ref{PqDisplayStructure}) and{\nobreakspace}\texttt{PqDisplayAutomorphisms} (\ref{PqDisplayAutomorphisms}) where this option may be used. 
\item[{\texttt{PrintAutomorphisms := \mbox{\texttt{\mdseries\slshape list}}} \label{option PrintAutomorphisms}\index{option PrintAutomorphisms}}]  Specifies that automorphism matrices be printed. 
\item[{\texttt{PrintPermutations := \mbox{\texttt{\mdseries\slshape list}}} \label{option PrintPermutations}\index{option PrintPermutations}}]  Specifies that permutations of the subgroups be printed. 
\item[{\texttt{Filename := \mbox{\texttt{\mdseries\slshape string}}} \label{option Filename}\index{option Filename}}]  Specifies that an output or input file to be written to or read from by the \texttt{pq} program should have the name \mbox{\texttt{\mdseries\slshape string}}. 
\end{description}
 }

 }

 
\chapter{\textcolor{Chapter }{Installing the ANUPQ Package}}\label{Installing-ANUPQ}
\logpage{[ 7, 0, 0 ]}
\hyperdef{L}{X7FB487238298CF42}{}
{
  The ANU \texttt{pq} program is written in C and the package can be installed under UNIX and in
environments similar to UNIX. In particular it is known to work on Linux and
Mac OS X, and also on Windows equipped with cygwin. 

 The current version of the \textsf{ANUPQ} package requires \textsf{GAP}{\nobreakspace}4.9, and version 1.2 of the \textsf{AutPGrp} package. However, we recommend using at least \textsf{GAP}{\nobreakspace}4.6 and \textsf{AutPGrp}{\nobreakspace}1.5. 

 To install the \textsf{ANUPQ} package, move the file \texttt{anupq\texttt{\symbol{45}}\mbox{\texttt{\mdseries\slshape XXX}}.tar.gz} for some version number \mbox{\texttt{\mdseries\slshape XXX}} into the \texttt{pkg} directory in which you plan to install \textsf{ANUPQ}. Usually, this will be the directory \texttt{pkg} in the hierarchy of your version of \textsf{GAP}; it is however also possible to keep an additional \texttt{pkg} directory in your private directories. The only essential difference with
installing \textsf{ANUPQ} in a \texttt{pkg} directory different to the \textsf{GAP} home directory is that one must start \textsf{GAP} with the \texttt{\texttt{\symbol{45}}l} switch (see Section{\nobreakspace} \textbf{Reference: Command Line Options}), e.g.{\nobreakspace}if your private \texttt{pkg} directory is a subdirectory of \texttt{mygap} in your home directory you might type: 
\begin{Verbatim}[commandchars=!@|,fontsize=\small,frame=single,label=]
  gap -l ";myhomedir/mygap"
\end{Verbatim}
 where \mbox{\texttt{\mdseries\slshape myhomedir}} is the path to your home directory, which may be replaced by a tilde. The
empty path before the semicolon is filled in by the default path of the \textsf{GAP} home directory. 

 Then, in your chosen \texttt{pkg} directory, unpack \texttt{anupq\texttt{\symbol{45}}\mbox{\texttt{\mdseries\slshape XXX}}.tar.gz} by 
\begin{Verbatim}[commandchars=!@|,fontsize=\small,frame=single,label=]
  tar xf anupq-<XXX>.tar.gz
\end{Verbatim}
 Change to the newly created \texttt{anupq} directory. Now you need to call \texttt{configure}. If you installed \textsf{ANUPQ} into the main \texttt{pkg} directory, call 
\begin{Verbatim}[commandchars=!@|,fontsize=\small,frame=single,label=]
  ./configure
\end{Verbatim}
 If you installed ANUPQ in another directory than the usual 'pkg' subdirectory,
instead call 
\begin{Verbatim}[commandchars=!@|,fontsize=\small,frame=single,label=]
  ./configure --with-gaproot=<path>
\end{Verbatim}
 where \mbox{\texttt{\mdseries\slshape path}} is the path to the \textsf{GAP} home directory. (You can also call 
\begin{Verbatim}[commandchars=!@|,fontsize=\small,frame=single,label=]
  ./configure --help
\end{Verbatim}
 for further options.) 

 What this does is look for a file \texttt{sysinfo.gap} in the root directory of \textsf{GAP} in order to determine an architecture name for the subdirectory of \texttt{bin} in which to put the compiled \texttt{pq} binary. This only makes sense if \textsf{GAP} was compiled for the same architecture that \texttt{pq} will be. If you have a shared file system mounted across different
architectures, then you should run \texttt{configure} and \texttt{make} for \textsf{ANUPQ} for each architecture immediately after compiling \textsf{GAP} on the same architecture. 

 If you had to install the package in your own directory but wish to use the
system \textsf{GAP} then you will need to find out what \mbox{\texttt{\mdseries\slshape path}} is. To do this, start up \textsf{GAP} and find out what \textsf{GAP}'s root path is from finding the value of the variable \texttt{GAPInfo.RootPaths}, e.g. 
\begin{Verbatim}[commandchars=!@|,fontsize=\small,frame=single,label=Example]
  !gapprompt@gap>| !gapinput@GAPInfo.RootPaths;|
  [ "/usr/local/lib/gap4r4/" ]
\end{Verbatim}
 would tell you to use \texttt{/usr/local/lib/gap4r4} for \mbox{\texttt{\mdseries\slshape path}}. 

 The \texttt{configure} command will fetch the architecture type for which \textsf{GAP} has been compiled last and create a \texttt{Makefile}. You can now simply call 
\begin{Verbatim}[commandchars=!@|,fontsize=\small,frame=single,label=]
  make
\end{Verbatim}
 to compile the binary and to install it in the appropriate place. 

 \index{ANUPQ_GAP_EXEC@\texttt{ANUPQ{\textunderscore}GAP{\textunderscore}EXEC}!environment variable} The path of \textsf{GAP} (see \emph{Note} below) used by the \texttt{pq} binary (the value \texttt{GAP} is set to in the \texttt{make} command) may be over\texttt{\symbol{45}}ridden by setting the environment
variable \texttt{ANUPQ{\textunderscore}GAP{\textunderscore}EXEC}. These values are only of interest when the \texttt{pq} program is run as a standalone; however, the \texttt{testPq} script assumes you have set one of these correctly (see Section{\nobreakspace}\hyperref[Testing your ANUPQ installation]{`Testing your ANUPQ installation'}). When the \texttt{pq} program is started from \textsf{GAP} communication occurs via an iostream, so that the \texttt{pq} binary does not actually need to know a valid path for \textsf{GAP} is this case. 

 \emph{Note.} By ``path of \textsf{GAP}'' we mean the path of the command used to invoke \textsf{GAP} (which should be a script, e.g. the \texttt{gap.sh} script generated in the \texttt{bin} directory for the version of \textsf{GAP} when \textsf{GAP} was compiled). The usual strategy is to copy the \texttt{gap.sh} script to a standard location, e.g. \texttt{/usr/local/bin/gap}. It is a mistake to copy the \textsf{GAP} executable \texttt{gap} (in a directory with name of form \texttt{bin/\mbox{\texttt{\mdseries\slshape compile\texttt{\symbol{45}}platform}}}) to the standard location, since direct invocation of the executable results
in \textsf{GAP} starting without being able to find its own library (a fatal error). 
\section{\textcolor{Chapter }{Testing your ANUPQ installation}}\label{Testing your ANUPQ installation}
\logpage{[ 7, 1, 0 ]}
\hyperdef{L}{X854577C8800DC7C2}{}
{
  \index{ANUPQ_GAP_EXEC@\texttt{ANUPQ{\textunderscore}GAP{\textunderscore}EXEC}!environment variable} Now it is time to test the installation. After doing \texttt{configure} and \texttt{make} you will have a \texttt{testPq} script. The script assumes that, if the environment variable \texttt{ANUPQ{\textunderscore}GAP{\textunderscore}EXEC} is set, it is a correct path for \textsf{GAP}, or otherwise that the \texttt{make} call that compiled the \texttt{pq} program set \texttt{GAP} to a correct path for \textsf{GAP} (see Section{\nobreakspace}\hyperref[Running the pq program as a standalone]{`Running the pq program as a standalone'} for more details). To run the tests, just type: 
\begin{Verbatim}[commandchars=!@|,fontsize=\small,frame=single,label=]
  ./testPq
\end{Verbatim}
 Some of the tests the script runs take a while. Please be patient. The script
checks that you not only have a correct \textsf{GAP} (at least version 4.4) installation that includes the \textsf{AutPGrp} package, but that the \textsf{ANUPQ} package and its \texttt{pq} binary interact correctly. You should see something like the following output: 
\begin{Verbatim}[commandchars=@|B,fontsize=\small,frame=single,label=Example]
  Made dir: /tmp/testPq
  Testing installation of ANUPQ Package (version 3.1)
   
  The first two tests check that the pq C program compiled ok.
  Testing the pq binary ... OK.
  Testing the pq binary's stack size ... OK.
  The pq C program compiled ok! We test it's the right one below.
   
  The next tests check that you have the right version of GAP
  for the ANUPQ package and that GAP is finding
  the right versions of the ANUPQ and AutPGrp packages.
   
  Checking GAP ...
   pq binary made with GAP set to: /usr/local/bin/gap
   Starting GAP to determine version and package availability ...
    GAP version (4.6.5) ... OK.
    GAP found ANUPQ package (version 3.1) ... good.
    GAP found pq binary (version 1.9) ... good.
    GAP found AutPGrp package (version 1.5) ... good.
   GAP is OK.
   
  Checking the link between the pq binary and GAP ... OK.
  Testing the standard presentation part of the pq binary ... OK.
  Doing p-group generation (final GAP/ANUPQ) test ... OK.
  Tests complete.
  Removed dir: /tmp/testPq
  Enjoy using your functional ANUPQ package!
\end{Verbatim}
 }

 
\section{\textcolor{Chapter }{Running the pq program as a standalone}}\label{Running the pq program as a standalone}
\logpage{[ 7, 2, 0 ]}
\hyperdef{L}{X82906B5184C61063}{}
{
  \index{ANUPQ_GAP_EXEC@\texttt{ANUPQ{\textunderscore}GAP{\textunderscore}EXEC}!environment variable} When the \texttt{pq} program is run as a standalone it sometimes needs to call \textsf{GAP} to compute stabilisers of subgroups; in doing so, it first checks the value of
the environment variable \texttt{ANUPQ{\textunderscore}GAP{\textunderscore}EXEC}, and uses that, if set, or otherwise the value of \texttt{GAP} it was compiled with, as the path for \textsf{GAP}. If you ran \texttt{testPq} (see Section{\nobreakspace}\hyperref[Testing your ANUPQ installation]{`Testing your ANUPQ installation'}) and you got both \textsf{GAP} is \texttt{OK} and the link between the \texttt{pq} binary and \textsf{GAP} is \texttt{OK}, you should be fine. Otherwise heed the recommendations of the error messages
you get and run the \texttt{testPq} until all tests are passed. 

 It is especially important that the \textsf{GAP}, whose path you gave, should know where to find the \textsf{ANUPQ} and \textsf{AutPGrp} packages. To ensure this the path should be to a shell script that invokes \textsf{GAP}. If you needed to install the needed packages in your own directory (because,
say, you are not a system administrator) then you should create your own shell
script that runs \textsf{GAP} with a correct setting of the \texttt{\texttt{\symbol{45}}l} option and set the path used by the \texttt{pq} binary to the path of that script. To create the script that runs \textsf{GAP} it is easiest to copy the system one and edit it, e.g. start by executing the
following UNIX commands (skip the second step if you already have a \texttt{bin} directory; \texttt{you@unix{\textgreater}} is your UNIX prompt): 
\begin{Verbatim}[commandchars=!|A,fontsize=\small,frame=single,label=]
  you@unix> cd
  you@unix> mkdir bin
  you@unix> cd bin
  you@unix> which gap
  /usr/local/bin/gap
  you@unix> cp /usr/local/bin/gap mygap
  you@unix> chmod +x mygap
\end{Verbatim}
 At the second\texttt{\symbol{45}}last step use the path of \textsf{GAP} returned by \texttt{which gap}. Now hopefully you will have a copy of the script that runs the system \textsf{GAP} in \texttt{mygap}. Now use your favourite editor to edit the \texttt{\texttt{\symbol{45}}l} part of the last line of \texttt{mygap} which should initially look something like: 
\begin{Verbatim}[commandchars=!@|,fontsize=\small,frame=single,label=]
  exec $GAP_DIR/bin/$GAP_PRG -m $GAP_MEM -o 970m -l $GAP_DIR $*
\end{Verbatim}
 so that it becomes (the tilde is a UNIX abbreviation for your home directory): 
\begin{Verbatim}[commandchars=!@|,fontsize=\small,frame=single,label=]
  exec $GAP_DIR/bin/$GAP_PRG -m $GAP_MEM -o 970m -l "$GAP_DIR;~/gapstuff" $*
\end{Verbatim}
 assuming that your personal \textsf{GAP} \texttt{pkg} directory is a subdirectory of \texttt{gapstuff} in your home directory. Finally, to let the \texttt{pq} program know where \textsf{GAP} is and also know where your \texttt{pkg} directory is that contains \textsf{ANUPQ}, set the environment variable \texttt{ANUPQ{\textunderscore}GAP{\textunderscore}EXEC} to the complete (i.e. absolute) path of your \texttt{mygap} script (do not use the tilde abbreviation). }

 }

 

\appendix


\chapter{\textcolor{Chapter }{Examples}}\label{Examples}
\logpage{[ "A", 0, 0 ]}
\hyperdef{L}{X7A489A5D79DA9E5C}{}
{
  There are a large number of examples provided with the \textsf{ANUPQ} package. These may be executed or displayed via the function \texttt{PqExample} (see{\nobreakspace}\texttt{PqExample} (\ref{PqExample})). Each example resides in a file of the same name in the directory \texttt{examples}. Most of the examples are translations to \textsf{GAP} of examples provided for the \texttt{pq} standalone by Eamonn O'Brien; the standalone examples are found in directories \texttt{standalone/examples} ($p$\texttt{\symbol{45}}quotient and $p$\texttt{\symbol{45}}group generation examples) and \texttt{standalone/isom} (standard presentation examples). The first line of each example indicates its
origin. All the examples seen in earlier chapters of this manual are also
available as examples, in a slightly modified form (the example which one can
run in order to see something very close to the text example ``live'' is always indicated near \texttt{\symbol{45}}\texttt{\symbol{45}} usually
immediately after \texttt{\symbol{45}}\texttt{\symbol{45}} the text example).
The format of the (\texttt{PqExample}) examples is such that they can be read by the standard \texttt{Read} function of \textsf{GAP}, but certain features and comments are interpreted by the function \texttt{PqExample} to do somewhat more than \texttt{Read} does. In particular, any function without a \texttt{\texttt{\symbol{45}}i}, \texttt{\texttt{\symbol{45}}ni} or \texttt{.g} suffix has both a non\texttt{\symbol{45}}interactive and interactive form; in
these cases, the default form is the non\texttt{\symbol{45}}interactive form,
and giving \texttt{PqStart} as second argument generates the interactive form. 

 Running \texttt{PqExample} without an argument or with a non\texttt{\symbol{45}}existent example \texttt{Info}s the available examples and some hints on usage: 
\begin{Verbatim}[commandchars=!@|,fontsize=\small,frame=single,label=Example]
  !gapprompt@gap>| !gapinput@PqExample();|
  #I                   PqExample Index (Table of Contents)
  #I                   -----------------------------------
  #I  This table of possible examples is displayed when calling `PqExample'
  #I  with no arguments, or with the argument: "index" (meant in the  sense
  #I  of ``list''), or with a non-existent example name.
  #I  
  #I  Examples that have a name ending in `-ni' are  non-interactive  only.
  #I  Examples that have a  name  ending  in  `-i'  are  interactive  only.
  #I  Examples with names ending in `.g' also have  only  one  form.  Other
  #I  examples have both a non-interactive and an  interactive  form;  call
  #I  `PqExample' with 2nd argument `PqStart' to get the  interactive  form
  #I  of the example. The substring `PG' in an  example  name  indicates  a
  #I  p-Group Generation example, `SP' indicates  a  Standard  Presentation
  #I  example, `Rel' indicates it uses  the  `Relators'  option,  and  `Id'
  #I  indicates it uses the `Identities' option.
  #I  
  #I  The following ANUPQ examples are available:
  #I  
  #I   p-Quotient examples:
  #I    general:
  #I     "Pq"                   "Pq-ni"                "PqEpimorphism"        
  #I     "PqPCover"             "PqSupplementInnerAutomorphisms"
  #I    2-groups:
  #I     "2gp-Rel"              "2gp-Rel-i"            "2gp-a-Rel-i"
  #I     "B2-4"                 "B2-4-Id"              "B2-8-i"
  #I     "B4-4-i"               "B4-4-a-i"             "B5-4.g"
  #I    3-groups:
  #I     "3gp-Rel-i"            "3gp-a-Rel"            "3gp-a-Rel-i"
  #I     "3gp-a-x-Rel-i"        "3gp-maxoccur-Rel-i"
  #I    5-groups:
  #I     "5gp-Rel-i"            "5gp-a-Rel-i"          "5gp-b-Rel-i"
  #I     "5gp-c-Rel-i"          "5gp-metabelian-Rel-i" "5gp-maxoccur-Rel-i"
  #I     "F2-5-i"               "B2-5-i"               "R2-5-i"
  #I     "R2-5-x-i"             "B5-5-Engel3-Id"
  #I    7-groups:
  #I     "7gp-Rel-i"
  #I    11-groups:
  #I     "11gp-i"               "11gp-Rel-i"           "11gp-a-Rel-i"
  #I     "11gp-3-Engel-Id"      "11gp-3-Engel-Id-i"
  #I  
  #I   p-Group Generation examples:
  #I    general:
  #I     "PqDescendants-1"      "PqDescendants-2"      "PqDescendants-3"
  #I     "PqDescendants-1-i"
  #I    2-groups:
  #I     "2gp-PG-i"             "2gp-PG-2-i"           "2gp-PG-3-i"
  #I     "2gp-PG-4-i"           "2gp-PG-e4-i"
  #I     "PqDescendantsTreeCoclassOne-16-i"
  #I    3-groups:
  #I     "3gp-PG-i"             "3gp-PG-4-i"           "3gp-PG-x-i"
  #I     "3gp-PG-x-1-i"         "PqDescendants-treetraverse-i"
  #I     "PqDescendantsTreeCoclassOne-9-i"
  #I    5-groups:
  #I     "5gp-PG-i"             "Nott-PG-Rel-i"        "Nott-APG-Rel-i"
  #I     "PqDescendantsTreeCoclassOne-25-i"
  #I    7,11-groups:
  #I     "7gp-PG-i"             "11gp-PG-i"
  #I  
  #I   Standard Presentation examples:
  #I    general:
  #I     "StandardPresentation" "StandardPresentation-i"
  #I     "EpimorphismStandardPresentation"
  #I     "EpimorphismStandardPresentation-i"           "IsIsomorphicPGroup-ni"
  #I    2-groups:
  #I     "2gp-SP-Rel-i"         "2gp-SP-1-Rel-i"       "2gp-SP-2-Rel-i"
  #I     "2gp-SP-3-Rel-i"       "2gp-SP-4-Rel-i"       "2gp-SP-d-Rel-i"
  #I     "gp-256-SP-Rel-i"      "B2-4-SP-i"            "G2-SP-Rel-i"
  #I    3-groups:
  #I     "3gp-SP-Rel-i"         "3gp-SP-1-Rel-i"       "3gp-SP-2-Rel-i"
  #I     "3gp-SP-3-Rel-i"       "3gp-SP-4-Rel-i"       "G3-SP-Rel-i"
  #I    5-groups:
  #I     "5gp-SP-Rel-i"         "5gp-SP-a-Rel-i"       "5gp-SP-b-Rel-i"
  #I     "5gp-SP-big-Rel-i"     "5gp-SP-d-Rel-i"       "G5-SP-Rel-i"
  #I     "G5-SP-a-Rel-i"        "Nott-SP-Rel-i"
  #I    7-groups:
  #I     "7gp-SP-Rel-i"         "7gp-SP-a-Rel-i"       "7gp-SP-b-Rel-i"
  #I    11-groups:
  #I     "11gp-SP-a-i"          "11gp-SP-a-Rel-i"      "11gp-SP-a-Rel-1-i"
  #I     "11gp-SP-b-i"          "11gp-SP-b-Rel-i"      "11gp-SP-c-Rel-i"
  #I  
  #I  Notes
  #I  -----
  #I  1. The example (first) argument of  `PqExample'  is  a  string;  each
  #I     example above is in double quotes to remind you to include them.
  #I  2. Some examples accept options. To find  out  whether  a  particular
  #I     example accepts options, display it first (by including  `Display'
  #I     as  last  argument)  which  will  also  indicate  how  `PqExample'
  #I     interprets the options, e.g. `PqExample("11gp-SP-a-i", Display);'.
  #I  3. Try `SetInfoLevel(InfoANUPQ, <n>);' for  some  <n>  in  [2  ..  4]
  #I     before calling PqExample, to see what's going on behind the scenes.
  #I  
\end{Verbatim}
 If on your terminal you are unable to scroll back, an alternative to typing \texttt{PqExample();} to see the displayed examples is to use on\texttt{\symbol{45}}line help,
i.e.{\nobreakspace} you may type: 
\begin{Verbatim}[commandchars=!@|,fontsize=\small,frame=single,label=Example]
  !gapprompt@gap>| !gapinput@?anupq:examples|
\end{Verbatim}
 which will display this appendix in a \textsf{GAP} session. If you are not fussed about the order in which the examples are
organised, \texttt{AllPqExamples();} lists the available examples relatively compactly (see{\nobreakspace}\texttt{AllPqExamples} (\ref{AllPqExamples})). 

 In the remainder of this appendix we will discuss particular aspects related
to the \texttt{Relators} (see{\nobreakspace}\ref{option Relators}) and \texttt{Identities} (see{\nobreakspace}\ref{option Identities}) options, and the construction of the Burnside group $B(5, 4)$. 
\section{\textcolor{Chapter }{The Relators Option}}\label{The Relators Option}
\logpage{[ "A", 1, 0 ]}
\hyperdef{L}{X80676E317BF4DF13}{}
{
  The \texttt{Relators} option was included because computations involving words containing
commutators that are pre\texttt{\symbol{45}}expanded by \textsf{GAP} before being passed to the \texttt{pq} program may run considerably more slowly, than the same computations being run
with \textsf{GAP} pre\texttt{\symbol{45}}expansions avoided. The following examples demonstrate
a case where the performance hit due to pre\texttt{\symbol{45}}expansion of
commutators by \textsf{GAP} is a factor of order 100 (in order to see timing information from the \texttt{pq} program, we set the \texttt{InfoANUPQ} level to 2). 

 Firstly, we run the example that allows pre\texttt{\symbol{45}}expansion of
commutators (the function \texttt{PqLeftNormComm} is provided by the \textsf{ANUPQ} package; see{\nobreakspace}\texttt{PqLeftNormComm} (\ref{PqLeftNormComm})). Note that since the two commutators of this example are \emph{very} long (taking more than an page to print), we have edited the output at this
point. 
\begin{Verbatim}[commandchars=!@|,fontsize=\small,frame=single,label=Example]
  !gapprompt@gap>| !gapinput@SetInfoLevel(InfoANUPQ, 2); #to see timing information|
  !gapprompt@gap>| !gapinput@PqExample("11gp-i");|
  #I  #Example: "11gp-i" . . . based on: examples/11gp
  #I  F, a, b, c, R, procId are local to `PqExample'
  !gapprompt@gap>| !gapinput@F := FreeGroup("a", "b", "c"); a := F.1; b := F.2; c := F.3;|
  <free group on the generators [ a, b, c ]>
  a
  b
  c
  !gapprompt@gap>| !gapinput@R := [PqLeftNormComm([b, a, a, b, c])^11, |
  !gapprompt@>| !gapinput@         PqLeftNormComm([a, b, b, a, b, c])^11, (a * b)^11];;|
  !gapprompt@gap>| !gapinput@procId := PqStart(F/R : Prime := 11);;|
  !gapprompt@gap>| !gapinput@PqPcPresentation(procId : ClassBound := 7, |
  !gapprompt@>| !gapinput@                             OutputLevel := 1);|
  #I  Lower exponent-11 central series for [grp]
  #I  Group: [grp] to lower exponent-11 central class 1 has order 11^3
  #I  Group: [grp] to lower exponent-11 central class 2 has order 11^8
  #I  Group: [grp] to lower exponent-11 central class 3 has order 11^19
  #I  Group: [grp] to lower exponent-11 central class 4 has order 11^42
  #I  Group: [grp] to lower exponent-11 central class 5 has order 11^98
  #I  Group: [grp] to lower exponent-11 central class 6 has order 11^228
  #I  Group: [grp] to lower exponent-11 central class 7 has order 11^563
  #I  Computation of presentation took 27.04 seconds
  !gapprompt@gap>| !gapinput@PqSavePcPresentation(procId, ANUPQData.outfile);|
  #I  Variables used in `PqExample' are saved in `ANUPQData.example.vars'.
\end{Verbatim}
 Now we do the same calculation using the \texttt{Relators} option. In this way, the commutators are passed directly as strings to the \texttt{pq} program, so that \textsf{GAP} does not ``see'' them and pre\texttt{\symbol{45}}expand them. 
\begin{Verbatim}[commandchars=!@|,fontsize=\small,frame=single,label=Example]
  !gapprompt@gap>| !gapinput@PqExample("11gp-Rel-i");|
  #I  #Example: "11gp-Rel-i" . . . based on: examples/11gp
  #I  #(equivalent to "11gp-i" example but uses `Relators' option)
  #I  F, rels, procId are local to `PqExample'
  !gapprompt@gap>| !gapinput@F := FreeGroup("a", "b", "c");|
  <free group on the generators [ a, b, c ]>
  !gapprompt@gap>| !gapinput@rels := ["[b, a, a, b, c]^11", "[a, b, b, a, b, c]^11", "(a * b)^11"];|
  [ "[b, a, a, b, c]^11", "[a, b, b, a, b, c]^11", "(a * b)^11" ]
  !gapprompt@gap>| !gapinput@procId := PqStart(F : Prime := 11, Relators := rels);;|
  !gapprompt@gap>| !gapinput@PqPcPresentation(procId : ClassBound := 7, |
  !gapprompt@>| !gapinput@                             OutputLevel := 1);|
  #I  Relators parsed ok.
  #I  Lower exponent-11 central series for [grp]
  #I  Group: [grp] to lower exponent-11 central class 1 has order 11^3
  #I  Group: [grp] to lower exponent-11 central class 2 has order 11^8
  #I  Group: [grp] to lower exponent-11 central class 3 has order 11^19
  #I  Group: [grp] to lower exponent-11 central class 4 has order 11^42
  #I  Group: [grp] to lower exponent-11 central class 5 has order 11^98
  #I  Group: [grp] to lower exponent-11 central class 6 has order 11^228
  #I  Group: [grp] to lower exponent-11 central class 7 has order 11^563
  #I  Computation of presentation took 0.27 seconds
  !gapprompt@gap>| !gapinput@PqSavePcPresentation(procId, ANUPQData.outfile);|
  #I  Variables used in `PqExample' are saved in `ANUPQData.example.vars'.
\end{Verbatim}
 }

 
\section{\textcolor{Chapter }{The Identities Option and PqEvaluateIdentities Function}}\label{The Identities Option and PqEvaluateIdentities Function}
\logpage{[ "A", 2, 0 ]}
\hyperdef{L}{X80B2AA7084C57F5D}{}
{
  Please pay heed to the warnings given for the \texttt{Identities} option (see{\nobreakspace}\ref{option Identities}); it is written mainly at the \textsf{GAP} level and is not particularly optimised. The \texttt{Identities} option allows one to compute $p$\texttt{\symbol{45}}quotients that satisfy an identity. A trivial example
better done using the \texttt{Exponent} option, but which nevertheless demonstrates the usage of the \texttt{Identities} option, is as follows: 
\begin{Verbatim}[commandchars=!@|,fontsize=\small,frame=single,label=Example]
  !gapprompt@gap>| !gapinput@SetInfoLevel(InfoANUPQ, 1);|
  !gapprompt@gap>| !gapinput@PqExample("B2-4-Id");|
  #I  #Example: "B2-4-Id" . . . alternative way to generate B(2, 4)
  #I  #Generates B(2, 4) by using the `Identities' option
  #I  #... this is not as efficient as using `Exponent' but
  #I  #demonstrates the usage of the `Identities' option.
  #I  F, f, procId are local to `PqExample'
  !gapprompt@gap>| !gapinput@F := FreeGroup("a", "b");|
  <free group on the generators [ a, b ]>
  !gapprompt@gap>| !gapinput@# All words w in the pc generators of B(2, 4) satisfy f(w) = 1 |
  !gapprompt@gap>| !gapinput@f := w -> w^4;|
  function( w ) ... end
  !gapprompt@gap>| !gapinput@Pq( F : Prime := 2, Identities := [ f ] );|
  #I  Class 1 with 2 generators.
  #I  Class 2 with 5 generators.
  #I  Class 3 with 7 generators.
  #I  Class 4 with 10 generators.
  #I  Class 5 with 12 generators.
  #I  Class 5 with 12 generators.
  <pc group of size 4096 with 12 generators>
  #I  Variables used in `PqExample' are saved in `ANUPQData.example.vars'.
  !gapprompt@gap>| !gapinput@time; |
  1400
\end{Verbatim}
 Note that the \texttt{time} statement gives the time in milliseconds spent by \textsf{GAP} in executing the \texttt{PqExample("B2\texttt{\symbol{45}}4\texttt{\symbol{45}}Id");} command (i.e.{\nobreakspace}everything up to the \texttt{Info}\texttt{\symbol{45}}ing of the variables used), but over 90\% of that time is
spent in the final \texttt{Pq} statement. The time spent by the \texttt{pq} program, which is negligible anyway (you can check this by running the example
while the \texttt{InfoANUPQ} level is set to 2), is not counted by \texttt{time}. 

 Since the identity used in the above construction of $B(2, 4)$ is just an exponent law, the ``right'' way to compute it is via the \texttt{Exponent} option (see{\nobreakspace}\ref{option Exponent}), which is implemented at the C level and \emph{is} highly optimised. Consequently, the \texttt{Exponent} option is significantly faster, generally by several orders of magnitude: 
\begin{Verbatim}[commandchars=!@|,fontsize=\small,frame=single,label=Example]
  !gapprompt@gap>| !gapinput@SetInfoLevel(InfoANUPQ, 2); # to see time spent by the `pq' program|
  !gapprompt@gap>| !gapinput@PqExample("B2-4");|
  #I  #Example: "B2-4" . . . the ``right'' way to generate B(2, 4)
  #I  #Generates B(2, 4) by using the `Exponent' option
  #I  F, procId are local to `PqExample'
  !gapprompt@gap>| !gapinput@F := FreeGroup("a", "b");|
  <free group on the generators [ a, b ]>
  !gapprompt@gap>| !gapinput@Pq( F : Prime := 2, Exponent := 4 );|
  #I  Computation of presentation took 0.00 seconds
  <pc group of size 4096 with 12 generators>
  #I  Variables used in `PqExample' are saved in `ANUPQData.example.vars'.
  !gapprompt@gap>| !gapinput@time; # time spent by GAP in executing `PqExample("B2-4");' |
  50
\end{Verbatim}
 The following example uses the \texttt{Identities} option to compute a 3\texttt{\symbol{45}}Engel group for the prime 11. As is
the case for the example \texttt{"B2\texttt{\symbol{45}}4\texttt{\symbol{45}}Id"}, the example has both a non\texttt{\symbol{45}}interactive and an interactive
form; below, we demonstrate the interactive form. 
\begin{Verbatim}[commandchars=!@|,fontsize=\small,frame=single,label=Example]
  !gapprompt@gap>| !gapinput@SetInfoLevel(InfoANUPQ, 1); # reset InfoANUPQ to default level|
  !gapprompt@gap>| !gapinput@PqExample("11gp-3-Engel-Id", PqStart);|
  #I  #Example: "11gp-3-Engel-Id" . . . 3-Engel group for prime 11
  #I  #Non-trivial example of using the `Identities' option
  #I  F, a, b, G, f, procId, Q are local to `PqExample'
  !gapprompt@gap>| !gapinput@F := FreeGroup("a", "b"); a := F.1; b := F.2;|
  <free group on the generators [ a, b ]>
  a
  b
  !gapprompt@gap>| !gapinput@G := F/[ a^11, b^11 ];|
  <fp group on the generators [ a, b ]>
  !gapprompt@gap>| !gapinput@# All word pairs u, v in the pc generators of the 11-quotient Q of G |
  !gapprompt@gap>| !gapinput@# must satisfy the Engel identity: [u, v, v, v] = 1.|
  !gapprompt@gap>| !gapinput@f := function(u, v) return PqLeftNormComm( [u, v, v, v] ); end;|
  function( u, v ) ... end
  !gapprompt@gap>| !gapinput@procId := PqStart( G );;|
  !gapprompt@gap>| !gapinput@Q := Pq( procId : Prime := 11, Identities := [ f ] );|
  #I  Class 1 with 2 generators.
  #I  Class 2 with 3 generators.
  #I  Class 3 with 5 generators.
  #I  Class 3 with 5 generators.
  <pc group of size 161051 with 5 generators>
  !gapprompt@gap>| !gapinput@# We do a ``sample'' check that pairs of elements of Q do satisfy|
  !gapprompt@gap>| !gapinput@# the given identity:|
  !gapprompt@gap>| !gapinput@f( Random(Q), Random(Q) );|
  <identity> of ...
  !gapprompt@gap>| !gapinput@f( Q.1, Q.2 );|
  <identity> of ...
  #I  Variables used in `PqExample' are saved in `ANUPQData.example.vars'.
\end{Verbatim}
 The (interactive) call to \texttt{Pq} above is essentially equivalent to a call to \texttt{PqPcPresentation} with the same arguments and options followed by a call to \texttt{PqCurrentGroup}. Moreover, the call to \texttt{PqPcPresentation} (as described in{\nobreakspace}\texttt{PqPcPresentation} (\ref{PqPcPresentation})) is equivalent to a ``class 1'' call to \texttt{PqPcPresentation} followed by the requisite number of calls to \texttt{PqNextClass}, and with the \texttt{Identities} option set, both \texttt{PqPcPresentation} and \texttt{PqNextClass} ``quietly'' perform the equivalent of a \texttt{PqEvaluateIdentities} call. In the following example we break down the \texttt{Pq} call into its low\texttt{\symbol{45}}level equivalents, and set and unset the \texttt{Identities} option to show where \texttt{PqEvaluateIdentities} fits into this scheme. 
\begin{Verbatim}[commandchars=!@|,fontsize=\small,frame=single,label=Example]
  !gapprompt@gap>| !gapinput@PqExample("11gp-3-Engel-Id-i");|
  #I  #Example: "11gp-3-Engel-Id-i" . . . 3-Engel grp for prime 11
  #I  #Variation of "11gp-3-Engel-Id" broken down into its lower-level component
  #I  #command parts.
  #I  F, a, b, G, f, procId, Q are local to `PqExample'
  !gapprompt@gap>| !gapinput@F := FreeGroup("a", "b"); a := F.1; b := F.2;|
  <free group on the generators [ a, b ]>
  a
  b
  !gapprompt@gap>| !gapinput@G := F/[ a^11, b^11 ];|
  <fp group on the generators [ a, b ]>
  !gapprompt@gap>| !gapinput@# All word pairs u, v in the pc generators of the 11-quotient Q of G |
  !gapprompt@gap>| !gapinput@# must satisfy the Engel identity: [u, v, v, v] = 1.|
  !gapprompt@gap>| !gapinput@f := function(u, v) return PqLeftNormComm( [u, v, v, v] ); end;|
  function( u, v ) ... end
  !gapprompt@gap>| !gapinput@procId := PqStart( G : Prime := 11 );;|
  !gapprompt@gap>| !gapinput@PqPcPresentation( procId : ClassBound := 1);|
  !gapprompt@gap>| !gapinput@PqEvaluateIdentities( procId : Identities := [f] );|
  #I  Class 1 with 2 generators.
  !gapprompt@gap>| !gapinput@for c in [2 .. 4] do|
  !gapprompt@>| !gapinput@     PqNextClass( procId : Identities := [] ); #reset `Identities' option|
  !gapprompt@>| !gapinput@     PqEvaluateIdentities( procId : Identities := [f] );|
  !gapprompt@>| !gapinput@   od;|
  #I  Class 2 with 3 generators.
  #I  Class 3 with 5 generators.
  #I  Class 3 with 5 generators.
  !gapprompt@gap>| !gapinput@Q := PqCurrentGroup( procId );|
  <pc group of size 161051 with 5 generators>
  !gapprompt@gap>| !gapinput@# We do a ``sample'' check that pairs of elements of Q do satisfy|
  !gapprompt@gap>| !gapinput@# the given identity:|
  !gapprompt@gap>| !gapinput@f( Random(Q), Random(Q) );|
  <identity> of ...
  !gapprompt@gap>| !gapinput@f( Q.1, Q.2 );|
  <identity> of ...
  #I  Variables used in `PqExample' are saved in `ANUPQData.example.vars'.
\end{Verbatim}
 }

 
\section{\textcolor{Chapter }{A Large Example}}\label{A Large Example}
\logpage{[ "A", 3, 0 ]}
\hyperdef{L}{X7A997D1A8532A84B}{}
{
  \index{B(5,4)@$B(5,4)$} An example demonstrating how a large computation can be organised with the \textsf{ANUPQ} package is the computation of the Burnside group $B(5, 4)$, the largest group of exponent 4 generated by 5 elements. It has order $2^{2728}$ and lower exponent\texttt{\symbol{45}}$p$ central class 13. The example \texttt{"B5\texttt{\symbol{45}}4.g"} computes $B(5, 4)$; it is based on a \texttt{pq} standalone input file written by M.{\nobreakspace}F.{\nobreakspace}Newman. 

 To be able to do examples like this was part of the motivation to provide
access to the low\texttt{\symbol{45}}level functions of the standalone program
from within \textsf{GAP}. 

 Please note that the construction uses the knowledge gained by Newman and
O'Brien in their initial construction of $B(5, 4)$, in particular, insight into the commutator structure of the group and the
knowledge of the $p$\texttt{\symbol{45}}central class and the order of $B(5, 4)$. Therefore, the construction cannot be used to prove that $B(5, 4)$ has the order and class mentioned above. It is merely a reconstruction of the
group. More information is contained in the header of the file \texttt{examples/B5\texttt{\symbol{45}}4.g}. 
\begin{Verbatim}[commandchars=!@|,fontsize=\small,frame=single,label=Example]
  procId := PqStart( FreeGroup(5) : Exponent := 4, Prime := 2 );
  Pq( procId : ClassBound := 2 );
  PqSupplyAutomorphisms( procId,
        [
          [ [ 1, 1, 0, 0, 0],      # first automorphism
            [ 0, 1, 0, 0, 0],
            [ 0, 0, 1, 0, 0],
            [ 0, 0, 0, 1, 0],
            [ 0, 0, 0, 0, 1] ],
  
          [ [ 0, 0, 0, 0, 1],      # second automorphism
            [ 1, 0, 0, 0, 0],
            [ 0, 1, 0, 0, 0],
            [ 0, 0, 1, 0, 0],
            [ 0, 0, 0, 1, 0] ]
                               ] );;
  
  Relations :=
    [ [],          ## class 1
      [],          ## class 2
      [],          ## class 3
      [],          ## class 4
      [],          ## class 5
      [],          ## class 6
      ## class 7     
      [ [ "x2","x1","x1","x3","x4","x4","x4" ] ],
      ## class 8
      [ [ "x2","x1","x1","x3","x4","x5","x5","x5" ] ],
      ## class 9
      [ [ "x2","x1","x1","x3","x4","x4","x5","x5","x5" ],
        [ "x2","x1","x1","x2","x3","x4","x5","x5","x5" ],
        [ "x2","x1","x1","x3","x3","x4","x5","x5","x5" ] ],
      ## class 10
      [ [ "x2","x1","x1","x2","x3","x3","x4","x5","x5","x5" ],
        [ "x2","x1","x1","x3","x3","x4","x4","x5","x5","x5" ] ],
      ## class 11
      [ [ "x2","x1","x1","x2","x3","x3","x4","x4","x5","x5","x5" ],
        [ "x2","x1","x1","x2","x3","x1","x3","x4","x2","x4","x3" ] ],
      ## class 12
      [ [ "x2","x1","x1","x2","x3","x1","x3","x4","x2","x5","x5","x5" ],
        [ "x2","x1","x1","x3","x2","x4","x3","x5","x4","x5","x5","x5" ] ],
      ## class 13
      [ [ "x2","x1","x1","x2","x3","x1","x3","x4","x2","x4","x5","x5","x5" 
          ] ]
  ];
  
  for class in [ 3 .. 13 ] do
      Print( "Computing class ", class, "\n" );
      PqSetupTablesForNextClass( procId );
  
      for w in [ class, class-1 .. 7 ] do
  
          PqAddTails( procId, w );   
          PqDisplayPcPresentation( procId );
  
          if Relations[ w ] <> [] then
              # recalculate automorphisms
              PqExtendAutomorphisms( procId );
  
              for r in Relations[ w ] do
                  Print( "Collecting ", r, "\n" );
                  PqCommutator( procId, r, 1 );
                  PqEchelonise( procId );
                  PqApplyAutomorphisms( procId, 15 ); #queue factor = 15
              od;
  
              PqEliminateRedundantGenerators( procId );
          fi;   
          PqComputeTails( procId, w );
      od;
      PqDisplayPcPresentation( procId );
  
      smallclass := Minimum( class, 6 );
      for w in [ smallclass, smallclass-1 .. 2 ] do
          PqTails( procId, w );
      od;
      # recalculate automorphisms
      PqExtendAutomorphisms( procId );
      PqCollect( procId, "x5^4" );
      PqEchelonise( procId );
      PqApplyAutomorphisms( procId, 15 ); #queue factor = 15
      PqEliminateRedundantGenerators( procId );
      PqDisplayPcPresentation( procId );
  od;
\end{Verbatim}
 }

 
\section{\textcolor{Chapter }{Developing descendants trees}}\label{Developing descendants trees}
\logpage{[ "A", 4, 0 ]}
\hyperdef{L}{X7FB6DA8D7BB90D8C}{}
{
  In the following example we will explore the 3\texttt{\symbol{45}}groups of
rank 2 and 3\texttt{\symbol{45}}coclass 1 up to 3\texttt{\symbol{45}}class 5.
This will be done using the $p$\texttt{\symbol{45}}group generation machinery of the package. We start with
the elementary abelian 3\texttt{\symbol{45}}group of rank 2. From within \textsf{GAP}, run the example \texttt{"PqDescendants\texttt{\symbol{45}}treetraverse\texttt{\symbol{45}}i"} via \texttt{PqExample} (see{\nobreakspace}\texttt{PqExample} (\ref{PqExample})). 
\begin{Verbatim}[commandchars=!@|,fontsize=\small,frame=single,label=Example]
  !gapprompt@gap>| !gapinput@G := ElementaryAbelianGroup( 9 );|
  <pc group of size 9 with 2 generators>
  !gapprompt@gap>| !gapinput@procId := PqStart( G );;|
  !gapprompt@gap>| !gapinput@#|
  !gapprompt@gap>| !gapinput@#  Below, we use the option StepSize in order to construct descendants|
  !gapprompt@gap>| !gapinput@#  of coclass 1. This is equivalent to setting the StepSize to 1 in|
  !gapprompt@gap>| !gapinput@#  each descendant calculation.|
  !gapprompt@gap>| !gapinput@#|
  !gapprompt@gap>| !gapinput@#  The elementary abelian group of order 9 has 3 descendants of|
  !gapprompt@gap>| !gapinput@#  3-class 2 and 3-coclass 1, as the result of the next command|
  !gapprompt@gap>| !gapinput@#  shows. |
  !gapprompt@gap>| !gapinput@#|
  !gapprompt@gap>| !gapinput@PqDescendants( procId : StepSize := 1 );|
  [ <pc group of size 27 with 3 generators>, 
    <pc group of size 27 with 3 generators>, 
    <pc group of size 27 with 3 generators> ]
  !gapprompt@gap>| !gapinput@#|
  !gapprompt@gap>| !gapinput@#  Now we will compute the descendants of coclass 1 for each of the|
  !gapprompt@gap>| !gapinput@#  groups above. Then we will compute the descendants  of coclass 1|
  !gapprompt@gap>| !gapinput@#  of each descendant and so  on.  Note  that the  pq program keeps|
  !gapprompt@gap>| !gapinput@#  one file for each class at a time.  For example, the descendants|
  !gapprompt@gap>| !gapinput@#  calculation for  the  second  group  of class  2  overwrites the|
  !gapprompt@gap>| !gapinput@#  descendant file  obtained  from  the  first  group  of  class 2.|
  !gapprompt@gap>| !gapinput@#  Hence,  we have to traverse the descendants tree  in depth first|
  !gapprompt@gap>| !gapinput@#  order.|
  !gapprompt@gap>| !gapinput@#|
  !gapprompt@gap>| !gapinput@PqPGSetDescendantToPcp( procId, 2, 1 );|
  !gapprompt@gap>| !gapinput@PqPGExtendAutomorphisms( procId );|
  !gapprompt@gap>| !gapinput@PqPGConstructDescendants( procId : StepSize := 1 );|
  2
  !gapprompt@gap>| !gapinput@PqPGSetDescendantToPcp( procId, 3, 1 );|
  !gapprompt@gap>| !gapinput@PqPGExtendAutomorphisms( procId );|
  !gapprompt@gap>| !gapinput@PqPGConstructDescendants( procId : StepSize := 1 );|
  2
  !gapprompt@gap>| !gapinput@PqPGSetDescendantToPcp( procId, 4, 1 );|
  !gapprompt@gap>| !gapinput@PqPGExtendAutomorphisms( procId );|
  !gapprompt@gap>| !gapinput@PqPGConstructDescendants( procId : StepSize := 1 );|
  2
  !gapprompt@gap>| !gapinput@#|
  !gapprompt@gap>| !gapinput@#  At this point we stop traversing the ``left most'' branch of the|
  !gapprompt@gap>| !gapinput@#  descendants tree and move upwards.|
  !gapprompt@gap>| !gapinput@#|
  !gapprompt@gap>| !gapinput@PqPGSetDescendantToPcp( procId, 4, 2 );|
  !gapprompt@gap>| !gapinput@PqPGExtendAutomorphisms( procId );|
  !gapprompt@gap>| !gapinput@PqPGConstructDescendants( procId : StepSize := 1 );|
  #I  group restored from file is incapable
  0
  !gapprompt@gap>| !gapinput@PqPGSetDescendantToPcp( procId, 3, 2 );|
  !gapprompt@gap>| !gapinput@PqPGExtendAutomorphisms( procId );|
  !gapprompt@gap>| !gapinput@PqPGConstructDescendants( procId : StepSize := 1 );|
  #I  group restored from file is incapable
  0
  !gapprompt@gap>| !gapinput@#  |
  !gapprompt@gap>| !gapinput@#  The computations above indicate that the descendants subtree under|
  !gapprompt@gap>| !gapinput@#  the first descendant of the elementary abelian group of order 9|
  !gapprompt@gap>| !gapinput@#  will have only one path of infinite length.|
  !gapprompt@gap>| !gapinput@#|
  !gapprompt@gap>| !gapinput@PqPGSetDescendantToPcp( procId, 2, 2 );|
  !gapprompt@gap>| !gapinput@PqPGExtendAutomorphisms( procId );|
  !gapprompt@gap>| !gapinput@PqPGConstructDescendants( procId : StepSize := 1 );|
  4
  !gapprompt@gap>| !gapinput@#|
  !gapprompt@gap>| !gapinput@#  We get four descendants here, three of which will turn out to be|
  !gapprompt@gap>| !gapinput@#  incapable, i.e., they have no descendants and are terminal nodes|
  !gapprompt@gap>| !gapinput@#  in the descendants tree.|
  !gapprompt@gap>| !gapinput@#|
  !gapprompt@gap>| !gapinput@PqPGSetDescendantToPcp( procId, 2, 3 );|
  !gapprompt@gap>| !gapinput@PqPGExtendAutomorphisms( procId );|
  !gapprompt@gap>| !gapinput@PqPGConstructDescendants( procId : StepSize := 1 );|
  #I  group restored from file is incapable
  0
  !gapprompt@gap>| !gapinput@#|
  !gapprompt@gap>| !gapinput@#  The third descendant of class three is incapable.  Let us return|
  !gapprompt@gap>| !gapinput@#  to the second descendant of class 2.|
  !gapprompt@gap>| !gapinput@#|
  !gapprompt@gap>| !gapinput@PqPGSetDescendantToPcp( procId, 2, 2 );|
  !gapprompt@gap>| !gapinput@PqPGExtendAutomorphisms( procId );|
  !gapprompt@gap>| !gapinput@PqPGConstructDescendants( procId : StepSize := 1 );|
  4
  !gapprompt@gap>| !gapinput@PqPGSetDescendantToPcp( procId, 3, 1 );|
  !gapprompt@gap>| !gapinput@PqPGExtendAutomorphisms( procId );|
  !gapprompt@gap>| !gapinput@PqPGConstructDescendants( procId : StepSize := 1 );|
  #I  group restored from file is incapable
  0
  !gapprompt@gap>| !gapinput@PqPGSetDescendantToPcp( procId, 3, 2 );|
  !gapprompt@gap>| !gapinput@PqPGExtendAutomorphisms( procId );|
  !gapprompt@gap>| !gapinput@PqPGConstructDescendants( procId : StepSize := 1 );|
  #I  group restored from file is incapable
  0
  !gapprompt@gap>| !gapinput@#|
  !gapprompt@gap>| !gapinput@#  We skip the third descendant for the moment ... |
  !gapprompt@gap>| !gapinput@#|
  !gapprompt@gap>| !gapinput@PqPGSetDescendantToPcp( procId, 3, 4 );|
  !gapprompt@gap>| !gapinput@PqPGExtendAutomorphisms( procId );|
  !gapprompt@gap>| !gapinput@PqPGConstructDescendants( procId : StepSize := 1 );|
  #I  group restored from file is incapable
  0
  !gapprompt@gap>| !gapinput@#|
  !gapprompt@gap>| !gapinput@#  ... and look at it now.|
  !gapprompt@gap>| !gapinput@#|
  !gapprompt@gap>| !gapinput@PqPGSetDescendantToPcp( procId, 3, 3 );|
  !gapprompt@gap>| !gapinput@PqPGExtendAutomorphisms( procId );|
  !gapprompt@gap>| !gapinput@PqPGConstructDescendants( procId : StepSize := 1 );|
  6
  !gapprompt@gap>| !gapinput@#|
  !gapprompt@gap>| !gapinput@#  In this branch of the descendant tree we get 6 descendants of class|
  !gapprompt@gap>| !gapinput@#  three.  Of those 5 will turn out to be incapable and one will have|
  !gapprompt@gap>| !gapinput@#  7 descendants.|
  !gapprompt@gap>| !gapinput@#|
  !gapprompt@gap>| !gapinput@PqPGSetDescendantToPcp( procId, 4, 1 );|
  !gapprompt@gap>| !gapinput@PqPGExtendAutomorphisms( procId );|
  !gapprompt@gap>| !gapinput@PqPGConstructDescendants( procId : StepSize := 1 );|
  #I  group restored from file is incapable
  0
  !gapprompt@gap>| !gapinput@PqPGSetDescendantToPcp( procId, 4, 2 );|
  !gapprompt@gap>| !gapinput@PqPGExtendAutomorphisms( procId );|
  !gapprompt@gap>| !gapinput@PqPGConstructDescendants( procId : StepSize := 1 );|
  7
  !gapprompt@gap>| !gapinput@PqPGSetDescendantToPcp( procId, 4, 3 );|
  !gapprompt@gap>| !gapinput@PqPGExtendAutomorphisms( procId );|
  !gapprompt@gap>| !gapinput@PqPGConstructDescendants( procId : StepSize := 1 );|
  #I  group restored from file is incapable
  0
\end{Verbatim}
 To automate the above procedure to some extent we provide: 

\subsection{\textcolor{Chapter }{PqDescendantsTreeCoclassOne}}
\logpage{[ "A", 4, 1 ]}\nobreak
\hyperdef{L}{X805E6B418193CF2D}{}
{\noindent\textcolor{FuncColor}{$\triangleright$\enspace\texttt{PqDescendantsTreeCoclassOne({\mdseries\slshape i})\index{PqDescendantsTreeCoclassOne@\texttt{PqDescendantsTreeCoclassOne}}
\label{PqDescendantsTreeCoclassOne}
}\hfill{\scriptsize (function)}}\\
\noindent\textcolor{FuncColor}{$\triangleright$\enspace\texttt{PqDescendantsTreeCoclassOne({\mdseries\slshape })\index{PqDescendantsTreeCoclassOne@\texttt{PqDescendantsTreeCoclassOne}!for default process}
\label{PqDescendantsTreeCoclassOne:for default process}
}\hfill{\scriptsize (function)}}\\


 for the \mbox{\texttt{\mdseries\slshape i}}th or default interactive \textsf{ANUPQ} process, generate a descendant tree for the group of the process (which must
be a pc $p$\texttt{\symbol{45}}group) consisting of descendants of $p$\texttt{\symbol{45}}coclass 1 and extending to the class determined by the
option \texttt{TreeDepth} (or 6 if the option is omitted). In an \textsf{XGAP} session, a graphical representation of the descendants tree appears in a
separate window. Subsequent calls to \texttt{PqDescendantsTreeCoclassOne} for the same process may be used to extend the descendant tree from the last
descendant computed that itself has more than one descendant. \texttt{PqDescendantsTreeCoclassOne} also accepts the options \texttt{CapableDescendants} (or \texttt{AllDescendants}) and any options accepted by the interactive \texttt{PqDescendants} function (see{\nobreakspace}\texttt{PqDescendants} (\ref{PqDescendants:interactive})). 

 \emph{Notes} 
\begin{enumerate}
\item  \texttt{PqDescendantsTreeCoclassOne} first calls \texttt{PqDescendants}. If \texttt{PqDescendants} has already been called for the process, the previous value computed is used
and a warning is \texttt{Info}\texttt{\symbol{45}}ed at \texttt{InfoANUPQ} level 1. 
\item  As each descendant is processed its unique label defined by the \texttt{pq} program and number of descendants is \texttt{Info}\texttt{\symbol{45}}ed at \texttt{InfoANUPQ} level 1. 
\item  \texttt{PqDescendantsTreeCoclassOne} is an ``experimental'' function that is included to demonstrate the sort of things that are possible
with the $p$\texttt{\symbol{45}}group generation machinery. 
\end{enumerate}
 Ignoring the extra functionality provided in an \textsf{XGAP} session, \texttt{PqDescendantsTreeCoclassOne}, with one argument that is the index of an interactive \textsf{ANUPQ} process, is approximately equivalent to: 
\begin{Verbatim}[commandchars=!@|,fontsize=\small,frame=single,label=]
  PqDescendantsTreeCoclassOne := function( procId )
      local des, i;
  
      des := PqDescendants( procId : StepSize := 1 );
      RecurseDescendants( procId, 2, Length(des) );
  end;
\end{Verbatim}
 where \texttt{RecurseDescendants} is (approximately) defined as follows: 
\begin{Verbatim}[commandchars=!@|,fontsize=\small,frame=single,label=]
  RecurseDescendants := function( procId, class, n )
      local i, nr;
  
      if class > ValueOption("TreeDepth") then return; fi;
  
      for i in [1..n] do
          PqPGSetDescendantToPcp( procId, class, i );
          PqPGExtendAutomorphisms( procId );
          nr := PqPGConstructDescendants( procId : StepSize := 1 );
          Print( "Number of descendants of group ", i,
                 " at class ", class, ": ", nr, "\n" );
          RecurseDescendants( procId, class+1, nr );
      od;
      return;
  end;
\end{Verbatim}
 The following examples (executed via \texttt{PqExample}; see{\nobreakspace}\texttt{PqExample} (\ref{PqExample})), demonstrate the use of \texttt{PqDescendantsTreeCoclassOne}: 
\begin{description}
\item[{\texttt{"PqDescendantsTreeCoclassOne\texttt{\symbol{45}}9\texttt{\symbol{45}}i"}}]  approximately does example \texttt{"PqDescendants\texttt{\symbol{45}}treetraverse\texttt{\symbol{45}}i"} again using \texttt{PqDescendantsTreeCoclassOne}; 
\item[{\texttt{"PqDescendantsTreeCoclassOne\texttt{\symbol{45}}16\texttt{\symbol{45}}i"}}]  uses the option \texttt{CapableDescendants}; and 
\item[{\texttt{"PqDescendantsTreeCoclassOne\texttt{\symbol{45}}25\texttt{\symbol{45}}i"}}]  calculates all descendants by omitting the \texttt{CapableDescendants} option. 
\end{description}
 The numbers \texttt{9}, \texttt{16} and \texttt{25} respectively, indicate the order of the elementary abelian group to which \texttt{PqDescendantsTreeCoclassOne} is applied for these examples. }

 }

 }

\def\bibname{References\logpage{[ "Bib", 0, 0 ]}
\hyperdef{L}{X7A6F98FD85F02BFE}{}
}

\bibliographystyle{alpha}
\bibliography{ANUPQ.bib}

\addcontentsline{toc}{chapter}{References}

\def\indexname{Index\logpage{[ "Ind", 0, 0 ]}
\hyperdef{L}{X83A0356F839C696F}{}
}

\cleardoublepage
\phantomsection
\addcontentsline{toc}{chapter}{Index}


\printindex

\immediate\write\pagenrlog{["Ind", 0, 0], \arabic{page},}
\newpage
\immediate\write\pagenrlog{["End"], \arabic{page}];}
\immediate\closeout\pagenrlog
\end{document}
